% Options for packages loaded elsewhere
\PassOptionsToPackage{unicode}{hyperref}
\PassOptionsToPackage{hyphens}{url}
\documentclass[
  12pt,
]{article}
\usepackage{xcolor}
\usepackage[left=1in,right=1in,top=1in,bottom=1in]{geometry}
\usepackage{amsmath,amssymb}
\setcounter{secnumdepth}{5}
\usepackage{iftex}
\ifPDFTeX
  \usepackage[T1]{fontenc}
  \usepackage[utf8]{inputenc}
  \usepackage{textcomp} % provide euro and other symbols
\else % if luatex or xetex
  \usepackage{unicode-math} % this also loads fontspec
  \defaultfontfeatures{Scale=MatchLowercase}
  \defaultfontfeatures[\rmfamily]{Ligatures=TeX,Scale=1}
\fi
\usepackage{lmodern}
\ifPDFTeX\else
  % xetex/luatex font selection
  \setmainfont[]{Times New Roman}
\fi
% Use upquote if available, for straight quotes in verbatim environments
\IfFileExists{upquote.sty}{\usepackage{upquote}}{}
\IfFileExists{microtype.sty}{% use microtype if available
  \usepackage[]{microtype}
  \UseMicrotypeSet[protrusion]{basicmath} % disable protrusion for tt fonts
}{}
\usepackage{setspace}
\makeatletter
\@ifundefined{KOMAClassName}{% if non-KOMA class
  \IfFileExists{parskip.sty}{%
    \usepackage{parskip}
  }{% else
    \setlength{\parindent}{0pt}
    \setlength{\parskip}{6pt plus 2pt minus 1pt}}
}{% if KOMA class
  \KOMAoptions{parskip=half}}
\makeatother
\usepackage{color}
\usepackage{fancyvrb}
\newcommand{\VerbBar}{|}
\newcommand{\VERB}{\Verb[commandchars=\\\{\}]}
\DefineVerbatimEnvironment{Highlighting}{Verbatim}{commandchars=\\\{\}}
% Add ',fontsize=\small' for more characters per line
\usepackage{framed}
\definecolor{shadecolor}{RGB}{248,248,248}
\newenvironment{Shaded}{\begin{snugshade}}{\end{snugshade}}
\newcommand{\AlertTok}[1]{\textcolor[rgb]{0.94,0.16,0.16}{#1}}
\newcommand{\AnnotationTok}[1]{\textcolor[rgb]{0.56,0.35,0.01}{\textbf{\textit{#1}}}}
\newcommand{\AttributeTok}[1]{\textcolor[rgb]{0.13,0.29,0.53}{#1}}
\newcommand{\BaseNTok}[1]{\textcolor[rgb]{0.00,0.00,0.81}{#1}}
\newcommand{\BuiltInTok}[1]{#1}
\newcommand{\CharTok}[1]{\textcolor[rgb]{0.31,0.60,0.02}{#1}}
\newcommand{\CommentTok}[1]{\textcolor[rgb]{0.56,0.35,0.01}{\textit{#1}}}
\newcommand{\CommentVarTok}[1]{\textcolor[rgb]{0.56,0.35,0.01}{\textbf{\textit{#1}}}}
\newcommand{\ConstantTok}[1]{\textcolor[rgb]{0.56,0.35,0.01}{#1}}
\newcommand{\ControlFlowTok}[1]{\textcolor[rgb]{0.13,0.29,0.53}{\textbf{#1}}}
\newcommand{\DataTypeTok}[1]{\textcolor[rgb]{0.13,0.29,0.53}{#1}}
\newcommand{\DecValTok}[1]{\textcolor[rgb]{0.00,0.00,0.81}{#1}}
\newcommand{\DocumentationTok}[1]{\textcolor[rgb]{0.56,0.35,0.01}{\textbf{\textit{#1}}}}
\newcommand{\ErrorTok}[1]{\textcolor[rgb]{0.64,0.00,0.00}{\textbf{#1}}}
\newcommand{\ExtensionTok}[1]{#1}
\newcommand{\FloatTok}[1]{\textcolor[rgb]{0.00,0.00,0.81}{#1}}
\newcommand{\FunctionTok}[1]{\textcolor[rgb]{0.13,0.29,0.53}{\textbf{#1}}}
\newcommand{\ImportTok}[1]{#1}
\newcommand{\InformationTok}[1]{\textcolor[rgb]{0.56,0.35,0.01}{\textbf{\textit{#1}}}}
\newcommand{\KeywordTok}[1]{\textcolor[rgb]{0.13,0.29,0.53}{\textbf{#1}}}
\newcommand{\NormalTok}[1]{#1}
\newcommand{\OperatorTok}[1]{\textcolor[rgb]{0.81,0.36,0.00}{\textbf{#1}}}
\newcommand{\OtherTok}[1]{\textcolor[rgb]{0.56,0.35,0.01}{#1}}
\newcommand{\PreprocessorTok}[1]{\textcolor[rgb]{0.56,0.35,0.01}{\textit{#1}}}
\newcommand{\RegionMarkerTok}[1]{#1}
\newcommand{\SpecialCharTok}[1]{\textcolor[rgb]{0.81,0.36,0.00}{\textbf{#1}}}
\newcommand{\SpecialStringTok}[1]{\textcolor[rgb]{0.31,0.60,0.02}{#1}}
\newcommand{\StringTok}[1]{\textcolor[rgb]{0.31,0.60,0.02}{#1}}
\newcommand{\VariableTok}[1]{\textcolor[rgb]{0.00,0.00,0.00}{#1}}
\newcommand{\VerbatimStringTok}[1]{\textcolor[rgb]{0.31,0.60,0.02}{#1}}
\newcommand{\WarningTok}[1]{\textcolor[rgb]{0.56,0.35,0.01}{\textbf{\textit{#1}}}}
\usepackage{longtable,booktabs,array}
\newcounter{none} % for unnumbered tables
\usepackage{calc} % for calculating minipage widths
% Correct order of tables after \paragraph or \subparagraph
\usepackage{etoolbox}
\makeatletter
\patchcmd\longtable{\par}{\if@noskipsec\mbox{}\fi\par}{}{}
\makeatother
% Allow footnotes in longtable head/foot
\IfFileExists{footnotehyper.sty}{\usepackage{footnotehyper}}{\usepackage{footnote}}
\makesavenoteenv{longtable}
\usepackage{graphicx}
\makeatletter
\newsavebox\pandoc@box
\newcommand*\pandocbounded[1]{% scales image to fit in text height/width
  \sbox\pandoc@box{#1}%
  \Gscale@div\@tempa{\textheight}{\dimexpr\ht\pandoc@box+\dp\pandoc@box\relax}%
  \Gscale@div\@tempb{\linewidth}{\wd\pandoc@box}%
  \ifdim\@tempb\p@<\@tempa\p@\let\@tempa\@tempb\fi% select the smaller of both
  \ifdim\@tempa\p@<\p@\scalebox{\@tempa}{\usebox\pandoc@box}%
  \else\usebox{\pandoc@box}%
  \fi%
}
% Set default figure placement to htbp
\def\fps@figure{htbp}
\makeatother
\setlength{\emergencystretch}{3em} % prevent overfull lines
\providecommand{\tightlist}{%
  \setlength{\itemsep}{0pt}\setlength{\parskip}{0pt}}
\usepackage[]{natbib}
\bibliographystyle{plainnat}
\usepackage{indentfirst}
\setlength{\parindent}{1.5cm}
\setlength{\parskip}{0pt}
\usepackage{titlesec}
\titlespacing*{\section}{0pt}{1.2ex}{0.5ex}
\titlespacing*{\subsection}{0pt}{1ex}{0.3ex}
\titlespacing*{\subsubsection}{0pt}{0.8ex}{0.2ex}
\usepackage[font=small]{caption}
\usepackage{multicol}
\usepackage[none]{hyphenat}
\setlength{\columnsep}{0.3cm}
\setlength{\columnseprule}{0pt}
\usepackage{bookmark}
\IfFileExists{xurl.sty}{\usepackage{xurl}}{} % add URL line breaks if available
\urlstyle{same}
\hypersetup{
  pdftitle={amylase production from bacteria},
  pdfauthor={nmz},
  hidelinks,
  pdfcreator={LaTeX via pandoc}}

\title{amylase production from bacteria}
\author{nmz}
\date{2026-08-07}

\begin{document}
\maketitle

{
\setcounter{tocdepth}{3}
\tableofcontents
}
\setstretch{1}
\section{Abstract}\label{abstract}

\begin{multicols}{2}
This study presents a systematic investigation of bacterial $\alpha$-amylase catalytic properties through integrated spectrophotometric, kinetic, and response surface methodologies [@ataallahiSimple2stepPurification2021]. The analytical workflow encompassed glucose calibration via 3,5-dinitrosalicylic acid (DNS) assay, fermentation kinetics under varying salt concentrations, response surface optimization of pH and temperature, comparative substrate specificity across four botanical starches, and Michaelis-Menten parameter estimation [@ataallahiSimple2stepPurification2021; @aENZYMICSTUDIESAlphaAMYLASE2019; @ataallahiSimple2stepPurification2021b]. The glucose standard curve demonstrated exceptional linearity ($R^2$ = 0.9909), establishing a robust empirical transfer function for reducing sugar quantification. Response surface methodology identified optimal catalytic conditions at pH 5.96 and 49.89°C, with both eigenvalues negative ($\lambda_1$ = -0.381, $\lambda_2$ = -0.471), confirming a true global maximum. Kinetic analysis revealed pronounced substrate-dependent variation in catalytic efficiency, with soluble starch exhibiting the highest $V_\text{max}$ (0.139 $\Delta A_{620}$ min$^{-1}$) and lowest $K_\text{m}$ (3.835% w/v), while rice starch demonstrated the lowest $V_\text{max}$ (0.077 $\Delta A_{620}$·min$^{-1}$) and $K_\text{m}$ (1.409% w/v). The catalytic efficiency index ($V_\text{max}$/$K_\text{m}$) ranged from 0.039 for rice starch to 0.054 for tapioca starch, indicating differential substrate accessibility governed by granular architecture and crystallinity. Progress curve analysis via second-order polynomial modeling revealed systematic deviations from linearity, with negative residuals at early timepoints (t = 0-15 min) indicating rapid initial substrate cleavage exceeding model predictions, while positive residuals at late timepoints (t = 45-60 min) reflected product inhibition and substrate depletion not captured by global quadratic fits. These findings demonstrate that bacterial $\alpha$-amylase exhibits optimal activity under moderately acidic conditions and displays marked substrate preference governed by starch botanical origin, with implications for industrial starch processing applications.

\end{multicols}

\section{Data Import}\label{data-import}

\begin{itemize}
\tightlist
\item
  load the libraries
\item
  load the data for standard curve
\item
  load the experimental data
\end{itemize}

\begin{Shaded}
\begin{Highlighting}[]
\CommentTok{\# loading libraries}
\FunctionTok{library}\NormalTok{(ggpubr)}
\FunctionTok{library}\NormalTok{(tidyverse)}
\FunctionTok{library}\NormalTok{(readxl)}
\FunctionTok{library}\NormalTok{(Myfuns)}
\FunctionTok{library}\NormalTok{(patchwork)}
\FunctionTok{library}\NormalTok{(broom)}

\CommentTok{\# loading standard curve data}
\NormalTok{std\_c }\OtherTok{\textless{}{-}} \FunctionTok{read\_excel}\NormalTok{(}\StringTok{"../Data/Salting out and DNS assay.xlsx"}\NormalTok{,}
  \AttributeTok{sheet =} \StringTok{"Sheet3"}
\NormalTok{)}
\FunctionTok{names}\NormalTok{(std\_c)}
\end{Highlighting}
\end{Shaded}

\begin{verbatim}
[1] "Glu_conc" "Abs"     
\end{verbatim}

\begin{Shaded}
\begin{Highlighting}[]
\CommentTok{\# loading Experiment Data}
\NormalTok{so }\OtherTok{\textless{}{-}} \FunctionTok{read\_excel}\NormalTok{(}\StringTok{"../Data/Salting out and DNS assay.xlsx"}\NormalTok{,}
  \AttributeTok{sheet =} \StringTok{"Sheet2"}
\NormalTok{)}
\FunctionTok{names}\NormalTok{(so)}
\end{Highlighting}
\end{Shaded}

\begin{verbatim}
[1] "Time"      "Salt_conc" "Abs"      
\end{verbatim}

\begin{Shaded}
\begin{Highlighting}[]
\NormalTok{buf }\OtherTok{\textless{}{-}} \FloatTok{0.6713}
\NormalTok{so2 }\OtherTok{\textless{}{-}}\NormalTok{ so }\SpecialCharTok{\%\textgreater{}\%} \FunctionTok{mutate}\NormalTok{(}\AttributeTok{Abs =}\NormalTok{ Abs }\SpecialCharTok{{-}}\NormalTok{ buf)}
\end{Highlighting}
\end{Shaded}

\section{Relationship of the Glucose concentration and
Absorbance}\label{relationship-of-the-glucose-concentration-and-absorbance}

\begin{itemize}
\tightlist
\item
  fit the linear regression
\item
  visualize the linear relationship
\end{itemize}

\begin{Shaded}
\begin{Highlighting}[]
\NormalTok{std\_c\_eq1 }\OtherTok{\textless{}{-}} \FunctionTok{lm}\NormalTok{(Glu\_conc }\SpecialCharTok{\textasciitilde{}}\NormalTok{ Abs, std\_c)}
\FunctionTok{summary}\NormalTok{(std\_c\_eq1)}
\end{Highlighting}
\end{Shaded}

\begin{verbatim}

Call:
lm(formula = Glu_conc ~ Abs, data = std_c)

Residuals:
      Min        1Q    Median        3Q       Max 
-0.041356 -0.027569 -0.006494  0.006797  0.065767 

Coefficients:
            Estimate Std. Error t value Pr(>|t|)    
(Intercept) -0.03242    0.01513  -2.143   0.0478 *  
Abs          0.46291    0.01109  41.760   <2e-16 ***
---
Signif. codes:  0 '***' 0.001 '**' 0.01 '*' 0.05 '.' 0.1 ' ' 1

Residual standard error: 0.03454 on 16 degrees of freedom
Multiple R-squared:  0.9909,    Adjusted R-squared:  0.9903 
F-statistic:  1744 on 1 and 16 DF,  p-value: < 2.2e-16
\end{verbatim}

\begin{Shaded}
\begin{Highlighting}[]
\NormalTok{std\_c\_eq2 }\OtherTok{\textless{}{-}} \FunctionTok{lm}\NormalTok{(Abs }\SpecialCharTok{\textasciitilde{}}\NormalTok{ Glu\_conc, std\_c)}
\FunctionTok{summary}\NormalTok{(std\_c\_eq2)}
\end{Highlighting}
\end{Shaded}

\begin{verbatim}

Call:
lm(formula = Abs ~ Glu_conc, data = std_c)

Residuals:
     Min       1Q   Median       3Q      Max 
-0.14011 -0.02057  0.01599  0.05024  0.09916 

Coefficients:
            Estimate Std. Error t value Pr(>|t|)    
(Intercept)  0.07986    0.03104   2.573   0.0204 *  
Glu_conc     2.14059    0.05126  41.760   <2e-16 ***
---
Signif. codes:  0 '***' 0.001 '**' 0.01 '*' 0.05 '.' 0.1 ' ' 1

Residual standard error: 0.07428 on 16 degrees of freedom
Multiple R-squared:  0.9909,    Adjusted R-squared:  0.9903 
F-statistic:  1744 on 1 and 16 DF,  p-value: < 2.2e-16
\end{verbatim}

\begin{Shaded}
\begin{Highlighting}[]
\NormalTok{p1 }\OtherTok{\textless{}{-}}\NormalTok{ std\_c }\SpecialCharTok{\%\textgreater{}\%} \FunctionTok{reframe}\NormalTok{(}\AttributeTok{Abs=}\FunctionTok{mean}\NormalTok{(Abs),}\AttributeTok{.by =}\NormalTok{ Glu\_conc) }\SpecialCharTok{\%\textgreater{}\%} 
  \FunctionTok{ggplot}\NormalTok{(}\FunctionTok{aes}\NormalTok{(Glu\_conc, Abs)) }\SpecialCharTok{+}
  \FunctionTok{geom\_point}\NormalTok{() }\SpecialCharTok{+}
  \FunctionTok{geom\_smooth}\NormalTok{(}\AttributeTok{method =} \StringTok{"lm"}\NormalTok{, }\AttributeTok{color =} \StringTok{"red"}\NormalTok{, }\AttributeTok{se =}\NormalTok{ F) }\SpecialCharTok{+}
  \CommentTok{\# geom\_smooth(linetype = "dashed", color = "blue") +}
  \FunctionTok{stat\_regline\_equation}\NormalTok{(}
    \FunctionTok{aes}\NormalTok{(}\AttributeTok{label =} \FunctionTok{paste}\NormalTok{(..eq.label..,}
\NormalTok{                      ..adj.rr.label..,}
                      \AttributeTok{sep =} \StringTok{"\textasciitilde{}\textasciitilde{}\textasciitilde{}\textasciitilde{}\textasciitilde{}"}\NormalTok{)),}
    \AttributeTok{decreasing =} \ConstantTok{FALSE}
\NormalTok{  )}\SpecialCharTok{+}
  \FunctionTok{labs}\NormalTok{(}\AttributeTok{x =} \StringTok{"Glucose Concentration (mg/ml)"}\NormalTok{,}
       \AttributeTok{y =} \StringTok{"Absorbance (at 570 nm)"}\NormalTok{)}\SpecialCharTok{+}
  \FunctionTok{Mytheme}\NormalTok{(}\AttributeTok{font =} \StringTok{"TeX Gyre Termes"}\NormalTok{)}
  

\NormalTok{p2 }\OtherTok{\textless{}{-}}\NormalTok{ std\_c }\SpecialCharTok{\%\textgreater{}\%} \FunctionTok{ggplot}\NormalTok{(}\FunctionTok{aes}\NormalTok{(Abs, Glu\_conc)) }\SpecialCharTok{+}
  \FunctionTok{geom\_point}\NormalTok{() }\SpecialCharTok{+}
  \FunctionTok{geom\_smooth}\NormalTok{(}\AttributeTok{method =} \StringTok{"lm"}\NormalTok{, }\AttributeTok{color =} \StringTok{"red"}\NormalTok{) }\SpecialCharTok{+}
  \CommentTok{\# geom\_smooth(linetype = "dashed", color = "blue") +}
  \FunctionTok{stat\_regline\_equation}\NormalTok{(}
    \FunctionTok{aes}\NormalTok{(}\AttributeTok{label =} \FunctionTok{paste}\NormalTok{(..eq.label..,}
\NormalTok{                      ..adj.rr.label..,}
                      \AttributeTok{sep =} \StringTok{"\textasciitilde{}\textasciitilde{}\textasciitilde{}\textasciitilde{}\textasciitilde{}"}\NormalTok{)),}
    \CommentTok{\# aes(label = paste0("atop(", ..eq.label.., ",", ..adj.rr.label.., ")")),}
    \AttributeTok{decreasing =} \ConstantTok{FALSE}
\NormalTok{  )}\SpecialCharTok{+}
  \FunctionTok{labs}\NormalTok{(}\AttributeTok{x=} \StringTok{"Absorbance (at 570 nm)"}\NormalTok{,}
       \AttributeTok{y=} \StringTok{"Glucose Concentration (mg/ml)"}\NormalTok{)}\SpecialCharTok{+}
  \FunctionTok{Mytheme}\NormalTok{(}\AttributeTok{font =} \StringTok{"TeX Gyre Termes"}\NormalTok{)}

\NormalTok{p1}\SpecialCharTok{+}\NormalTok{p2}\SpecialCharTok{+}
  \FunctionTok{plot\_annotation}\NormalTok{(}\AttributeTok{tag\_levels =} \StringTok{"A"}\NormalTok{)}
\end{Highlighting}
\end{Shaded}

\begin{verbatim}
`geom_smooth()` using formula = 'y ~ x'
`geom_smooth()` using formula = 'y ~ x'
\end{verbatim}

\begin{center}\includegraphics[width=1\linewidth]{../Figs/std_curve-1} \end{center}

\section{Prepare the experiment data for
analysis}\label{prepare-the-experiment-data-for-analysis}

\begin{Shaded}
\begin{Highlighting}[]
\FunctionTok{names}\NormalTok{(so)}
\end{Highlighting}
\end{Shaded}

\begin{verbatim}
[1] "Time"      "Salt_conc" "Abs"      
\end{verbatim}

\begin{Shaded}
\begin{Highlighting}[]
\NormalTok{buf }\OtherTok{\textless{}{-}} \FloatTok{0.6713}
\NormalTok{so2 }\OtherTok{\textless{}{-}}\NormalTok{ so }\SpecialCharTok{\%\textgreater{}\%} \FunctionTok{mutate}\NormalTok{(}\AttributeTok{Abs =}\NormalTok{ Abs }\SpecialCharTok{{-}}\NormalTok{ buf)}

\NormalTok{pre\_glu }\OtherTok{\textless{}{-}} \FunctionTok{predict}\NormalTok{(std\_c\_eq1, }\FunctionTok{data.frame}\NormalTok{(}\AttributeTok{Abs =}\NormalTok{ so2}\SpecialCharTok{$}\NormalTok{Abs))}
\CommentTok{\# pre\_glu2 \textless{}{-} predict(std\_c\_eq2, data.frame(Abs = so2$Abs))}

\NormalTok{so2}\SpecialCharTok{$}\NormalTok{glu }\OtherTok{\textless{}{-}}\NormalTok{ pre\_glu}
\NormalTok{so2 }\OtherTok{\textless{}{-}}\NormalTok{ so2[}\SpecialCharTok{{-}}\FunctionTok{c}\NormalTok{(}\DecValTok{1}\SpecialCharTok{:}\DecValTok{6}\NormalTok{), ] }
\NormalTok{p3 }\OtherTok{\textless{}{-}}\NormalTok{ so2 }\SpecialCharTok{\%\textgreater{}\%}
  \FunctionTok{mutate}\NormalTok{(}\AttributeTok{Salt\_conc =} \FunctionTok{factor}\NormalTok{(Salt\_conc)) }\SpecialCharTok{\%\textgreater{}\%}
  \FunctionTok{ggplot}\NormalTok{(}\FunctionTok{aes}\NormalTok{(Time, glu, }\AttributeTok{group =} \FunctionTok{factor}\NormalTok{(Salt\_conc))) }\SpecialCharTok{+}
  \FunctionTok{geom\_point}\NormalTok{(}\FunctionTok{aes}\NormalTok{(}\AttributeTok{color =}\NormalTok{ Salt\_conc)) }\SpecialCharTok{+}
  \FunctionTok{geom\_smooth}\NormalTok{(}\FunctionTok{aes}\NormalTok{(}\AttributeTok{color =}\NormalTok{ Salt\_conc),}
    \AttributeTok{se =}\NormalTok{ F,}
    \CommentTok{\# method = "lm",}
    \CommentTok{\# formula = y \textasciitilde{} poly(x, 2, raw = T)}
\NormalTok{  ) }\SpecialCharTok{+}
  \FunctionTok{labs}\NormalTok{(}
    \AttributeTok{x =} \StringTok{"Growth time (hour)"}\NormalTok{,}
    \AttributeTok{y =} \StringTok{"Glucose Concentration}
\StringTok{    (mg/ml)"}\NormalTok{,}
    \AttributeTok{color =} \StringTok{"Salt Concentration"}
\NormalTok{  ) }\SpecialCharTok{+}
  \FunctionTok{scale\_color\_discrete}\NormalTok{(}
    \AttributeTok{breaks =} \FunctionTok{c}\NormalTok{(}\DecValTok{10}\NormalTok{, }\DecValTok{20}\NormalTok{, }\DecValTok{30}\NormalTok{),}
    \AttributeTok{labels =} \FunctionTok{c}\NormalTok{(}\StringTok{"10 \%"}\NormalTok{, }\StringTok{"20 \%"}\NormalTok{, }\StringTok{"30 \%"}\NormalTok{)}
\NormalTok{  ) }\SpecialCharTok{+}
  \FunctionTok{theme\_minimal}\NormalTok{() }\SpecialCharTok{+}
  \FunctionTok{theme}\NormalTok{(}
    \AttributeTok{text =} \FunctionTok{element\_text}\NormalTok{(}
      \AttributeTok{family =}\NormalTok{ mfont,}
      \AttributeTok{color =} \StringTok{"black"}\NormalTok{,}
      \AttributeTok{size =} \DecValTok{12}
\NormalTok{    ),}
    \AttributeTok{axis.line =} \FunctionTok{element\_line}\NormalTok{(}\AttributeTok{color =} \StringTok{"black"}\NormalTok{),}
    \AttributeTok{axis.text =} \FunctionTok{element\_text}\NormalTok{(}\AttributeTok{color =} \StringTok{"black"}\NormalTok{),}
    \AttributeTok{legend.position =} \StringTok{"bottom"}\NormalTok{,}
    \AttributeTok{legend.margin =} \FunctionTok{margin}\NormalTok{(}\SpecialCharTok{{-}}\DecValTok{10}\NormalTok{, }\DecValTok{0}\NormalTok{, }\DecValTok{0}\NormalTok{, }\DecValTok{0}\NormalTok{),}
    \AttributeTok{plot.margin =} \FunctionTok{margin}\NormalTok{(}\DecValTok{5}\NormalTok{, }\DecValTok{5}\NormalTok{, }\DecValTok{5}\NormalTok{, }\DecValTok{5}\NormalTok{)}
\NormalTok{  )}

\DocumentationTok{\#\# For Contour plot}
\NormalTok{salt\_range }\OtherTok{\textless{}{-}} \FunctionTok{range}\NormalTok{(so2}\SpecialCharTok{$}\NormalTok{Salt\_conc)}
\NormalTok{time\_range }\OtherTok{\textless{}{-}} \FunctionTok{range}\NormalTok{(so2}\SpecialCharTok{$}\NormalTok{Time)}

\NormalTok{salt\_seq }\OtherTok{\textless{}{-}} \FunctionTok{seq}\NormalTok{(salt\_range[}\DecValTok{1}\NormalTok{], salt\_range[}\DecValTok{2}\NormalTok{], }\AttributeTok{length.out =} \DecValTok{30}\NormalTok{)}
\NormalTok{time\_seq }\OtherTok{\textless{}{-}} \FunctionTok{seq}\NormalTok{(time\_range[}\DecValTok{1}\NormalTok{], time\_range[}\DecValTok{2}\NormalTok{], }\AttributeTok{length.out =} \DecValTok{30}\NormalTok{)}
\NormalTok{prediction\_grid }\OtherTok{\textless{}{-}} \FunctionTok{expand.grid}\NormalTok{(}\AttributeTok{Salt\_conc =}\NormalTok{ salt\_seq, }\AttributeTok{Time =}\NormalTok{ time\_seq)}

\NormalTok{prediction\_grid}\SpecialCharTok{$}\NormalTok{glu }\OtherTok{\textless{}{-}} \FunctionTok{predict}\NormalTok{(}\FunctionTok{lm}\NormalTok{(glu }\SpecialCharTok{\textasciitilde{}} \FunctionTok{poly}\NormalTok{(Salt\_conc, }\DecValTok{2}\NormalTok{, }\AttributeTok{raw =}\NormalTok{ T) }\SpecialCharTok{*} 
                                    \FunctionTok{poly}\NormalTok{(Time, }\DecValTok{2}\NormalTok{, }\AttributeTok{raw =}\NormalTok{ T), }\AttributeTok{data =}\NormalTok{ so2),}
  \AttributeTok{newdata =}\NormalTok{ prediction\_grid}
\NormalTok{)}

\NormalTok{p4 }\OtherTok{\textless{}{-}}\NormalTok{ prediction\_grid }\SpecialCharTok{\%\textgreater{}\%}
  \FunctionTok{ggplot}\NormalTok{(}\FunctionTok{aes}\NormalTok{(Time, Salt\_conc)) }\SpecialCharTok{+}
  \FunctionTok{geom\_tile}\NormalTok{(}\FunctionTok{aes}\NormalTok{(}\AttributeTok{fill =}\NormalTok{ glu)) }\SpecialCharTok{+}
  \FunctionTok{geom\_contour}\NormalTok{(}\FunctionTok{aes}\NormalTok{(}\AttributeTok{z =}\NormalTok{ glu), }\AttributeTok{color =} \StringTok{"white"}\NormalTok{) }\SpecialCharTok{+}
  \FunctionTok{scale\_fill\_viridis\_c}\NormalTok{(}
    \AttributeTok{name =} \StringTok{"Glucose}\SpecialCharTok{\textbackslash{}n}\StringTok{(mg/ml)"}\NormalTok{, }\AttributeTok{direction =} \SpecialCharTok{{-}}\DecValTok{1}\NormalTok{,}
    \AttributeTok{limits =} \FunctionTok{c}\NormalTok{(}\FloatTok{0.3}\NormalTok{, }\FloatTok{0.67}\NormalTok{),}
    \AttributeTok{breaks =} \FunctionTok{seq}\NormalTok{(}\FloatTok{0.3}\NormalTok{, }\FloatTok{0.67}\NormalTok{, }\FloatTok{0.1}\NormalTok{)}
\NormalTok{  ) }\SpecialCharTok{+}
  \FunctionTok{labs}\NormalTok{(}
    \AttributeTok{x =} \StringTok{"Growth time (hour)"}\NormalTok{,}
    \AttributeTok{y =} \StringTok{"Salt Concentration"}
    
\NormalTok{  ) }\SpecialCharTok{+}
  \FunctionTok{theme}\NormalTok{(}
    \AttributeTok{text =} \FunctionTok{element\_text}\NormalTok{(}
      \AttributeTok{family =}\NormalTok{ mfont,}
      \AttributeTok{size =} \DecValTok{12}\NormalTok{,}
      \AttributeTok{color =} \StringTok{"black"}
\NormalTok{    ),}
    \AttributeTok{axis.text =} \FunctionTok{element\_text}\NormalTok{(}
      \AttributeTok{family =}\NormalTok{ mfont,}
      \AttributeTok{size =} \DecValTok{12}\NormalTok{,}
      \AttributeTok{color =} \StringTok{"black"}
\NormalTok{    ),}
    \AttributeTok{axis.line =} \FunctionTok{element\_blank}\NormalTok{(),}
    \AttributeTok{legend.spacing.x =} \FunctionTok{unit}\NormalTok{(}\DecValTok{5}\NormalTok{, }\StringTok{"pt"}\NormalTok{),}
    \AttributeTok{legend.text =} \FunctionTok{element\_text}\NormalTok{(}
      \AttributeTok{margin =} \FunctionTok{margin}\NormalTok{(}\AttributeTok{l =} \DecValTok{5}\NormalTok{),}
      \AttributeTok{vjust =} \SpecialCharTok{{-}}\NormalTok{.}\DecValTok{2}
\NormalTok{    ),}
    \AttributeTok{panel.background =} \FunctionTok{element\_blank}\NormalTok{()}
\NormalTok{  )}

\NormalTok{p3}\SpecialCharTok{+}\NormalTok{p4}\SpecialCharTok{+}\FunctionTok{plot\_annotation}\NormalTok{(}\AttributeTok{tag\_levels =} \StringTok{"A"}\NormalTok{)}
\end{Highlighting}
\end{Shaded}

\begin{verbatim}
`geom_smooth()` using method = 'loess' and formula = 'y ~ x'
\end{verbatim}

\begin{center}\includegraphics[width=1\linewidth]{../Figs/exp_data-1} \end{center}

\subsection{model used for contour
plot}\label{model-used-for-contour-plot}

\begin{Shaded}
\begin{Highlighting}[]
\NormalTok{mod\_used }\OtherTok{\textless{}{-}} \FunctionTok{lm}\NormalTok{(glu }\SpecialCharTok{\textasciitilde{}} \FunctionTok{poly}\NormalTok{(Salt\_conc, }\DecValTok{2}\NormalTok{, }\AttributeTok{raw =}\NormalTok{ T) }\SpecialCharTok{*} \FunctionTok{poly}\NormalTok{(Time, }\DecValTok{2}\NormalTok{, }\AttributeTok{raw =}\NormalTok{ T), }
               \AttributeTok{data =}\NormalTok{ so2)}
\FunctionTok{summary}\NormalTok{(mod\_used)}
\end{Highlighting}
\end{Shaded}

\begin{verbatim}

Call:
lm(formula = glu ~ poly(Salt_conc, 2, raw = T) * poly(Time, 2, 
    raw = T), data = so2)

Residuals:
      Min        1Q    Median        3Q       Max 
-0.083070 -0.019811 -0.003844  0.024736  0.078950 

Coefficients:
                                                       Estimate Std. Error
(Intercept)                                           1.317e+01  4.184e+00
poly(Salt_conc, 2, raw = T)1                         -1.161e+00  5.006e-01
poly(Salt_conc, 2, raw = T)2                          2.013e-02  1.351e-02
poly(Time, 2, raw = T)1                              -4.081e-01  1.401e-01
poly(Time, 2, raw = T)2                               3.296e-03  1.156e-03
poly(Salt_conc, 2, raw = T)1:poly(Time, 2, raw = T)1  3.592e-02  1.666e-02
poly(Salt_conc, 2, raw = T)2:poly(Time, 2, raw = T)1 -5.921e-04  4.458e-04
poly(Salt_conc, 2, raw = T)1:poly(Time, 2, raw = T)2 -2.799e-04  1.368e-04
poly(Salt_conc, 2, raw = T)2:poly(Time, 2, raw = T)2  4.436e-06  3.634e-06
                                                     t value Pr(>|t|)   
(Intercept)                                            3.149  0.00529 **
poly(Salt_conc, 2, raw = T)1                          -2.320  0.03163 * 
poly(Salt_conc, 2, raw = T)2                           1.490  0.15269   
poly(Time, 2, raw = T)1                               -2.913  0.00892 **
poly(Time, 2, raw = T)2                                2.851  0.01022 * 
poly(Salt_conc, 2, raw = T)1:poly(Time, 2, raw = T)1   2.155  0.04416 * 
poly(Salt_conc, 2, raw = T)2:poly(Time, 2, raw = T)1  -1.328  0.19988   
poly(Salt_conc, 2, raw = T)1:poly(Time, 2, raw = T)2  -2.046  0.05487 . 
poly(Salt_conc, 2, raw = T)2:poly(Time, 2, raw = T)2   1.221  0.23719   
---
Signif. codes:  0 '***' 0.001 '**' 0.01 '*' 0.05 '.' 0.1 ' ' 1

Residual standard error: 0.04749 on 19 degrees of freedom
Multiple R-squared:  0.841, Adjusted R-squared:  0.774 
F-statistic: 12.56 on 8 and 19 DF,  p-value: 4.099e-06
\end{verbatim}

\section{Optimization of pH and Temperature for maximum enzyme
activity}\label{optimization-of-ph-and-temperature-for-maximum-enzyme-activity}

\begin{itemize}
\tightlist
\item
  load data
\item
  fit to response surface
\item
  visualize the response surface
\item
  calculation of optimal condition \#\# load library and data
\end{itemize}

\begin{Shaded}
\begin{Highlighting}[]
\FunctionTok{library}\NormalTok{(rsm)}
\NormalTok{RSM }\OtherTok{\textless{}{-}} \FunctionTok{read\_excel}\NormalTok{(}\StringTok{"../Data/(26.06.2026) absorbance data.xlsx"}\NormalTok{)}
\NormalTok{RSM2 }\OtherTok{\textless{}{-}}\NormalTok{ RSM }\SpecialCharTok{\%\textgreater{}\%} \FunctionTok{mutate}\NormalTok{(}\AttributeTok{Temp =} \FunctionTok{str\_remove}\NormalTok{(Temp, }\StringTok{"°C"}\NormalTok{) }\SpecialCharTok{\%\textgreater{}\%} \FunctionTok{as.numeric}\NormalTok{())}

\NormalTok{codedrsm2 }\OtherTok{\textless{}{-}} \FunctionTok{coded.data}\NormalTok{(RSM2, }
\NormalTok{                        x1}\SpecialCharTok{\textasciitilde{}}\NormalTok{(Temp }\SpecialCharTok{{-}}\DecValTok{50}\NormalTok{)}\SpecialCharTok{/}\DecValTok{10}\NormalTok{,}
\NormalTok{                        x2 }\SpecialCharTok{\textasciitilde{}}\NormalTok{(pH }\SpecialCharTok{{-}} \DecValTok{6}\NormalTok{)}\SpecialCharTok{/}\FloatTok{0.5}\NormalTok{)}
\end{Highlighting}
\end{Shaded}

\subsection{fit the model}\label{fit-the-model}

\begin{Shaded}
\begin{Highlighting}[]
\NormalTok{rsmmod }\OtherTok{\textless{}{-}} \FunctionTok{rsm}\NormalTok{(avg}\SpecialCharTok{\textasciitilde{}}\FunctionTok{SO}\NormalTok{(x1, x2), codedrsm2)}
\FunctionTok{print}\NormalTok{(rsmmod)}
\end{Highlighting}
\end{Shaded}

\begin{verbatim}

Call:
rsm(formula = avg ~ SO(x1, x2), data = codedrsm2)

Coefficients:
   (Intercept)    FO(x1, x2)x1    FO(x1, x2)x2     TWI(x1, x2)  PQ(x1, x2)x1^2  
       1.22668        -0.01519        -0.05774        -0.06495        -0.45733  
PQ(x1, x2)x2^2  
      -0.39470  
\end{verbatim}

\begin{Shaded}
\begin{Highlighting}[]
\FunctionTok{summary}\NormalTok{(rsmmod)}
\end{Highlighting}
\end{Shaded}

\begin{verbatim}

Call:
rsm(formula = avg ~ SO(x1, x2), data = codedrsm2)

             Estimate Std. Error  t value  Pr(>|t|)    
(Intercept)  1.226682   0.019478  62.9775 1.078e-09 ***
x1          -0.015189   0.013642  -1.1134  0.308139    
x2          -0.057740   0.013851  -4.1686  0.005888 ** 
x1:x2       -0.064950   0.019491  -3.3324  0.015760 *  
x1^2        -0.457332   0.014798 -30.9044 7.622e-08 ***
x2^2        -0.394703   0.015681 -25.1709 2.589e-07 ***
---
Signif. codes:  0 '***' 0.001 '**' 0.01 '*' 0.05 '.' 0.1 ' ' 1

Multiple R-squared:  0.9956,    Adjusted R-squared:  0.9919 
F-statistic: 270.6 on 5 and 6 DF,  p-value: 5.62e-07

Analysis of Variance Table

Response: avg
            Df  Sum Sq Mean Sq  F value    Pr(>F)
FO(x1, x2)   2 0.02779 0.01389   9.1435   0.01508
TWI(x1, x2)  1 0.01687 0.01687  11.1048   0.01576
PQ(x1, x2)   2 2.01124 1.00562 661.8022 9.189e-08
Residuals    6 0.00912 0.00152                   
Lack of fit  3 0.00261 0.00087   0.4017   0.76322
Pure error   3 0.00650 0.00217                   

Stationary point of response surface:
         x1          x2 
-0.01147885 -0.07219898 

Stationary point in original units:
     Temp        pH 
49.885212  5.963901 

Eigenanalysis:
eigen() decomposition
$values
[1] -0.3809044 -0.4711314

$vectors
         [,1]       [,2]
x1  0.3910698 -0.9203610
x2 -0.9203610 -0.3910698
\end{verbatim}

\begin{Shaded}
\begin{Highlighting}[]
\FunctionTok{canonical}\NormalTok{(rsmmod)}
\end{Highlighting}
\end{Shaded}

\begin{verbatim}
$xs
         x1          x2 
-0.01147885 -0.07219898 

$eigen
eigen() decomposition
$values
[1] -0.3809044 -0.4711314

$vectors
         [,1]       [,2]
x1  0.3910698 -0.9203610
x2 -0.9203610 -0.3910698
\end{verbatim}

\begin{Shaded}
\begin{Highlighting}[]
\FunctionTok{canonical.path}\NormalTok{(rsmmod)}
\end{Highlighting}
\end{Shaded}

\begin{verbatim}
   dist     x1     x2 |  Temp     pH |   yhat
1  -5.0 -1.967  4.530 | 30.33 8.2650 | -8.295
2  -4.5 -1.771  4.069 | 32.29 8.0345 | -6.483
3  -4.0 -1.576  3.609 | 34.24 7.8045 | -4.865
4  -3.5 -1.380  3.149 | 36.20 7.5745 | -3.437
5  -3.0 -1.185  2.689 | 38.15 7.3445 | -2.200
6  -2.5 -0.989  2.229 | 40.11 7.1145 | -1.152
7  -2.0 -0.794  1.769 | 42.06 6.8845 | -0.296
8  -1.5 -0.598  1.308 | 44.02 6.6540 |  0.372
9  -1.0 -0.403  0.848 | 45.97 6.4240 |  0.848
10 -0.5 -0.207  0.388 | 47.93 6.1940 |  1.134
11  0.0 -0.011 -0.072 | 49.89 5.9640 |  1.229
12  0.5  0.184 -0.532 | 51.84 5.7340 |  1.134
13  1.0  0.380 -0.993 | 53.80 5.5035 |  0.848
14  1.5  0.575 -1.453 | 55.75 5.2735 |  0.372
15  2.0  0.771 -1.913 | 57.71 5.0435 | -0.295
16  2.5  0.966 -2.373 | 59.66 4.8135 | -1.151
17  3.0  1.162 -2.833 | 61.62 4.5835 | -2.199
18  3.5  1.357 -3.293 | 63.57 4.3535 | -3.436
19  4.0  1.553 -3.754 | 65.53 4.1230 | -4.867
20  4.5  1.748 -4.214 | 67.48 3.8930 | -6.485
21  5.0  1.944 -4.674 | 69.44 3.6630 | -8.294
\end{verbatim}

\subsection{visualize}\label{visualize}

\begin{Shaded}
\begin{Highlighting}[]
\FunctionTok{contour}\NormalTok{(rsmmod, }\SpecialCharTok{\textasciitilde{}}\NormalTok{ x1 }\SpecialCharTok{+}\NormalTok{ x2,}
        \AttributeTok{at =} \FunctionTok{canonical}\NormalTok{(rsmmod)}\SpecialCharTok{$}\NormalTok{xs,}
        \AttributeTok{image =}\NormalTok{ T, }\AttributeTok{nlevels =} \DecValTok{15}\NormalTok{, }\AttributeTok{decode =}\NormalTok{ T}
\NormalTok{)}
\end{Highlighting}
\end{Shaded}

\begin{center}\includegraphics[width=1\linewidth]{../Figs/rsm_plot-1} \end{center}

\begin{Shaded}
\begin{Highlighting}[]
\CommentTok{\# from Myfuns library}
\FunctionTok{my\_gg\_rsm}\NormalTok{(rsmmod, }\SpecialCharTok{\textasciitilde{}}\NormalTok{ x1 }\SpecialCharTok{+}\NormalTok{ x2,}
          \AttributeTok{at =} \FunctionTok{canonical}\NormalTok{(rsmmod)}\SpecialCharTok{$}\NormalTok{xs,}
          \AttributeTok{color =} \StringTok{\textquotesingle{}black\textquotesingle{}}\NormalTok{,}
          \AttributeTok{text\_size =} \DecValTok{3}\NormalTok{,}
          \AttributeTok{decode =}\NormalTok{T)}
\end{Highlighting}
\end{Shaded}

\begin{verbatim}
[[1]]
\end{verbatim}

\begin{center}\includegraphics[width=1\linewidth]{../Figs/rsm_plot-2} \end{center}

\begin{Shaded}
\begin{Highlighting}[]
\FunctionTok{persp}\NormalTok{(rsmmod, }\SpecialCharTok{\textasciitilde{}}\NormalTok{ x1 }\SpecialCharTok{+}\NormalTok{ x2,}
      \AttributeTok{at =} \FunctionTok{canonical}\NormalTok{(rsmmod)}\SpecialCharTok{$}\NormalTok{xs,}
      \AttributeTok{contours =} \FunctionTok{list}\NormalTok{(}\AttributeTok{z =} \StringTok{"bottom"}\NormalTok{), }\AttributeTok{ticktype =} \StringTok{"detailed"}\NormalTok{, }\AttributeTok{nticks =} \DecValTok{7}\NormalTok{,}
      \AttributeTok{shade =}\NormalTok{ .}\DecValTok{2}\NormalTok{,}
      \AttributeTok{col =} \FunctionTok{hcl.colors}\NormalTok{(}\DecValTok{20}\NormalTok{, }\StringTok{"Purple{-}Yellow"}\NormalTok{),}
      \AttributeTok{family =}\NormalTok{ mfont,}
      \AttributeTok{zlab =} \StringTok{"glu"}
\NormalTok{)}
\end{Highlighting}
\end{Shaded}

\begin{center}\includegraphics[width=1\linewidth]{../Figs/rsm_plot-3} \end{center}

\begin{Shaded}
\begin{Highlighting}[]
\CommentTok{\# library(rgl)}
\CommentTok{\# x1\_seq \textless{}{-} seq(min(rsmmod$data$x1), max(rsmmod$data$x1), length.out = 30)}
\CommentTok{\# x2\_seq \textless{}{-} seq(min(rsmmod$data$x2), max(rsmmod$data$x2), length.out = 30)}
\CommentTok{\# grid \textless{}{-} expand.grid(x1 = x1\_seq, x2 = x2\_seq)}
\CommentTok{\# grid$z \textless{}{-} predict(rsmmod, newdata = grid)}
\CommentTok{\# }
\CommentTok{\# \# Convert to matrix}
\CommentTok{\# z\_matrix \textless{}{-} matrix(grid$z, nrow = length(x1\_seq), ncol = length(x2\_seq))}
\CommentTok{\# }
\CommentTok{\# persp3d(x1\_seq, x2\_seq, z\_matrix,}
\CommentTok{\#   col = hcl.colors(20, "Purple{-}Yellow")[cut(z\_matrix, 20)],}
\CommentTok{\#   xlab = "Temperature",}
\CommentTok{\#   ylab = "pH",}
\CommentTok{\#   zlab = "Absorbance",}
\CommentTok{\#   ticktype = "detailed"}
\CommentTok{\# )}
\CommentTok{\# }
\CommentTok{\# \# Calculate contour lines}
\CommentTok{\# contour\_lines \textless{}{-} contourLines(x1\_seq, x2\_seq, z\_matrix)}
\CommentTok{\# }
\CommentTok{\# \# Add contours at the BOTTOM}
\CommentTok{\# for (i in seq\_along(contour\_lines)) \{}
\CommentTok{\#   lines3d(}
\CommentTok{\#     x = contour\_lines[[i]]$x,}
\CommentTok{\#     y = contour\_lines[[i]]$y,}
\CommentTok{\#     z = min(z\_matrix),}
\CommentTok{\#     col = "gray50",}
\CommentTok{\#     lwd = 1,}
\CommentTok{\#     alpha = 0.5}
\CommentTok{\#   )}
\CommentTok{\# \}}
\CommentTok{\# }
\CommentTok{\# \# Add contours ON the SURFACE}
\CommentTok{\# for (i in seq\_along(contour\_lines)) \{}
\CommentTok{\#   z\_vals \textless{}{-} rep(contour\_lines[[i]]$level, length(contour\_lines[[i]]$x))}
\CommentTok{\#   lines3d(}
\CommentTok{\#     x = contour\_lines[[i]]$x,}
\CommentTok{\#     y = contour\_lines[[i]]$y,}
\CommentTok{\#     z = z\_vals,}
\CommentTok{\#     col = "black",}
\CommentTok{\#     lwd = 2}
\CommentTok{\#   )}
\CommentTok{\# \}}
\CommentTok{\# \# Add canonical point}
\CommentTok{\# canon \textless{}{-} canonical(rsmmod)$xs}
\CommentTok{\# canon\_z \textless{}{-} predict(rsmmod, newdata = as.data.frame(t(canon)))}
\CommentTok{\# points3d(canon["x1"], canon["x2"], canon\_z,}
\CommentTok{\#   col = "red", size = 10}
\CommentTok{\# )}
\CommentTok{\# }
\CommentTok{\# \# Add label for canonical point}
\CommentTok{\# text3d(canon["x1"], canon["x2"], canon\_z * 1.1,}
\CommentTok{\#   texts = "Optimum", col = "red"}
\CommentTok{\# )}
\CommentTok{\# axes3d()}
\CommentTok{\# rgl.snapshot("my\_3d\_plot.png", fmt = "png")}
\end{Highlighting}
\end{Shaded}

\begin{Shaded}
\begin{Highlighting}[]
\CommentTok{\# knitr::include\_graphics("my\_3d\_plot.png")}
\end{Highlighting}
\end{Shaded}

\subsection{What is optimum condition and its predicted
response}\label{what-is-optimum-condition-and-its-predicted-response}

\begin{Shaded}
\begin{Highlighting}[]
\CommentTok{\# optimum condition}
\FunctionTok{code2val}\NormalTok{(}\FunctionTok{canonical}\NormalTok{((rsmmod))}\SpecialCharTok{$}\NormalTok{xs, }\AttributeTok{codings =} \FunctionTok{codings}\NormalTok{(codedrsm2))}
\end{Highlighting}
\end{Shaded}

\begin{verbatim}
     Temp        pH 
49.885212  5.963901 
\end{verbatim}

\begin{Shaded}
\begin{Highlighting}[]
\NormalTok{prediction }\OtherTok{\textless{}{-}} \FunctionTok{predict}\NormalTok{(rsmmod, }\AttributeTok{newdata  =} \FunctionTok{data.frame}\NormalTok{(}
  \AttributeTok{x1 =} \FunctionTok{canonical}\NormalTok{(rsmmod)}\SpecialCharTok{$}\NormalTok{xs[}\DecValTok{1}\NormalTok{],}
  \AttributeTok{x2 =} \FunctionTok{canonical}\NormalTok{(rsmmod)}\SpecialCharTok{$}\NormalTok{xs[}\DecValTok{2}\NormalTok{]),}
  \AttributeTok{interval =} \StringTok{"confidence"}\NormalTok{)}

\CommentTok{\#predicted response}
\FunctionTok{kable}\NormalTok{(prediction)}
\end{Highlighting}
\end{Shaded}

{\def\LTcaptype{none} % do not increment counter
\begin{longtable}[]{@{}lrrr@{}}
\toprule\noalign{}
& fit & lwr & upr \\
\midrule\noalign{}
\endhead
\bottomrule\noalign{}
\endlastfoot
x1 & 1.228853 & 1.181245 & 1.276462 \\
\end{longtable}
}

\section{Amylase activity on four different starch of different
concentration}\label{amylase-activity-on-four-different-starch-of-different-concentration}

\begin{itemize}
\tightlist
\item
  load and tidy the data
\end{itemize}

\begin{Shaded}
\begin{Highlighting}[]
\NormalTok{UV\_04\_08\_2006\_ }\OtherTok{\textless{}{-}} \FunctionTok{read\_excel}\NormalTok{(}\StringTok{"../Data/UV (04.08.2006) {-} completed.xlsx"}\NormalTok{, }
                             \AttributeTok{sheet =} \StringTok{"Sheet1"}\NormalTok{)}
\end{Highlighting}
\end{Shaded}

\begin{verbatim}
New names:
* `` -> `...7`
* `` -> `...8`
* `` -> `...9`
\end{verbatim}

\begin{Shaded}
\begin{Highlighting}[]
\NormalTok{UV }\OtherTok{\textless{}{-}}\NormalTok{ UV\_04\_08\_2006\_[}\DecValTok{1}\SpecialCharTok{:}\DecValTok{42}\NormalTok{,}\DecValTok{1}\SpecialCharTok{:}\DecValTok{6}\NormalTok{]}
\NormalTok{UV }\OtherTok{\textless{}{-}}\NormalTok{ UV }\SpecialCharTok{\%\textgreater{}\%} \FunctionTok{filter}\NormalTok{(}\SpecialCharTok{!}\FunctionTok{is.na}\NormalTok{(}\StringTok{\textasciigrave{}}\AttributeTok{0 min}\StringTok{\textasciigrave{}}\NormalTok{))}

\NormalTok{ns }\OtherTok{\textless{}{-}} \FunctionTok{names}\NormalTok{(UV)}
\NormalTok{ns[}\DecValTok{1}\NormalTok{] }\OtherTok{\textless{}{-}} \StringTok{"Starch"}
\FunctionTok{names}\NormalTok{(UV) }\OtherTok{\textless{}{-}}\NormalTok{ ns}

\NormalTok{UV }\OtherTok{\textless{}{-}}\NormalTok{ UV }\SpecialCharTok{\%\textgreater{}\%} \FunctionTok{filter}\NormalTok{(}\SpecialCharTok{!}\FunctionTok{str\_detect}\NormalTok{(}\StringTok{\textasciigrave{}}\AttributeTok{0 min}\StringTok{\textasciigrave{}}\NormalTok{, }\StringTok{"min"}\NormalTok{))}

\NormalTok{starch }\OtherTok{\textless{}{-}} \FunctionTok{c}\NormalTok{(}\StringTok{"Soluble Starch"}\NormalTok{,}\StringTok{"Topioca Starch"}\NormalTok{,}\StringTok{"Corn Starch"}\NormalTok{,}\StringTok{"Rice Starch"}\NormalTok{)}
\NormalTok{starch }\OtherTok{\textless{}{-}} \FunctionTok{c}\NormalTok{(}\FunctionTok{rep}\NormalTok{(starch, }\AttributeTok{each=}\DecValTok{2}\NormalTok{), }\StringTok{"Negative control"}\NormalTok{)}
\NormalTok{starch }\OtherTok{\textless{}{-}} \FunctionTok{rep}\NormalTok{(starch, }\AttributeTok{times =} \DecValTok{4}\NormalTok{)}
\NormalTok{UV[,}\DecValTok{1}\NormalTok{] }\OtherTok{\textless{}{-}}\NormalTok{ starch}

\NormalTok{lvl }\OtherTok{\textless{}{-}} \FunctionTok{c}\NormalTok{(}
  \StringTok{"Soluble Starch"}\NormalTok{,}
  \StringTok{"Corn Starch"}\NormalTok{,}
  \StringTok{"Rice Starch"}\NormalTok{,}
  \StringTok{"Topioca Starch"}\NormalTok{,}
  \StringTok{"Negative control"}
\NormalTok{)}

\NormalTok{UV2 }\OtherTok{\textless{}{-}}\NormalTok{ UV }\SpecialCharTok{\%\textgreater{}\%} 
  \FunctionTok{mutate}\NormalTok{(}\AttributeTok{sub\_conc=}\FunctionTok{rep}\NormalTok{(}\FunctionTok{c}\NormalTok{(}\DecValTok{1}\NormalTok{, }\FloatTok{0.5}\NormalTok{, }\FloatTok{1.5}\NormalTok{, }\DecValTok{2}\NormalTok{), }\AttributeTok{each =} \DecValTok{9}\NormalTok{)) }\SpecialCharTok{\%\textgreater{}\%} 
  \FunctionTok{pivot\_longer}\NormalTok{(}\DecValTok{2}\SpecialCharTok{:}\DecValTok{6}\NormalTok{,}
                    \AttributeTok{names\_to =} \StringTok{"Time"}\NormalTok{,}
                    \AttributeTok{values\_to =} \StringTok{"ABS"}\NormalTok{) }\SpecialCharTok{\%\textgreater{}\%} 
  \FunctionTok{mutate}\NormalTok{(}\AttributeTok{Time=}\FunctionTok{str\_remove}\NormalTok{(Time, }\StringTok{" min"}\NormalTok{),}
         \AttributeTok{Time =} \FunctionTok{as.double}\NormalTok{(Time)) }\SpecialCharTok{\%\textgreater{}\%} 
  \FunctionTok{mutate}\NormalTok{(}\AttributeTok{Starch =} \FunctionTok{factor}\NormalTok{(Starch,}
                         \AttributeTok{levels =}\NormalTok{ lvl),}
         \AttributeTok{ABS=}\FunctionTok{as.numeric}\NormalTok{(ABS)) }
\FunctionTok{kable}\NormalTok{(UV2)}
\end{Highlighting}
\end{Shaded}

{\def\LTcaptype{none} % do not increment counter
\begin{longtable}[]{@{}lrrr@{}}
\toprule\noalign{}
Starch & sub\_conc & Time & ABS \\
\midrule\noalign{}
\endhead
\bottomrule\noalign{}
\endlastfoot
Soluble Starch & 1.0 & 0 & 0.9021 \\
Soluble Starch & 1.0 & 15 & 0.4526 \\
Soluble Starch & 1.0 & 30 & 0.2145 \\
Soluble Starch & 1.0 & 45 & 0.0124 \\
Soluble Starch & 1.0 & 60 & 0.0102 \\
Soluble Starch & 1.0 & 0 & 0.8968 \\
Soluble Starch & 1.0 & 15 & 0.4469 \\
Soluble Starch & 1.0 & 30 & 0.2078 \\
Soluble Starch & 1.0 & 45 & 0.0111 \\
Soluble Starch & 1.0 & 60 & 0.0095 \\
Topioca Starch & 1.0 & 0 & 1.0618 \\
Topioca Starch & 1.0 & 15 & 0.6215 \\
Topioca Starch & 1.0 & 30 & 0.3574 \\
Topioca Starch & 1.0 & 45 & 0.1401 \\
Topioca Starch & 1.0 & 60 & 0.0935 \\
Topioca Starch & 1.0 & 0 & 1.0576 \\
Topioca Starch & 1.0 & 15 & 0.6188 \\
Topioca Starch & 1.0 & 30 & 0.3559 \\
Topioca Starch & 1.0 & 45 & 0.1368 \\
Topioca Starch & 1.0 & 60 & 0.0962 \\
Corn Starch & 1.0 & 0 & 0.9842 \\
Corn Starch & 1.0 & 15 & 0.5059 \\
Corn Starch & 1.0 & 30 & 0.2764 \\
Corn Starch & 1.0 & 45 & 0.0187 \\
Corn Starch & 1.0 & 60 & 0.0139 \\
Corn Starch & 1.0 & 0 & 0.9915 \\
Corn Starch & 1.0 & 15 & 0.5126 \\
Corn Starch & 1.0 & 30 & 0.2698 \\
Corn Starch & 1.0 & 45 & 0.0202 \\
Corn Starch & 1.0 & 60 & 0.0147 \\
Rice Starch & 1.0 & 0 & 1.1250 \\
Rice Starch & 1.0 & 15 & 0.6748 \\
Rice Starch & 1.0 & 30 & 0.3629 \\
Rice Starch & 1.0 & 45 & 0.0524 \\
Rice Starch & 1.0 & 60 & 0.0426 \\
Rice Starch & 1.0 & 0 & 1.3380 \\
Rice Starch & 1.0 & 15 & 0.6785 \\
Rice Starch & 1.0 & 30 & 0.3859 \\
Rice Starch & 1.0 & 45 & 0.0479 \\
Rice Starch & 1.0 & 60 & 0.0335 \\
Negative control & 1.0 & 0 & 0.9328 \\
Negative control & 1.0 & 15 & 0.9326 \\
Negative control & 1.0 & 30 & 0.9331 \\
Negative control & 1.0 & 45 & 0.9321 \\
Negative control & 1.0 & 60 & 0.9348 \\
Soluble Starch & 0.5 & 0 & 0.5124 \\
Soluble Starch & 0.5 & 15 & 0.2866 \\
Soluble Starch & 0.5 & 30 & 0.1428 \\
Soluble Starch & 0.5 & 45 & 0.0084 \\
Soluble Starch & 0.5 & 60 & 0.0068 \\
Soluble Starch & 0.5 & 0 & 0.5181 \\
Soluble Starch & 0.5 & 15 & 0.2818 \\
Soluble Starch & 0.5 & 30 & 0.1379 \\
Soluble Starch & 0.5 & 45 & 0.0097 \\
Soluble Starch & 0.5 & 60 & 0.0059 \\
Topioca Starch & 0.5 & 0 & 0.6128 \\
Topioca Starch & 0.5 & 15 & 0.3522 \\
Topioca Starch & 0.5 & 30 & 0.1998 \\
Topioca Starch & 0.5 & 45 & 0.0825 \\
Topioca Starch & 0.5 & 60 & 0.0514 \\
Topioca Starch & 0.5 & 0 & 0.6085 \\
Topioca Starch & 0.5 & 15 & 0.3484 \\
Topioca Starch & 0.5 & 30 & 0.1948 \\
Topioca Starch & 0.5 & 45 & 0.0796 \\
Topioca Starch & 0.5 & 60 & 0.0478 \\
Corn Starch & 0.5 & 0 & 0.5642 \\
Corn Starch & 0.5 & 15 & 0.2928 \\
Corn Starch & 0.5 & 30 & 0.1546 \\
Corn Starch & 0.5 & 45 & 0.0104 \\
Corn Starch & 0.5 & 60 & 0.0082 \\
Corn Starch & 0.5 & 0 & 0.5695 \\
Corn Starch & 0.5 & 15 & 0.2964 \\
Corn Starch & 0.5 & 30 & 0.1498 \\
Corn Starch & 0.5 & 45 & 0.0118 \\
Corn Starch & 0.5 & 60 & 0.0071 \\
Rice Starch & 0.5 & 0 & 0.7184 \\
Rice Starch & 0.5 & 15 & 0.4562 \\
Rice Starch & 0.5 & 30 & 0.2485 \\
Rice Starch & 0.5 & 45 & 0.0368 \\
Rice Starch & 0.5 & 60 & 0.0245 \\
Rice Starch & 0.5 & 0 & 0.7246 \\
Rice Starch & 0.5 & 15 & 0.4628 \\
Rice Starch & 0.5 & 30 & 0.2554 \\
Rice Starch & 0.5 & 45 & 0.0412 \\
Rice Starch & 0.5 & 60 & 0.0281 \\
Negative control & 0.5 & 0 & 0.4688 \\
Negative control & 0.5 & 15 & 0.4682 \\
Negative control & 0.5 & 30 & 0.4695 \\
Negative control & 0.5 & 45 & 0.4689 \\
Negative control & 0.5 & 60 & 0.4696 \\
Soluble Starch & 1.5 & 0 & 1.2486 \\
Soluble Starch & 1.5 & 15 & 0.6735 \\
Soluble Starch & 1.5 & 30 & 0.3387 \\
Soluble Starch & 1.5 & 45 & 0.0286 \\
Soluble Starch & 1.5 & 60 & 0.0198 \\
Soluble Starch & 1.5 & 0 & 1.2418 \\
Soluble Starch & 1.5 & 15 & 0.6658 \\
Soluble Starch & 1.5 & 30 & 0.3314 \\
Soluble Starch & 1.5 & 45 & 0.0259 \\
Soluble Starch & 1.5 & 60 & 0.0176 \\
Topioca Starch & 1.5 & 0 & 1.4328 \\
Topioca Starch & 1.5 & 15 & 0.8468 \\
Topioca Starch & 1.5 & 30 & 0.4986 \\
Topioca Starch & 1.5 & 45 & 0.2068 \\
Topioca Starch & 1.5 & 60 & 0.1428 \\
Topioca Starch & 1.5 & 0 & 1.4268 \\
Topioca Starch & 1.5 & 15 & 0.8385 \\
Topioca Starch & 1.5 & 30 & 0.4912 \\
Topioca Starch & 1.5 & 45 & 0.1985 \\
Topioca Starch & 1.5 & 60 & 0.1362 \\
Corn Starch & 1.5 & 0 & 1.3258 \\
Corn Starch & 1.5 & 15 & 0.7165 \\
Corn Starch & 1.5 & 30 & 0.3868 \\
Corn Starch & 1.5 & 45 & 0.0352 \\
Corn Starch & 1.5 & 60 & 0.0248 \\
Corn Starch & 1.5 & 0 & 1.3342 \\
Corn Starch & 1.5 & 15 & 0.7248 \\
Corn Starch & 1.5 & 30 & 0.3796 \\
Corn Starch & 1.5 & 45 & 0.0358 \\
Corn Starch & 1.5 & 60 & 0.0221 \\
Rice Starch & 1.5 & 0 & 1.4862 \\
Rice Starch & 1.5 & 15 & 0.9386 \\
Rice Starch & 1.5 & 30 & 0.5568 \\
Rice Starch & 1.5 & 45 & 0.1126 \\
Rice Starch & 1.5 & 60 & 0.0826 \\
Rice Starch & 1.5 & 0 & 1.4988 \\
Rice Starch & 1.5 & 15 & 0.9462 \\
Rice Starch & 1.5 & 30 & 0.5486 \\
Rice Starch & 1.5 & 45 & 0.1048 \\
Rice Starch & 1.5 & 60 & 0.0765 \\
Negative control & 1.5 & 0 & 1.4082 \\
Negative control & 1.5 & 15 & 1.4076 \\
Negative control & 1.5 & 30 & 1.4078 \\
Negative control & 1.5 & 45 & 1.4080 \\
Negative control & 1.5 & 60 & 1.4070 \\
Soluble Starch & 2.0 & 0 & 1.5624 \\
Soluble Starch & 2.0 & 15 & 0.8426 \\
Soluble Starch & 2.0 & 30 & 0.4528 \\
Soluble Starch & 2.0 & 45 & 0.0412 \\
Soluble Starch & 2.0 & 60 & 0.0284 \\
Soluble Starch & 2.0 & 0 & 1.5558 \\
Soluble Starch & 2.0 & 15 & 0.8354 \\
Soluble Starch & 2.0 & 30 & 0.4465 \\
Soluble Starch & 2.0 & 45 & 0.0376 \\
Soluble Starch & 2.0 & 60 & 0.0259 \\
Topioca Starch & 2.0 & 0 & 1.7385 \\
Topioca Starch & 2.0 & 15 & 1.0224 \\
Topioca Starch & 2.0 & 30 & 0.6254 \\
Topioca Starch & 2.0 & 45 & 0.2826 \\
Topioca Starch & 2.0 & 60 & 0.1885 \\
Topioca Starch & 2.0 & 0 & 1.7316 \\
Topioca Starch & 2.0 & 15 & 1.0158 \\
Topioca Starch & 2.0 & 30 & 0.6187 \\
Topioca Starch & 2.0 & 45 & 0.2758 \\
Topioca Starch & 2.0 & 60 & 0.1824 \\
Corn Starch & 2.0 & 0 & 1.6488 \\
Corn Starch & 2.0 & 15 & 0.8958 \\
Corn Starch & 2.0 & 30 & 0.5186 \\
Corn Starch & 2.0 & 45 & 0.0485 \\
Corn Starch & 2.0 & 60 & 0.0354 \\
Corn Starch & 2.0 & 0 & 1.6554 \\
Corn Starch & 2.0 & 15 & 0.8886 \\
Corn Starch & 2.0 & 30 & 0.5114 \\
Corn Starch & 2.0 & 45 & 0.0458 \\
Corn Starch & 2.0 & 60 & 0.0328 \\
Rice Starch & 2.0 & 0 & 1.8624 \\
Rice Starch & 2.0 & 15 & 1.1842 \\
Rice Starch & 2.0 & 30 & 0.7128 \\
Rice Starch & 2.0 & 45 & 0.1658 \\
Rice Starch & 2.0 & 60 & 0.1184 \\
Rice Starch & 2.0 & 0 & 1.8756 \\
Rice Starch & 2.0 & 15 & 1.1786 \\
Rice Starch & 2.0 & 30 & 0.7056 \\
Rice Starch & 2.0 & 45 & 0.1586 \\
Rice Starch & 2.0 & 60 & 0.1126 \\
Negative control & 2.0 & 0 & 1.8658 \\
Negative control & 2.0 & 15 & 1.8646 \\
Negative control & 2.0 & 30 & 1.8662 \\
Negative control & 2.0 & 45 & 1.8655 \\
Negative control & 2.0 & 60 & 1.8668 \\
\end{longtable}
}

\subsection{Visualizing the enzyme
activity}\label{visualizing-the-enzyme-activity}

\begin{Shaded}
\begin{Highlighting}[]
\NormalTok{UV2 }\SpecialCharTok{\%\textgreater{}\%} 
  \CommentTok{\#filter(Starch!="Negative control")\%\textgreater{}\% }
  \FunctionTok{ggplot}\NormalTok{(}\FunctionTok{aes}\NormalTok{(}\AttributeTok{x=}\NormalTok{Time, }\AttributeTok{y=}\NormalTok{ABS, }\AttributeTok{color =}\NormalTok{ Starch))}\SpecialCharTok{+}
  \FunctionTok{geom\_point}\NormalTok{()}\SpecialCharTok{+}
  \FunctionTok{coord\_cartesian}\NormalTok{(}\AttributeTok{ylim =} \FunctionTok{c}\NormalTok{(}\DecValTok{0}\NormalTok{,}\DecValTok{2}\NormalTok{))}\SpecialCharTok{+}
  \FunctionTok{geom\_smooth}\NormalTok{(}\AttributeTok{method =} \StringTok{\textquotesingle{}lm\textquotesingle{}}\NormalTok{, }\AttributeTok{se =}\NormalTok{ F,}
              \AttributeTok{formula =}\NormalTok{ y}\SpecialCharTok{\textasciitilde{}}\FunctionTok{poly}\NormalTok{(x, }\DecValTok{2}\NormalTok{, }\AttributeTok{raw =}\NormalTok{T))}\SpecialCharTok{+}
  \FunctionTok{scale\_x\_continuous}\NormalTok{(}\AttributeTok{breaks =} \FunctionTok{seq}\NormalTok{(}\DecValTok{0}\NormalTok{, }\DecValTok{60}\NormalTok{, }\DecValTok{15}\NormalTok{))}\SpecialCharTok{+}
  \FunctionTok{scale\_y\_continuous}\NormalTok{(}\AttributeTok{breaks =} \FunctionTok{seq}\NormalTok{(}\DecValTok{0}\NormalTok{, }\DecValTok{2}\NormalTok{, }\FloatTok{0.5}\NormalTok{))}\SpecialCharTok{+}
  \FunctionTok{scale\_color\_manual}\NormalTok{(}
    \AttributeTok{breaks =}\NormalTok{ lvl,}
    \AttributeTok{values =} \FunctionTok{c}\NormalTok{(}\StringTok{"green"}\NormalTok{,}
               \StringTok{"yellow"}\NormalTok{,}
               \StringTok{"black"}\NormalTok{,}
               \StringTok{"blue"}\NormalTok{,}
               \StringTok{"red"}\NormalTok{)}
\NormalTok{  )}\SpecialCharTok{+}
  \FunctionTok{labs}\NormalTok{(}
    \AttributeTok{x =} \StringTok{"Time (minutes)"}\NormalTok{,}
    \AttributeTok{y =} \StringTok{"Absorbance at 620 nm"}
\NormalTok{  )}\SpecialCharTok{+}
  \FunctionTok{facet\_grid}\NormalTok{(.}\SpecialCharTok{\textasciitilde{}}\NormalTok{sub\_conc,}
             \AttributeTok{labeller =} \FunctionTok{as\_labeller}\NormalTok{(}\ControlFlowTok{function}\NormalTok{(x)\{}
               \FunctionTok{paste}\NormalTok{(x, }\StringTok{"\%"}\NormalTok{)}
\NormalTok{             \}))}\SpecialCharTok{+}
  \FunctionTok{theme\_classic}\NormalTok{()}\SpecialCharTok{+}
  \FunctionTok{theme}\NormalTok{(}
    \AttributeTok{text =} \FunctionTok{element\_text}\NormalTok{(}\AttributeTok{family=}\NormalTok{ mfont,}
                        \AttributeTok{size =} \DecValTok{12}\NormalTok{,}
                        \AttributeTok{color =} \StringTok{\textquotesingle{}black\textquotesingle{}}\NormalTok{),}
    \AttributeTok{legend.position =} \StringTok{"bottom"}\NormalTok{,}
    \AttributeTok{strip.background =} \FunctionTok{element\_rect}\NormalTok{(}\AttributeTok{color =} \StringTok{"grey85"}\NormalTok{)}
\NormalTok{    )}
\end{Highlighting}
\end{Shaded}

\begin{center}\includegraphics[width=1\linewidth]{../Figs/enz_activity-1} \end{center}

\begin{Shaded}
\begin{Highlighting}[]
\NormalTok{UV2}\FloatTok{.1} \OtherTok{\textless{}{-}}\NormalTok{ UV2 }\SpecialCharTok{\%\textgreater{}\%} \FunctionTok{reframe}\NormalTok{(}\AttributeTok{Absmean=}\FunctionTok{mean}\NormalTok{(ABS),}
                \AttributeTok{.by =} \FunctionTok{c}\NormalTok{(Time, sub\_conc, Starch)) }\SpecialCharTok{\%\textgreater{}\%} 
  \CommentTok{\# group\_by(Starch, Time, sub\_conc) \%\textgreater{}\% }
  \FunctionTok{mutate}\NormalTok{(}\AttributeTok{ABSper=}\NormalTok{(Absmean)}\SpecialCharTok{/}\FunctionTok{first}\NormalTok{(Absmean)}\SpecialCharTok{*}\DecValTok{100}\NormalTok{,}
         \AttributeTok{.by =} \FunctionTok{c}\NormalTok{(sub\_conc, Starch)) }
\FunctionTok{kable}\NormalTok{(UV2}\FloatTok{.1}\NormalTok{)}
\end{Highlighting}
\end{Shaded}

{\def\LTcaptype{none} % do not increment counter
\begin{longtable}[]{@{}rrlrr@{}}
\toprule\noalign{}
Time & sub\_conc & Starch & Absmean & ABSper \\
\midrule\noalign{}
\endhead
\bottomrule\noalign{}
\endlastfoot
0 & 1.0 & Soluble Starch & 0.89945 & 100.000000 \\
15 & 1.0 & Soluble Starch & 0.44975 & 50.002780 \\
30 & 1.0 & Soluble Starch & 0.21115 & 23.475457 \\
45 & 1.0 & Soluble Starch & 0.01175 & 1.306354 \\
60 & 1.0 & Soluble Starch & 0.00985 & 1.095114 \\
0 & 1.0 & Topioca Starch & 1.05970 & 100.000000 \\
15 & 1.0 & Topioca Starch & 0.62015 & 58.521280 \\
30 & 1.0 & Topioca Starch & 0.35665 & 33.655752 \\
45 & 1.0 & Topioca Starch & 0.13845 & 13.065018 \\
60 & 1.0 & Topioca Starch & 0.09485 & 8.950646 \\
0 & 1.0 & Corn Starch & 0.98785 & 100.000000 \\
15 & 1.0 & Corn Starch & 0.50925 & 51.551349 \\
30 & 1.0 & Corn Starch & 0.27310 & 27.645898 \\
45 & 1.0 & Corn Starch & 0.01945 & 1.968922 \\
60 & 1.0 & Corn Starch & 0.01430 & 1.447588 \\
0 & 1.0 & Rice Starch & 1.23150 & 100.000000 \\
15 & 1.0 & Rice Starch & 0.67665 & 54.945189 \\
30 & 1.0 & Rice Starch & 0.37440 & 30.401949 \\
45 & 1.0 & Rice Starch & 0.05015 & 4.072270 \\
60 & 1.0 & Rice Starch & 0.03805 & 3.089728 \\
0 & 1.0 & Negative control & 0.93280 & 100.000000 \\
15 & 1.0 & Negative control & 0.93260 & 99.978559 \\
30 & 1.0 & Negative control & 0.93310 & 100.032161 \\
45 & 1.0 & Negative control & 0.93210 & 99.924957 \\
60 & 1.0 & Negative control & 0.93480 & 100.214408 \\
0 & 0.5 & Soluble Starch & 0.51525 & 100.000000 \\
15 & 0.5 & Soluble Starch & 0.28420 & 55.157690 \\
30 & 0.5 & Soluble Starch & 0.14035 & 27.239204 \\
45 & 0.5 & Soluble Starch & 0.00905 & 1.756429 \\
60 & 0.5 & Soluble Starch & 0.00635 & 1.232411 \\
0 & 0.5 & Topioca Starch & 0.61065 & 100.000000 \\
15 & 0.5 & Topioca Starch & 0.35030 & 57.365103 \\
30 & 0.5 & Topioca Starch & 0.19730 & 32.309834 \\
45 & 0.5 & Topioca Starch & 0.08105 & 13.272742 \\
60 & 0.5 & Topioca Starch & 0.04960 & 8.122492 \\
0 & 0.5 & Corn Starch & 0.56685 & 100.000000 \\
15 & 0.5 & Corn Starch & 0.29460 & 51.971421 \\
30 & 0.5 & Corn Starch & 0.15220 & 26.850137 \\
45 & 0.5 & Corn Starch & 0.01110 & 1.958190 \\
60 & 0.5 & Corn Starch & 0.00765 & 1.349563 \\
0 & 0.5 & Rice Starch & 0.72150 & 100.000000 \\
15 & 0.5 & Rice Starch & 0.45950 & 63.686764 \\
30 & 0.5 & Rice Starch & 0.25195 & 34.920305 \\
45 & 0.5 & Rice Starch & 0.03900 & 5.405405 \\
60 & 0.5 & Rice Starch & 0.02630 & 3.645184 \\
0 & 0.5 & Negative control & 0.46880 & 100.000000 \\
15 & 0.5 & Negative control & 0.46820 & 99.872014 \\
30 & 0.5 & Negative control & 0.46950 & 100.149317 \\
45 & 0.5 & Negative control & 0.46890 & 100.021331 \\
60 & 0.5 & Negative control & 0.46960 & 100.170648 \\
0 & 1.5 & Soluble Starch & 1.24520 & 100.000000 \\
15 & 1.5 & Soluble Starch & 0.66965 & 53.778509 \\
30 & 1.5 & Soluble Starch & 0.33505 & 26.907324 \\
45 & 1.5 & Soluble Starch & 0.02725 & 2.188404 \\
60 & 1.5 & Soluble Starch & 0.01870 & 1.501767 \\
0 & 1.5 & Topioca Starch & 1.42980 & 100.000000 \\
15 & 1.5 & Topioca Starch & 0.84265 & 58.934816 \\
30 & 1.5 & Topioca Starch & 0.49490 & 34.613233 \\
45 & 1.5 & Topioca Starch & 0.20265 & 14.173311 \\
60 & 1.5 & Topioca Starch & 0.13950 & 9.756609 \\
0 & 1.5 & Corn Starch & 1.33000 & 100.000000 \\
15 & 1.5 & Corn Starch & 0.72065 & 54.184210 \\
30 & 1.5 & Corn Starch & 0.38320 & 28.812030 \\
45 & 1.5 & Corn Starch & 0.03550 & 2.669173 \\
60 & 1.5 & Corn Starch & 0.02345 & 1.763158 \\
0 & 1.5 & Rice Starch & 1.49250 & 100.000000 \\
15 & 1.5 & Rice Starch & 0.94240 & 63.142379 \\
30 & 1.5 & Rice Starch & 0.55270 & 37.031826 \\
45 & 1.5 & Rice Starch & 0.10870 & 7.283082 \\
60 & 1.5 & Rice Starch & 0.07955 & 5.329983 \\
0 & 1.5 & Negative control & 1.40820 & 100.000000 \\
15 & 1.5 & Negative control & 1.40760 & 99.957392 \\
30 & 1.5 & Negative control & 1.40780 & 99.971595 \\
45 & 1.5 & Negative control & 1.40800 & 99.985798 \\
60 & 1.5 & Negative control & 1.40700 & 99.914785 \\
0 & 2.0 & Soluble Starch & 1.55910 & 100.000000 \\
15 & 2.0 & Soluble Starch & 0.83900 & 53.813097 \\
30 & 2.0 & Soluble Starch & 0.44965 & 28.840357 \\
45 & 2.0 & Soluble Starch & 0.03940 & 2.527099 \\
60 & 2.0 & Soluble Starch & 0.02715 & 1.741389 \\
0 & 2.0 & Topioca Starch & 1.73505 & 100.000000 \\
15 & 2.0 & Topioca Starch & 1.01910 & 58.736060 \\
30 & 2.0 & Topioca Starch & 0.62205 & 35.851993 \\
45 & 2.0 & Topioca Starch & 0.27920 & 16.091755 \\
60 & 2.0 & Topioca Starch & 0.18545 & 10.688453 \\
0 & 2.0 & Corn Starch & 1.65210 & 100.000000 \\
15 & 2.0 & Corn Starch & 0.89220 & 54.003995 \\
30 & 2.0 & Corn Starch & 0.51500 & 31.172447 \\
45 & 2.0 & Corn Starch & 0.04715 & 2.853944 \\
60 & 2.0 & Corn Starch & 0.03410 & 2.064040 \\
0 & 2.0 & Rice Starch & 1.86900 & 100.000000 \\
15 & 2.0 & Rice Starch & 1.18140 & 63.210273 \\
30 & 2.0 & Rice Starch & 0.70920 & 37.945425 \\
45 & 2.0 & Rice Starch & 0.16220 & 8.678438 \\
60 & 2.0 & Rice Starch & 0.11550 & 6.179775 \\
0 & 2.0 & Negative control & 1.86580 & 100.000000 \\
15 & 2.0 & Negative control & 1.86460 & 99.935684 \\
30 & 2.0 & Negative control & 1.86620 & 100.021439 \\
45 & 2.0 & Negative control & 1.86550 & 99.983921 \\
60 & 2.0 & Negative control & 1.86680 & 100.053596 \\
\end{longtable}
}

\begin{Shaded}
\begin{Highlighting}[]
\NormalTok{UV2}\FloatTok{.1}\SpecialCharTok{\%\textgreater{}\%} 
  \FunctionTok{ggplot}\NormalTok{(}\FunctionTok{aes}\NormalTok{(}\AttributeTok{x=}\NormalTok{Time, }\AttributeTok{y=}\NormalTok{ABSper, }\AttributeTok{color =}\NormalTok{ Starch))}\SpecialCharTok{+}
  \CommentTok{\#geom\_smooth()}
  \FunctionTok{geom\_smooth}\NormalTok{(}\AttributeTok{method =} \StringTok{\textquotesingle{}lm\textquotesingle{}}\NormalTok{, }\AttributeTok{se =}\NormalTok{ F,}
              \AttributeTok{formula =}\NormalTok{ y}\SpecialCharTok{\textasciitilde{}}\FunctionTok{poly}\NormalTok{(x, }\DecValTok{2}\NormalTok{, }\AttributeTok{raw =}\NormalTok{T))}\SpecialCharTok{+}
  \FunctionTok{geom\_point}\NormalTok{()}\SpecialCharTok{+}
  \FunctionTok{scale\_x\_continuous}\NormalTok{(}\AttributeTok{breaks =} \FunctionTok{seq}\NormalTok{(}\DecValTok{0}\NormalTok{, }\DecValTok{60}\NormalTok{, }\DecValTok{15}\NormalTok{))}\SpecialCharTok{+}
  \FunctionTok{scale\_y\_continuous}\NormalTok{(}\AttributeTok{breaks =} \FunctionTok{seq}\NormalTok{(}\DecValTok{0}\NormalTok{, }\DecValTok{100}\NormalTok{, }\DecValTok{20}\NormalTok{))}\SpecialCharTok{+}
  \FunctionTok{scale\_color\_manual}\NormalTok{(}
    \AttributeTok{breaks =}\NormalTok{ lvl,}
    \AttributeTok{values =} \FunctionTok{c}\NormalTok{(}\StringTok{"green"}\NormalTok{,}
               \StringTok{"yellow"}\NormalTok{,}
               \StringTok{"black"}\NormalTok{,}
               \StringTok{"blue"}\NormalTok{,}
               \StringTok{"red"}\NormalTok{)}
\NormalTok{  )}\SpecialCharTok{+}
  \FunctionTok{labs}\NormalTok{(}
    \AttributeTok{x =} \StringTok{"Time (minutes)"}\NormalTok{,}
    \AttributeTok{y =} \StringTok{"Reduction percentage of}
\StringTok{    Absorbance at 620 nm"}
\NormalTok{  )}\SpecialCharTok{+}
  \FunctionTok{facet\_grid}\NormalTok{(.}\SpecialCharTok{\textasciitilde{}}\NormalTok{sub\_conc,}
             \AttributeTok{labeller =} \FunctionTok{as\_labeller}\NormalTok{(}\ControlFlowTok{function}\NormalTok{(x)\{}
               \FunctionTok{paste}\NormalTok{(x, }\StringTok{"\%"}\NormalTok{)}
\NormalTok{             \}))}\SpecialCharTok{+}
  \FunctionTok{theme\_classic}\NormalTok{()}\SpecialCharTok{+}
  \FunctionTok{theme}\NormalTok{(}
    \AttributeTok{text =} \FunctionTok{element\_text}\NormalTok{(}\AttributeTok{family=}\NormalTok{ mfont,}
                        \AttributeTok{size =} \DecValTok{12}\NormalTok{,}
                        \AttributeTok{color =} \StringTok{\textquotesingle{}black\textquotesingle{}}\NormalTok{),}
    \AttributeTok{legend.position =} \StringTok{"bottom"}\NormalTok{,}
    \AttributeTok{strip.background =} \FunctionTok{element\_rect}\NormalTok{(}\AttributeTok{color =} \StringTok{"grey85"}\NormalTok{)}
    \CommentTok{\# ,}
    \CommentTok{\# panel.grid.major = element\_line(}
    \CommentTok{\#   color = "grey80",}
    \CommentTok{\#   linewidth = 0.3,}
    \CommentTok{\#   linetype = 3}
\NormalTok{  )}
\end{Highlighting}
\end{Shaded}

\begin{center}\includegraphics[width=1\linewidth]{../Figs/enz_activity_percentage_by_sub_conc-1} \end{center}

\section{Numerical estimation of the curvature of the enzyme
activity}\label{numerical-estimation-of-the-curvature-of-the-enzyme-activity}

\begin{Shaded}
\begin{Highlighting}[]
\NormalTok{mod }\OtherTok{\textless{}{-}} \FunctionTok{lm}\NormalTok{(ABS}\SpecialCharTok{\textasciitilde{}}\NormalTok{Starch}\SpecialCharTok{+}\FunctionTok{poly}\NormalTok{(Time, }\DecValTok{2}\NormalTok{, }\AttributeTok{raw =}\NormalTok{T), }\AttributeTok{data=}\NormalTok{ UV2 }\SpecialCharTok{\%\textgreater{}\%} 
            \FunctionTok{filter}\NormalTok{(Starch}\SpecialCharTok{!=}\StringTok{"Negative control"}\NormalTok{) )}
\FunctionTok{summary}\NormalTok{(mod)}
\end{Highlighting}
\end{Shaded}

\begin{verbatim}

Call:
lm(formula = ABS ~ Starch + poly(Time, 2, raw = T), data = UV2 %>% 
    filter(Starch != "Negative control"))

Residuals:
    Min      1Q  Median      3Q     Max 
-0.6249 -0.1074 -0.0081  0.1042  0.6136 

Coefficients:
                          Estimate Std. Error t value Pr(>|t|)    
(Intercept)              1.095e+00  5.011e-02  21.859  < 2e-16 ***
StarchCorn Starch        3.612e-02  5.200e-02   0.695  0.48832    
StarchRice Starch        1.667e-01  5.200e-02   3.207  0.00163 ** 
StarchTopioca Starch     1.381e-01  5.200e-02   2.656  0.00875 ** 
poly(Time, 2, raw = T)1 -3.744e-02  3.055e-03 -12.253  < 2e-16 ***
poly(Time, 2, raw = T)2  3.076e-04  4.883e-05   6.299    3e-09 ***
---
Signif. codes:  0 '***' 0.001 '**' 0.01 '*' 0.05 '.' 0.1 ' ' 1

Residual standard error: 0.2325 on 154 degrees of freedom
Multiple R-squared:  0.776, Adjusted R-squared:  0.7688 
F-statistic: 106.7 on 5 and 154 DF,  p-value: < 2.2e-16
\end{verbatim}

\begin{Shaded}
\begin{Highlighting}[]
\CommentTok{\# estimattion vs lab{-}tested value}
\NormalTok{uv3 }\OtherTok{\textless{}{-}}\NormalTok{ UV2 }\SpecialCharTok{\%\textgreater{}\%} \FunctionTok{filter}\NormalTok{(Starch }\SpecialCharTok{!=} \StringTok{"Negative control"}\NormalTok{)}
\NormalTok{uv3}\SpecialCharTok{$}\NormalTok{fitted }\OtherTok{\textless{}{-}} \FunctionTok{fitted}\NormalTok{(mod)}
\NormalTok{uv3}\SpecialCharTok{$}\NormalTok{residual }\OtherTok{\textless{}{-}} \FunctionTok{resid}\NormalTok{(mod)}
\FunctionTok{kable}\NormalTok{(uv3)}
\end{Highlighting}
\end{Shaded}

{\def\LTcaptype{none} % do not increment counter
\begin{longtable}[]{@{}lrrrrr@{}}
\toprule\noalign{}
Starch & sub\_conc & Time & ABS & fitted & residual \\
\midrule\noalign{}
\endhead
\bottomrule\noalign{}
\endlastfoot
Soluble Starch & 1.0 & 0 & 0.9021 & 1.0952586 & -0.1931586 \\
Soluble Starch & 1.0 & 15 & 0.4526 & 0.6029145 & -0.1503145 \\
Soluble Starch & 1.0 & 30 & 0.2145 & 0.2489677 & -0.0344677 \\
Soluble Starch & 1.0 & 45 & 0.0124 & 0.0334182 & -0.0210182 \\
Soluble Starch & 1.0 & 60 & 0.0102 & -0.0437339 & 0.0539339 \\
Soluble Starch & 1.0 & 0 & 0.8968 & 1.0952586 & -0.1984586 \\
Soluble Starch & 1.0 & 15 & 0.4469 & 0.6029145 & -0.1560145 \\
Soluble Starch & 1.0 & 30 & 0.2078 & 0.2489677 & -0.0411677 \\
Soluble Starch & 1.0 & 45 & 0.0111 & 0.0334182 & -0.0223182 \\
Soluble Starch & 1.0 & 60 & 0.0095 & -0.0437339 & 0.0532339 \\
Topioca Starch & 1.0 & 0 & 1.0618 & 1.2333461 & -0.1715461 \\
Topioca Starch & 1.0 & 15 & 0.6215 & 0.7410020 & -0.1195020 \\
Topioca Starch & 1.0 & 30 & 0.3574 & 0.3870552 & -0.0296552 \\
Topioca Starch & 1.0 & 45 & 0.1401 & 0.1715057 & -0.0314057 \\
Topioca Starch & 1.0 & 60 & 0.0935 & 0.0943536 & -0.0008536 \\
Topioca Starch & 1.0 & 0 & 1.0576 & 1.2333461 & -0.1757461 \\
Topioca Starch & 1.0 & 15 & 0.6188 & 0.7410020 & -0.1222020 \\
Topioca Starch & 1.0 & 30 & 0.3559 & 0.3870552 & -0.0311552 \\
Topioca Starch & 1.0 & 45 & 0.1368 & 0.1715057 & -0.0347057 \\
Topioca Starch & 1.0 & 60 & 0.0962 & 0.0943536 & 0.0018464 \\
Corn Starch & 1.0 & 0 & 0.9842 & 1.1313786 & -0.1471786 \\
Corn Starch & 1.0 & 15 & 0.5059 & 0.6390345 & -0.1331345 \\
Corn Starch & 1.0 & 30 & 0.2764 & 0.2850877 & -0.0086877 \\
Corn Starch & 1.0 & 45 & 0.0187 & 0.0695382 & -0.0508382 \\
Corn Starch & 1.0 & 60 & 0.0139 & -0.0076139 & 0.0215139 \\
Corn Starch & 1.0 & 0 & 0.9915 & 1.1313786 & -0.1398786 \\
Corn Starch & 1.0 & 15 & 0.5126 & 0.6390345 & -0.1264345 \\
Corn Starch & 1.0 & 30 & 0.2698 & 0.2850877 & -0.0152877 \\
Corn Starch & 1.0 & 45 & 0.0202 & 0.0695382 & -0.0493382 \\
Corn Starch & 1.0 & 60 & 0.0147 & -0.0076139 & 0.0223139 \\
Rice Starch & 1.0 & 0 & 1.1250 & 1.2620011 & -0.1370011 \\
Rice Starch & 1.0 & 15 & 0.6748 & 0.7696570 & -0.0948570 \\
Rice Starch & 1.0 & 30 & 0.3629 & 0.4157102 & -0.0528102 \\
Rice Starch & 1.0 & 45 & 0.0524 & 0.2001607 & -0.1477607 \\
Rice Starch & 1.0 & 60 & 0.0426 & 0.1230086 & -0.0804086 \\
Rice Starch & 1.0 & 0 & 1.3380 & 1.2620011 & 0.0759989 \\
Rice Starch & 1.0 & 15 & 0.6785 & 0.7696570 & -0.0911570 \\
Rice Starch & 1.0 & 30 & 0.3859 & 0.4157102 & -0.0298102 \\
Rice Starch & 1.0 & 45 & 0.0479 & 0.2001607 & -0.1522607 \\
Rice Starch & 1.0 & 60 & 0.0335 & 0.1230086 & -0.0895086 \\
Soluble Starch & 0.5 & 0 & 0.5124 & 1.0952586 & -0.5828586 \\
Soluble Starch & 0.5 & 15 & 0.2866 & 0.6029145 & -0.3163145 \\
Soluble Starch & 0.5 & 30 & 0.1428 & 0.2489677 & -0.1061677 \\
Soluble Starch & 0.5 & 45 & 0.0084 & 0.0334182 & -0.0250182 \\
Soluble Starch & 0.5 & 60 & 0.0068 & -0.0437339 & 0.0505339 \\
Soluble Starch & 0.5 & 0 & 0.5181 & 1.0952586 & -0.5771586 \\
Soluble Starch & 0.5 & 15 & 0.2818 & 0.6029145 & -0.3211145 \\
Soluble Starch & 0.5 & 30 & 0.1379 & 0.2489677 & -0.1110677 \\
Soluble Starch & 0.5 & 45 & 0.0097 & 0.0334182 & -0.0237182 \\
Soluble Starch & 0.5 & 60 & 0.0059 & -0.0437339 & 0.0496339 \\
Topioca Starch & 0.5 & 0 & 0.6128 & 1.2333461 & -0.6205461 \\
Topioca Starch & 0.5 & 15 & 0.3522 & 0.7410020 & -0.3888020 \\
Topioca Starch & 0.5 & 30 & 0.1998 & 0.3870552 & -0.1872552 \\
Topioca Starch & 0.5 & 45 & 0.0825 & 0.1715057 & -0.0890057 \\
Topioca Starch & 0.5 & 60 & 0.0514 & 0.0943536 & -0.0429536 \\
Topioca Starch & 0.5 & 0 & 0.6085 & 1.2333461 & -0.6248461 \\
Topioca Starch & 0.5 & 15 & 0.3484 & 0.7410020 & -0.3926020 \\
Topioca Starch & 0.5 & 30 & 0.1948 & 0.3870552 & -0.1922552 \\
Topioca Starch & 0.5 & 45 & 0.0796 & 0.1715057 & -0.0919057 \\
Topioca Starch & 0.5 & 60 & 0.0478 & 0.0943536 & -0.0465536 \\
Corn Starch & 0.5 & 0 & 0.5642 & 1.1313786 & -0.5671786 \\
Corn Starch & 0.5 & 15 & 0.2928 & 0.6390345 & -0.3462345 \\
Corn Starch & 0.5 & 30 & 0.1546 & 0.2850877 & -0.1304877 \\
Corn Starch & 0.5 & 45 & 0.0104 & 0.0695382 & -0.0591382 \\
Corn Starch & 0.5 & 60 & 0.0082 & -0.0076139 & 0.0158139 \\
Corn Starch & 0.5 & 0 & 0.5695 & 1.1313786 & -0.5618786 \\
Corn Starch & 0.5 & 15 & 0.2964 & 0.6390345 & -0.3426345 \\
Corn Starch & 0.5 & 30 & 0.1498 & 0.2850877 & -0.1352877 \\
Corn Starch & 0.5 & 45 & 0.0118 & 0.0695382 & -0.0577382 \\
Corn Starch & 0.5 & 60 & 0.0071 & -0.0076139 & 0.0147139 \\
Rice Starch & 0.5 & 0 & 0.7184 & 1.2620011 & -0.5436011 \\
Rice Starch & 0.5 & 15 & 0.4562 & 0.7696570 & -0.3134570 \\
Rice Starch & 0.5 & 30 & 0.2485 & 0.4157102 & -0.1672102 \\
Rice Starch & 0.5 & 45 & 0.0368 & 0.2001607 & -0.1633607 \\
Rice Starch & 0.5 & 60 & 0.0245 & 0.1230086 & -0.0985086 \\
Rice Starch & 0.5 & 0 & 0.7246 & 1.2620011 & -0.5374011 \\
Rice Starch & 0.5 & 15 & 0.4628 & 0.7696570 & -0.3068570 \\
Rice Starch & 0.5 & 30 & 0.2554 & 0.4157102 & -0.1603102 \\
Rice Starch & 0.5 & 45 & 0.0412 & 0.2001607 & -0.1589607 \\
Rice Starch & 0.5 & 60 & 0.0281 & 0.1230086 & -0.0949086 \\
Soluble Starch & 1.5 & 0 & 1.2486 & 1.0952586 & 0.1533414 \\
Soluble Starch & 1.5 & 15 & 0.6735 & 0.6029145 & 0.0705855 \\
Soluble Starch & 1.5 & 30 & 0.3387 & 0.2489677 & 0.0897323 \\
Soluble Starch & 1.5 & 45 & 0.0286 & 0.0334182 & -0.0048182 \\
Soluble Starch & 1.5 & 60 & 0.0198 & -0.0437339 & 0.0635339 \\
Soluble Starch & 1.5 & 0 & 1.2418 & 1.0952586 & 0.1465414 \\
Soluble Starch & 1.5 & 15 & 0.6658 & 0.6029145 & 0.0628855 \\
Soluble Starch & 1.5 & 30 & 0.3314 & 0.2489677 & 0.0824323 \\
Soluble Starch & 1.5 & 45 & 0.0259 & 0.0334182 & -0.0075182 \\
Soluble Starch & 1.5 & 60 & 0.0176 & -0.0437339 & 0.0613339 \\
Topioca Starch & 1.5 & 0 & 1.4328 & 1.2333461 & 0.1994539 \\
Topioca Starch & 1.5 & 15 & 0.8468 & 0.7410020 & 0.1057980 \\
Topioca Starch & 1.5 & 30 & 0.4986 & 0.3870552 & 0.1115448 \\
Topioca Starch & 1.5 & 45 & 0.2068 & 0.1715057 & 0.0352943 \\
Topioca Starch & 1.5 & 60 & 0.1428 & 0.0943536 & 0.0484464 \\
Topioca Starch & 1.5 & 0 & 1.4268 & 1.2333461 & 0.1934539 \\
Topioca Starch & 1.5 & 15 & 0.8385 & 0.7410020 & 0.0974980 \\
Topioca Starch & 1.5 & 30 & 0.4912 & 0.3870552 & 0.1041448 \\
Topioca Starch & 1.5 & 45 & 0.1985 & 0.1715057 & 0.0269943 \\
Topioca Starch & 1.5 & 60 & 0.1362 & 0.0943536 & 0.0418464 \\
Corn Starch & 1.5 & 0 & 1.3258 & 1.1313786 & 0.1944214 \\
Corn Starch & 1.5 & 15 & 0.7165 & 0.6390345 & 0.0774655 \\
Corn Starch & 1.5 & 30 & 0.3868 & 0.2850877 & 0.1017123 \\
Corn Starch & 1.5 & 45 & 0.0352 & 0.0695382 & -0.0343382 \\
Corn Starch & 1.5 & 60 & 0.0248 & -0.0076139 & 0.0324139 \\
Corn Starch & 1.5 & 0 & 1.3342 & 1.1313786 & 0.2028214 \\
Corn Starch & 1.5 & 15 & 0.7248 & 0.6390345 & 0.0857655 \\
Corn Starch & 1.5 & 30 & 0.3796 & 0.2850877 & 0.0945123 \\
Corn Starch & 1.5 & 45 & 0.0358 & 0.0695382 & -0.0337382 \\
Corn Starch & 1.5 & 60 & 0.0221 & -0.0076139 & 0.0297139 \\
Rice Starch & 1.5 & 0 & 1.4862 & 1.2620011 & 0.2241989 \\
Rice Starch & 1.5 & 15 & 0.9386 & 0.7696570 & 0.1689430 \\
Rice Starch & 1.5 & 30 & 0.5568 & 0.4157102 & 0.1410898 \\
Rice Starch & 1.5 & 45 & 0.1126 & 0.2001607 & -0.0875607 \\
Rice Starch & 1.5 & 60 & 0.0826 & 0.1230086 & -0.0404086 \\
Rice Starch & 1.5 & 0 & 1.4988 & 1.2620011 & 0.2367989 \\
Rice Starch & 1.5 & 15 & 0.9462 & 0.7696570 & 0.1765430 \\
Rice Starch & 1.5 & 30 & 0.5486 & 0.4157102 & 0.1328898 \\
Rice Starch & 1.5 & 45 & 0.1048 & 0.2001607 & -0.0953607 \\
Rice Starch & 1.5 & 60 & 0.0765 & 0.1230086 & -0.0465086 \\
Soluble Starch & 2.0 & 0 & 1.5624 & 1.0952586 & 0.4671414 \\
Soluble Starch & 2.0 & 15 & 0.8426 & 0.6029145 & 0.2396855 \\
Soluble Starch & 2.0 & 30 & 0.4528 & 0.2489677 & 0.2038323 \\
Soluble Starch & 2.0 & 45 & 0.0412 & 0.0334182 & 0.0077818 \\
Soluble Starch & 2.0 & 60 & 0.0284 & -0.0437339 & 0.0721339 \\
Soluble Starch & 2.0 & 0 & 1.5558 & 1.0952586 & 0.4605414 \\
Soluble Starch & 2.0 & 15 & 0.8354 & 0.6029145 & 0.2324855 \\
Soluble Starch & 2.0 & 30 & 0.4465 & 0.2489677 & 0.1975323 \\
Soluble Starch & 2.0 & 45 & 0.0376 & 0.0334182 & 0.0041818 \\
Soluble Starch & 2.0 & 60 & 0.0259 & -0.0437339 & 0.0696339 \\
Topioca Starch & 2.0 & 0 & 1.7385 & 1.2333461 & 0.5051539 \\
Topioca Starch & 2.0 & 15 & 1.0224 & 0.7410020 & 0.2813980 \\
Topioca Starch & 2.0 & 30 & 0.6254 & 0.3870552 & 0.2383448 \\
Topioca Starch & 2.0 & 45 & 0.2826 & 0.1715057 & 0.1110943 \\
Topioca Starch & 2.0 & 60 & 0.1885 & 0.0943536 & 0.0941464 \\
Topioca Starch & 2.0 & 0 & 1.7316 & 1.2333461 & 0.4982539 \\
Topioca Starch & 2.0 & 15 & 1.0158 & 0.7410020 & 0.2747980 \\
Topioca Starch & 2.0 & 30 & 0.6187 & 0.3870552 & 0.2316448 \\
Topioca Starch & 2.0 & 45 & 0.2758 & 0.1715057 & 0.1042943 \\
Topioca Starch & 2.0 & 60 & 0.1824 & 0.0943536 & 0.0880464 \\
Corn Starch & 2.0 & 0 & 1.6488 & 1.1313786 & 0.5174214 \\
Corn Starch & 2.0 & 15 & 0.8958 & 0.6390345 & 0.2567655 \\
Corn Starch & 2.0 & 30 & 0.5186 & 0.2850877 & 0.2335123 \\
Corn Starch & 2.0 & 45 & 0.0485 & 0.0695382 & -0.0210382 \\
Corn Starch & 2.0 & 60 & 0.0354 & -0.0076139 & 0.0430139 \\
Corn Starch & 2.0 & 0 & 1.6554 & 1.1313786 & 0.5240214 \\
Corn Starch & 2.0 & 15 & 0.8886 & 0.6390345 & 0.2495655 \\
Corn Starch & 2.0 & 30 & 0.5114 & 0.2850877 & 0.2263123 \\
Corn Starch & 2.0 & 45 & 0.0458 & 0.0695382 & -0.0237382 \\
Corn Starch & 2.0 & 60 & 0.0328 & -0.0076139 & 0.0404139 \\
Rice Starch & 2.0 & 0 & 1.8624 & 1.2620011 & 0.6003989 \\
Rice Starch & 2.0 & 15 & 1.1842 & 0.7696570 & 0.4145430 \\
Rice Starch & 2.0 & 30 & 0.7128 & 0.4157102 & 0.2970898 \\
Rice Starch & 2.0 & 45 & 0.1658 & 0.2001607 & -0.0343607 \\
Rice Starch & 2.0 & 60 & 0.1184 & 0.1230086 & -0.0046086 \\
Rice Starch & 2.0 & 0 & 1.8756 & 1.2620011 & 0.6135989 \\
Rice Starch & 2.0 & 15 & 1.1786 & 0.7696570 & 0.4089430 \\
Rice Starch & 2.0 & 30 & 0.7056 & 0.4157102 & 0.2898898 \\
Rice Starch & 2.0 & 45 & 0.1586 & 0.2001607 & -0.0415607 \\
Rice Starch & 2.0 & 60 & 0.1126 & 0.1230086 & -0.0104086 \\
\end{longtable}
}

\section{Estimation of kinetic parameters Vo, Km and
Vmax}\label{estimation-of-kinetic-parameters-vo-km-and-vmax}

\begin{Shaded}
\begin{Highlighting}[]
\CommentTok{\# 1. Calculate Initial Velocity (v0)}
\CommentTok{\# We calculate v0 using the change btw Time = 0 and Time = 15}
\NormalTok{v0\_data }\OtherTok{\textless{}{-}}\NormalTok{ UV2 }\SpecialCharTok{\%\textgreater{}\%}
  \CommentTok{\# Exclude the negative control }
  \FunctionTok{filter}\NormalTok{(Starch }\SpecialCharTok{!=} \StringTok{"Negative control"}\NormalTok{) }\SpecialCharTok{\%\textgreater{}\%}
  \CommentTok{\# Filter only for the initial linear phase timepoints}
  \FunctionTok{filter}\NormalTok{(Time }\SpecialCharTok{\%in\%} \FunctionTok{c}\NormalTok{(}\DecValTok{0}\NormalTok{, }\DecValTok{15}\NormalTok{)) }\SpecialCharTok{\%\textgreater{}\%}
  \CommentTok{\# Group to average the two replicates in your dataset}
  \FunctionTok{group\_by}\NormalTok{(Starch, sub\_conc, Time) }\SpecialCharTok{\%\textgreater{}\%}
  \FunctionTok{summarize}\NormalTok{(}\AttributeTok{mean\_ABS =} \FunctionTok{mean}\NormalTok{(ABS), }\AttributeTok{.groups =} \StringTok{"drop"}\NormalTok{) }\SpecialCharTok{\%\textgreater{}\%}
  \CommentTok{\# Reshape data to have Time 0 and Time 15 in separate columns}
  \FunctionTok{pivot\_wider}\NormalTok{(}\AttributeTok{names\_from =}\NormalTok{ Time, }\AttributeTok{values\_from =}\NormalTok{ mean\_ABS,}
              \AttributeTok{names\_prefix =} \StringTok{"T\_"}\NormalTok{) }\SpecialCharTok{\%\textgreater{}\%}
  \FunctionTok{mutate}\NormalTok{(}
    \CommentTok{\# Calculate initial velocity}
    \AttributeTok{v0 =}\NormalTok{ (T\_0 }\SpecialCharTok{{-}}\NormalTok{ T\_15) }\SpecialCharTok{/} \DecValTok{15}
\NormalTok{  )}

\FunctionTok{print}\NormalTok{(}\StringTok{"Calculated Initial Velocities (v0):"}\NormalTok{)}
\end{Highlighting}
\end{Shaded}

\begin{verbatim}
[1] "Calculated Initial Velocities (v0):"
\end{verbatim}

\begin{Shaded}
\begin{Highlighting}[]
\FunctionTok{kable}\NormalTok{(v0\_data, }\AttributeTok{format =} \StringTok{\textquotesingle{}latex\textquotesingle{}}\NormalTok{, }\AttributeTok{booktabs =} \ConstantTok{TRUE}\NormalTok{, }\AttributeTok{escape =}\NormalTok{ T)}
\end{Highlighting}
\end{Shaded}

\begin{tabular}{lrrrr}
\toprule
Starch & sub\_conc & T\_0 & T\_15 & v0\\
\midrule
Soluble Starch & 0.5 & 0.51525 & 0.28420 & 0.0154033\\
Soluble Starch & 1.0 & 0.89945 & 0.44975 & 0.0299800\\
Soluble Starch & 1.5 & 1.24520 & 0.66965 & 0.0383700\\
Soluble Starch & 2.0 & 1.55910 & 0.83900 & 0.0480067\\
Corn Starch & 0.5 & 0.56685 & 0.29460 & 0.0181500\\
\addlinespace
Corn Starch & 1.0 & 0.98785 & 0.50925 & 0.0319067\\
Corn Starch & 1.5 & 1.33000 & 0.72065 & 0.0406233\\
Corn Starch & 2.0 & 1.65210 & 0.89220 & 0.0506600\\
Rice Starch & 0.5 & 0.72150 & 0.45950 & 0.0174667\\
Rice Starch & 1.0 & 1.23150 & 0.67665 & 0.0369900\\
\addlinespace
Rice Starch & 1.5 & 1.49250 & 0.94240 & 0.0366733\\
Rice Starch & 2.0 & 1.86900 & 1.18140 & 0.0458400\\
Topioca Starch & 0.5 & 0.61065 & 0.35030 & 0.0173567\\
Topioca Starch & 1.0 & 1.05970 & 0.62015 & 0.0293033\\
Topioca Starch & 1.5 & 1.42980 & 0.84265 & 0.0391433\\
\addlinespace
Topioca Starch & 2.0 & 1.73505 & 1.01910 & 0.0477300\\
\bottomrule
\end{tabular}

\subsection{Fitting Michaeli-Menten Model for the four types of
starch}\label{fitting-michaeli-menten-model-for-the-four-types-of-starch}

\begin{Shaded}
\begin{Highlighting}[]
\NormalTok{mm\_results }\OtherTok{\textless{}{-}}\NormalTok{ v0\_data }\SpecialCharTok{\%\textgreater{}\%}
  \FunctionTok{group\_by}\NormalTok{(Starch) }\SpecialCharTok{\%\textgreater{}\%}
  \FunctionTok{nest}\NormalTok{() }\SpecialCharTok{\%\textgreater{}\%}
  \FunctionTok{mutate}\NormalTok{(}
    \AttributeTok{model =} \FunctionTok{map}\NormalTok{(data, }\SpecialCharTok{\textasciitilde{}} \FunctionTok{nls}\NormalTok{(}
\NormalTok{      v0 }\SpecialCharTok{\textasciitilde{}}\NormalTok{ (Vmax }\SpecialCharTok{*}\NormalTok{ sub\_conc) }\SpecialCharTok{/}\NormalTok{ (Km }\SpecialCharTok{+}\NormalTok{ sub\_conc),}
      \AttributeTok{data =}\NormalTok{ .x,}
      \CommentTok{\# Dynamic start values based on the data}
      \AttributeTok{start =} \FunctionTok{list}\NormalTok{(}\AttributeTok{Vmax =} \FunctionTok{max}\NormalTok{(.x}\SpecialCharTok{$}\NormalTok{v0), }\AttributeTok{Km =} \FunctionTok{median}\NormalTok{(.x}\SpecialCharTok{$}\NormalTok{sub\_conc))}
\NormalTok{    )),}
    \CommentTok{\# Extract the Vmax and Km estimates cleanly}
    \AttributeTok{params =} \FunctionTok{map}\NormalTok{(model, tidy)}
\NormalTok{  ) }\SpecialCharTok{\%\textgreater{}\%}
  \FunctionTok{unnest}\NormalTok{(params) }\SpecialCharTok{\%\textgreater{}\%}
  \FunctionTok{select}\NormalTok{(Starch, term, estimate, std.error, p.value)}

\FunctionTok{print}\NormalTok{(}\StringTok{"Estimated Vmax and Km parameters:"}\NormalTok{)}
\end{Highlighting}
\end{Shaded}

\begin{verbatim}
[1] "Estimated Vmax and Km parameters:"
\end{verbatim}

\begin{Shaded}
\begin{Highlighting}[]
\FunctionTok{kable}\NormalTok{(mm\_results)}
\end{Highlighting}
\end{Shaded}

{\def\LTcaptype{none} % do not increment counter
\begin{longtable}[]{@{}llrrr@{}}
\toprule\noalign{}
Starch & term & estimate & std.error & p.value \\
\midrule\noalign{}
\endhead
\bottomrule\noalign{}
\endlastfoot
Soluble Starch & Vmax & 0.1394114 & 0.0248543 & 0.0303446 \\
Soluble Starch & Km & 3.8346015 & 0.9569834 & 0.0570089 \\
Corn Starch & Vmax & 0.1237107 & 0.0148865 & 0.0141730 \\
Corn Starch & Km & 2.9409360 & 0.5344320 & 0.0314722 \\
Rice Starch & Vmax & 0.0772405 & 0.0263232 & 0.0991661 \\
Rice Starch & Km & 1.4089156 & 0.9497896 & 0.2762165 \\
Topioca Starch & Vmax & 0.1199012 & 0.0080030 & 0.0044255 \\
Topioca Starch & Km & 3.0542478 & 0.3043381 & 0.0097835 \\
\end{longtable}
}

\begin{Shaded}
\begin{Highlighting}[]
\DocumentationTok{\#\# Another way of fitting MM model}
\FunctionTok{suppressMessages}\NormalTok{(}
  \FunctionTok{suppressWarnings}\NormalTok{(}
    \FunctionTok{library}\NormalTok{(drc, }\AttributeTok{quietly =} \ConstantTok{TRUE}\NormalTok{)}
\NormalTok{  )}
\NormalTok{)}
\NormalTok{fit }\OtherTok{\textless{}{-}} \FunctionTok{drm}\NormalTok{(v0}\SpecialCharTok{\textasciitilde{}}\NormalTok{sub\_conc, }\AttributeTok{data =}\NormalTok{ v0\_data, }\AttributeTok{fct =} \FunctionTok{MM.2}\NormalTok{())}
\FunctionTok{summary}\NormalTok{(fit)}
\end{Highlighting}
\end{Shaded}

\begin{verbatim}

Model fitted: Michaelis-Menten (2 parms)

Parameter estimates:

              Estimate Std. Error t-value   p-value    
d:(Intercept) 0.107786   0.014422  7.4736 2.996e-06 ***
e:(Intercept) 2.530430   0.536467  4.7168 0.0003305 ***
---
Signif. codes:  0 '***' 0.001 '**' 0.01 '*' 0.05 '.' 0.1 ' ' 1

Residual standard error:

 0.002395029 (14 degrees of freedom)
\end{verbatim}

\begin{Shaded}
\begin{Highlighting}[]
\CommentTok{\# Fit separate MM curves for each Starch group}
\NormalTok{fit\_grouped }\OtherTok{\textless{}{-}} \FunctionTok{drm}\NormalTok{(v0 }\SpecialCharTok{\textasciitilde{}}\NormalTok{ sub\_conc, }\AttributeTok{curveid =}\NormalTok{ Starch, }
                   \AttributeTok{data =}\NormalTok{ v0\_data, }\AttributeTok{fct =} \FunctionTok{MM.2}\NormalTok{())}

\CommentTok{\# Summary of Vmax (d) and Km (e) for each starch type}
\FunctionTok{summary}\NormalTok{(fit\_grouped)}
\end{Highlighting}
\end{Shaded}

\begin{verbatim}

Model fitted: Michaelis-Menten (2 parms)

Parameter estimates:

                 Estimate Std. Error t-value   p-value    
d:Soluble Starch 0.138942   0.053124  2.6154 0.0308689 *  
d:Corn Starch    0.123652   0.036633  3.3754 0.0097078 ** 
d:Rice Starch    0.077240   0.013407  5.7611 0.0004236 ***
d:Topioca Starch 0.119825   0.038667  3.0989 0.0146907 *  
e:Soluble Starch 3.816484   2.045083  1.8662 0.0989871 .  
e:Corn Starch    2.938814   1.315300  2.2343 0.0559180 .  
e:Rice Starch    1.408915   0.483061  2.9166 0.0193907 *  
e:Topioca Starch 3.051361   1.470592  2.0749 0.0716789 .  
---
Signif. codes:  0 '***' 0.001 '**' 0.01 '*' 0.05 '.' 0.1 ' ' 1

Residual standard error:

 0.002427663 (8 degrees of freedom)
\end{verbatim}

\begin{Shaded}
\begin{Highlighting}[]
\CommentTok{\# Predict smooth curves from fit\_grouped model}
\NormalTok{pred\_grid }\OtherTok{\textless{}{-}} \FunctionTok{expand.grid}\NormalTok{(}
  \AttributeTok{sub\_conc =} \FunctionTok{seq}\NormalTok{(}\FloatTok{0.1}\NormalTok{, }\FloatTok{2.5}\NormalTok{, }\AttributeTok{length.out =} \DecValTok{100}\NormalTok{),}
  \AttributeTok{Starch =} \FunctionTok{unique}\NormalTok{(v0\_data}\SpecialCharTok{$}\NormalTok{Starch)}
\NormalTok{)}
\NormalTok{pred\_grid}\SpecialCharTok{$}\NormalTok{pred\_v0 }\OtherTok{\textless{}{-}} \FunctionTok{predict}\NormalTok{(fit\_grouped, }\AttributeTok{newdata =}\NormalTok{ pred\_grid)}
\FunctionTok{kable}\NormalTok{(pred\_grid, }\AttributeTok{type =} \StringTok{\textquotesingle{}latex\textquotesingle{}}\NormalTok{, }\AttributeTok{escape =}\NormalTok{ T, }\AttributeTok{booktabs=} \ConstantTok{TRUE}\NormalTok{)}
\end{Highlighting}
\end{Shaded}

{\def\LTcaptype{none} % do not increment counter
\begin{longtable}[]{@{}rlr@{}}
\toprule\noalign{}
sub\_conc & Starch & pred\_v0 \\
\midrule\noalign{}
\endhead
\bottomrule\noalign{}
\endlastfoot
0.1000000 & Soluble Starch & 0.0035476 \\
0.1242424 & Soluble Starch & 0.0043805 \\
0.1484848 & Soluble Starch & 0.0052033 \\
0.1727273 & Soluble Starch & 0.0060160 \\
0.1969697 & Soluble Starch & 0.0068189 \\
0.2212121 & Soluble Starch & 0.0076122 \\
0.2454545 & Soluble Starch & 0.0083960 \\
0.2696970 & Soluble Starch & 0.0091705 \\
0.2939394 & Soluble Starch & 0.0099358 \\
0.3181818 & Soluble Starch & 0.0106922 \\
0.3424242 & Soluble Starch & 0.0114398 \\
0.3666667 & Soluble Starch & 0.0121787 \\
0.3909091 & Soluble Starch & 0.0129091 \\
0.4151515 & Soluble Starch & 0.0136311 \\
0.4393939 & Soluble Starch & 0.0143449 \\
0.4636364 & Soluble Starch & 0.0150506 \\
0.4878788 & Soluble Starch & 0.0157484 \\
0.5121212 & Soluble Starch & 0.0164383 \\
0.5363636 & Soluble Starch & 0.0171206 \\
0.5606061 & Soluble Starch & 0.0177953 \\
0.5848485 & Soluble Starch & 0.0184625 \\
0.6090909 & Soluble Starch & 0.0191225 \\
0.6333333 & Soluble Starch & 0.0197753 \\
0.6575758 & Soluble Starch & 0.0204210 \\
0.6818182 & Soluble Starch & 0.0210597 \\
0.7060606 & Soluble Starch & 0.0216916 \\
0.7303030 & Soluble Starch & 0.0223167 \\
0.7545455 & Soluble Starch & 0.0229353 \\
0.7787879 & Soluble Starch & 0.0235473 \\
0.8030303 & Soluble Starch & 0.0241528 \\
0.8272727 & Soluble Starch & 0.0247521 \\
0.8515152 & Soluble Starch & 0.0253451 \\
0.8757576 & Soluble Starch & 0.0259320 \\
0.9000000 & Soluble Starch & 0.0265129 \\
0.9242424 & Soluble Starch & 0.0270878 \\
0.9484848 & Soluble Starch & 0.0276568 \\
0.9727273 & Soluble Starch & 0.0282202 \\
0.9969697 & Soluble Starch & 0.0287778 \\
1.0212121 & Soluble Starch & 0.0293298 \\
1.0454545 & Soluble Starch & 0.0298764 \\
1.0696970 & Soluble Starch & 0.0304175 \\
1.0939394 & Soluble Starch & 0.0309533 \\
1.1181818 & Soluble Starch & 0.0314838 \\
1.1424242 & Soluble Starch & 0.0320091 \\
1.1666667 & Soluble Starch & 0.0325293 \\
1.1909091 & Soluble Starch & 0.0330445 \\
1.2151515 & Soluble Starch & 0.0335547 \\
1.2393939 & Soluble Starch & 0.0340600 \\
1.2636364 & Soluble Starch & 0.0345605 \\
1.2878788 & Soluble Starch & 0.0350563 \\
1.3121212 & Soluble Starch & 0.0355473 \\
1.3363636 & Soluble Starch & 0.0360338 \\
1.3606061 & Soluble Starch & 0.0365156 \\
1.3848485 & Soluble Starch & 0.0369930 \\
1.4090909 & Soluble Starch & 0.0374660 \\
1.4333333 & Soluble Starch & 0.0379346 \\
1.4575758 & Soluble Starch & 0.0383989 \\
1.4818182 & Soluble Starch & 0.0388589 \\
1.5060606 & Soluble Starch & 0.0393147 \\
1.5303030 & Soluble Starch & 0.0397665 \\
1.5545455 & Soluble Starch & 0.0402141 \\
1.5787879 & Soluble Starch & 0.0406577 \\
1.6030303 & Soluble Starch & 0.0410973 \\
1.6272727 & Soluble Starch & 0.0415331 \\
1.6515152 & Soluble Starch & 0.0419649 \\
1.6757576 & Soluble Starch & 0.0423930 \\
1.7000000 & Soluble Starch & 0.0428173 \\
1.7242424 & Soluble Starch & 0.0432378 \\
1.7484848 & Soluble Starch & 0.0436547 \\
1.7727273 & Soluble Starch & 0.0440680 \\
1.7969697 & Soluble Starch & 0.0444778 \\
1.8212121 & Soluble Starch & 0.0448840 \\
1.8454545 & Soluble Starch & 0.0452867 \\
1.8696970 & Soluble Starch & 0.0456860 \\
1.8939394 & Soluble Starch & 0.0460819 \\
1.9181818 & Soluble Starch & 0.0464744 \\
1.9424242 & Soluble Starch & 0.0468637 \\
1.9666667 & Soluble Starch & 0.0472496 \\
1.9909091 & Soluble Starch & 0.0476324 \\
2.0151515 & Soluble Starch & 0.0480120 \\
2.0393939 & Soluble Starch & 0.0483884 \\
2.0636364 & Soluble Starch & 0.0487617 \\
2.0878788 & Soluble Starch & 0.0491320 \\
2.1121212 & Soluble Starch & 0.0494993 \\
2.1363636 & Soluble Starch & 0.0498635 \\
2.1606061 & Soluble Starch & 0.0502248 \\
2.1848485 & Soluble Starch & 0.0505832 \\
2.2090909 & Soluble Starch & 0.0509386 \\
2.2333333 & Soluble Starch & 0.0512913 \\
2.2575758 & Soluble Starch & 0.0516411 \\
2.2818182 & Soluble Starch & 0.0519882 \\
2.3060606 & Soluble Starch & 0.0523325 \\
2.3303030 & Soluble Starch & 0.0526740 \\
2.3545455 & Soluble Starch & 0.0530129 \\
2.3787879 & Soluble Starch & 0.0533492 \\
2.4030303 & Soluble Starch & 0.0536828 \\
2.4272727 & Soluble Starch & 0.0540138 \\
2.4515152 & Soluble Starch & 0.0543423 \\
2.4757576 & Soluble Starch & 0.0546682 \\
2.5000000 & Soluble Starch & 0.0549917 \\
0.1000000 & Corn Starch & 0.0040691 \\
0.1242424 & Corn Starch & 0.0050155 \\
0.1484848 & Corn Starch & 0.0059471 \\
0.1727273 & Corn Starch & 0.0068641 \\
0.1969697 & Corn Starch & 0.0077670 \\
0.2212121 & Corn Starch & 0.0086560 \\
0.2454545 & Corn Starch & 0.0095315 \\
0.2696970 & Corn Starch & 0.0103938 \\
0.2939394 & Corn Starch & 0.0112431 \\
0.3181818 & Corn Starch & 0.0120798 \\
0.3424242 & Corn Starch & 0.0129041 \\
0.3666667 & Corn Starch & 0.0137163 \\
0.3909091 & Corn Starch & 0.0145167 \\
0.4151515 & Corn Starch & 0.0153055 \\
0.4393939 & Corn Starch & 0.0160830 \\
0.4636364 & Corn Starch & 0.0168495 \\
0.4878788 & Corn Starch & 0.0176051 \\
0.5121212 & Corn Starch & 0.0183500 \\
0.5363636 & Corn Starch & 0.0190846 \\
0.5606061 & Corn Starch & 0.0198090 \\
0.5848485 & Corn Starch & 0.0205234 \\
0.6090909 & Corn Starch & 0.0212281 \\
0.6333333 & Corn Starch & 0.0219232 \\
0.6575758 & Corn Starch & 0.0226089 \\
0.6818182 & Corn Starch & 0.0232855 \\
0.7060606 & Corn Starch & 0.0239530 \\
0.7303030 & Corn Starch & 0.0246117 \\
0.7545455 & Corn Starch & 0.0252618 \\
0.7787879 & Corn Starch & 0.0259034 \\
0.8030303 & Corn Starch & 0.0265367 \\
0.8272727 & Corn Starch & 0.0271618 \\
0.8515152 & Corn Starch & 0.0277790 \\
0.8757576 & Corn Starch & 0.0283883 \\
0.9000000 & Corn Starch & 0.0289899 \\
0.9242424 & Corn Starch & 0.0295839 \\
0.9484848 & Corn Starch & 0.0301705 \\
0.9727273 & Corn Starch & 0.0307499 \\
0.9969697 & Corn Starch & 0.0313221 \\
1.0212121 & Corn Starch & 0.0318874 \\
1.0454545 & Corn Starch & 0.0324457 \\
1.0696970 & Corn Starch & 0.0329973 \\
1.0939394 & Corn Starch & 0.0335422 \\
1.1181818 & Corn Starch & 0.0340807 \\
1.1424242 & Corn Starch & 0.0346127 \\
1.1666667 & Corn Starch & 0.0351385 \\
1.1909091 & Corn Starch & 0.0356581 \\
1.2151515 & Corn Starch & 0.0361716 \\
1.2393939 & Corn Starch & 0.0366792 \\
1.2636364 & Corn Starch & 0.0371809 \\
1.2878788 & Corn Starch & 0.0376769 \\
1.3121212 & Corn Starch & 0.0381672 \\
1.3363636 & Corn Starch & 0.0386519 \\
1.3606061 & Corn Starch & 0.0391312 \\
1.3848485 & Corn Starch & 0.0396051 \\
1.4090909 & Corn Starch & 0.0400737 \\
1.4333333 & Corn Starch & 0.0405371 \\
1.4575758 & Corn Starch & 0.0409954 \\
1.4818182 & Corn Starch & 0.0414487 \\
1.5060606 & Corn Starch & 0.0418971 \\
1.5303030 & Corn Starch & 0.0423405 \\
1.5545455 & Corn Starch & 0.0427792 \\
1.5787879 & Corn Starch & 0.0432132 \\
1.6030303 & Corn Starch & 0.0436425 \\
1.6272727 & Corn Starch & 0.0440673 \\
1.6515152 & Corn Starch & 0.0444876 \\
1.6757576 & Corn Starch & 0.0449035 \\
1.7000000 & Corn Starch & 0.0453151 \\
1.7242424 & Corn Starch & 0.0457223 \\
1.7484848 & Corn Starch & 0.0461254 \\
1.7727273 & Corn Starch & 0.0465243 \\
1.7969697 & Corn Starch & 0.0469191 \\
1.8212121 & Corn Starch & 0.0473099 \\
1.8454545 & Corn Starch & 0.0476967 \\
1.8696970 & Corn Starch & 0.0480796 \\
1.8939394 & Corn Starch & 0.0484587 \\
1.9181818 & Corn Starch & 0.0488340 \\
1.9424242 & Corn Starch & 0.0492056 \\
1.9666667 & Corn Starch & 0.0495735 \\
1.9909091 & Corn Starch & 0.0499378 \\
2.0151515 & Corn Starch & 0.0502985 \\
2.0393939 & Corn Starch & 0.0506557 \\
2.0636364 & Corn Starch & 0.0510095 \\
2.0878788 & Corn Starch & 0.0513598 \\
2.1121212 & Corn Starch & 0.0517068 \\
2.1363636 & Corn Starch & 0.0520504 \\
2.1606061 & Corn Starch & 0.0523908 \\
2.1848485 & Corn Starch & 0.0527280 \\
2.2090909 & Corn Starch & 0.0530620 \\
2.2333333 & Corn Starch & 0.0533929 \\
2.2575758 & Corn Starch & 0.0537206 \\
2.2818182 & Corn Starch & 0.0540454 \\
2.3060606 & Corn Starch & 0.0543671 \\
2.3303030 & Corn Starch & 0.0546859 \\
2.3545455 & Corn Starch & 0.0550017 \\
2.3787879 & Corn Starch & 0.0553147 \\
2.4030303 & Corn Starch & 0.0556248 \\
2.4272727 & Corn Starch & 0.0559321 \\
2.4515152 & Corn Starch & 0.0562367 \\
2.4757576 & Corn Starch & 0.0565385 \\
2.5000000 & Corn Starch & 0.0568377 \\
0.1000000 & Rice Starch & 0.0051189 \\
0.1242424 & Rice Starch & 0.0062593 \\
0.1484848 & Rice Starch & 0.0073642 \\
0.1727273 & Rice Starch & 0.0084352 \\
0.1969697 & Rice Starch & 0.0094739 \\
0.2212121 & Rice Starch & 0.0104816 \\
0.2454545 & Rice Starch & 0.0114599 \\
0.2696970 & Rice Starch & 0.0124099 \\
0.2939394 & Rice Starch & 0.0133328 \\
0.3181818 & Rice Starch & 0.0142299 \\
0.3424242 & Rice Starch & 0.0151021 \\
0.3666667 & Rice Starch & 0.0159504 \\
0.3909091 & Rice Starch & 0.0167760 \\
0.4151515 & Rice Starch & 0.0175796 \\
0.4393939 & Rice Starch & 0.0183621 \\
0.4636364 & Rice Starch & 0.0191243 \\
0.4878788 & Rice Starch & 0.0198671 \\
0.5121212 & Rice Starch & 0.0205911 \\
0.5363636 & Rice Starch & 0.0212970 \\
0.5606061 & Rice Starch & 0.0219856 \\
0.5848485 & Rice Starch & 0.0226575 \\
0.6090909 & Rice Starch & 0.0233132 \\
0.6333333 & Rice Starch & 0.0239533 \\
0.6575758 & Rice Starch & 0.0245784 \\
0.6818182 & Rice Starch & 0.0251891 \\
0.7060606 & Rice Starch & 0.0257857 \\
0.7303030 & Rice Starch & 0.0263688 \\
0.7545455 & Rice Starch & 0.0269388 \\
0.7787879 & Rice Starch & 0.0274962 \\
0.8030303 & Rice Starch & 0.0280414 \\
0.8272727 & Rice Starch & 0.0285747 \\
0.8515152 & Rice Starch & 0.0290967 \\
0.8757576 & Rice Starch & 0.0296075 \\
0.9000000 & Rice Starch & 0.0301076 \\
0.9242424 & Rice Starch & 0.0305973 \\
0.9484848 & Rice Starch & 0.0310770 \\
0.9727273 & Rice Starch & 0.0315469 \\
0.9969697 & Rice Starch & 0.0320073 \\
1.0212121 & Rice Starch & 0.0324585 \\
1.0454545 & Rice Starch & 0.0329009 \\
1.0696970 & Rice Starch & 0.0333345 \\
1.0939394 & Rice Starch & 0.0337598 \\
1.1181818 & Rice Starch & 0.0341769 \\
1.1424242 & Rice Starch & 0.0345861 \\
1.1666667 & Rice Starch & 0.0349875 \\
1.1909091 & Rice Starch & 0.0353815 \\
1.2151515 & Rice Starch & 0.0357682 \\
1.2393939 & Rice Starch & 0.0361479 \\
1.2636364 & Rice Starch & 0.0365206 \\
1.2878788 & Rice Starch & 0.0368866 \\
1.3121212 & Rice Starch & 0.0372462 \\
1.3363636 & Rice Starch & 0.0375993 \\
1.3606061 & Rice Starch & 0.0379463 \\
1.3848485 & Rice Starch & 0.0382873 \\
1.4090909 & Rice Starch & 0.0386224 \\
1.4333333 & Rice Starch & 0.0389518 \\
1.4575758 & Rice Starch & 0.0392756 \\
1.4818182 & Rice Starch & 0.0395940 \\
1.5060606 & Rice Starch & 0.0399070 \\
1.5303030 & Rice Starch & 0.0402150 \\
1.5545455 & Rice Starch & 0.0405178 \\
1.5787879 & Rice Starch & 0.0408158 \\
1.6030303 & Rice Starch & 0.0411090 \\
1.6272727 & Rice Starch & 0.0413975 \\
1.6515152 & Rice Starch & 0.0416814 \\
1.6757576 & Rice Starch & 0.0419608 \\
1.7000000 & Rice Starch & 0.0422359 \\
1.7242424 & Rice Starch & 0.0425068 \\
1.7484848 & Rice Starch & 0.0427734 \\
1.7727273 & Rice Starch & 0.0430361 \\
1.7969697 & Rice Starch & 0.0432947 \\
1.8212121 & Rice Starch & 0.0435495 \\
1.8454545 & Rice Starch & 0.0438004 \\
1.8696970 & Rice Starch & 0.0440477 \\
1.8939394 & Rice Starch & 0.0442913 \\
1.9181818 & Rice Starch & 0.0445314 \\
1.9424242 & Rice Starch & 0.0447680 \\
1.9666667 & Rice Starch & 0.0450012 \\
1.9909091 & Rice Starch & 0.0452311 \\
2.0151515 & Rice Starch & 0.0454577 \\
2.0393939 & Rice Starch & 0.0456811 \\
2.0636364 & Rice Starch & 0.0459015 \\
2.0878788 & Rice Starch & 0.0461187 \\
2.1121212 & Rice Starch & 0.0463330 \\
2.1363636 & Rice Starch & 0.0465443 \\
2.1606061 & Rice Starch & 0.0467528 \\
2.1848485 & Rice Starch & 0.0469585 \\
2.2090909 & Rice Starch & 0.0471614 \\
2.2333333 & Rice Starch & 0.0473616 \\
2.2575758 & Rice Starch & 0.0475591 \\
2.2818182 & Rice Starch & 0.0477541 \\
2.3060606 & Rice Starch & 0.0479465 \\
2.3303030 & Rice Starch & 0.0481364 \\
2.3545455 & Rice Starch & 0.0483239 \\
2.3787879 & Rice Starch & 0.0485089 \\
2.4030303 & Rice Starch & 0.0486917 \\
2.4272727 & Rice Starch & 0.0488721 \\
2.4515152 & Rice Starch & 0.0490502 \\
2.4757576 & Rice Starch & 0.0492261 \\
2.5000000 & Rice Starch & 0.0493999 \\
0.1000000 & Topioca Starch & 0.0038023 \\
0.1242424 & Topioca Starch & 0.0046880 \\
0.1484848 & Topioca Starch & 0.0055603 \\
0.1727273 & Topioca Starch & 0.0064195 \\
0.1969697 & Topioca Starch & 0.0072658 \\
0.2212121 & Topioca Starch & 0.0080996 \\
0.2454545 & Topioca Starch & 0.0089212 \\
0.2696970 & Topioca Starch & 0.0097307 \\
0.2939394 & Topioca Starch & 0.0105286 \\
0.3181818 & Topioca Starch & 0.0113149 \\
0.3424242 & Topioca Starch & 0.0120900 \\
0.3666667 & Topioca Starch & 0.0128541 \\
0.3909091 & Topioca Starch & 0.0136075 \\
0.4151515 & Topioca Starch & 0.0143503 \\
0.4393939 & Topioca Starch & 0.0150828 \\
0.4636364 & Topioca Starch & 0.0158052 \\
0.4878788 & Topioca Starch & 0.0165177 \\
0.5121212 & Topioca Starch & 0.0172205 \\
0.5363636 & Topioca Starch & 0.0179138 \\
0.5606061 & Topioca Starch & 0.0185978 \\
0.5848485 & Topioca Starch & 0.0192726 \\
0.6090909 & Topioca Starch & 0.0199386 \\
0.6333333 & Topioca Starch & 0.0205957 \\
0.6575758 & Topioca Starch & 0.0212443 \\
0.6818182 & Topioca Starch & 0.0218845 \\
0.7060606 & Topioca Starch & 0.0225164 \\
0.7303030 & Topioca Starch & 0.0231402 \\
0.7545455 & Topioca Starch & 0.0237560 \\
0.7787879 & Topioca Starch & 0.0243641 \\
0.8030303 & Topioca Starch & 0.0249645 \\
0.8272727 & Topioca Starch & 0.0255574 \\
0.8515152 & Topioca Starch & 0.0261429 \\
0.8757576 & Topioca Starch & 0.0267212 \\
0.9000000 & Topioca Starch & 0.0272924 \\
0.9242424 & Topioca Starch & 0.0278567 \\
0.9484848 & Topioca Starch & 0.0284141 \\
0.9727273 & Topioca Starch & 0.0289648 \\
0.9969697 & Topioca Starch & 0.0295089 \\
1.0212121 & Topioca Starch & 0.0300465 \\
1.0454545 & Topioca Starch & 0.0305777 \\
1.0696970 & Topioca Starch & 0.0311027 \\
1.0939394 & Topioca Starch & 0.0316216 \\
1.1181818 & Topioca Starch & 0.0321344 \\
1.1424242 & Topioca Starch & 0.0326413 \\
1.1666667 & Topioca Starch & 0.0331424 \\
1.1909091 & Topioca Starch & 0.0336377 \\
1.2151515 & Topioca Starch & 0.0341275 \\
1.2393939 & Topioca Starch & 0.0346116 \\
1.2636364 & Topioca Starch & 0.0350904 \\
1.2878788 & Topioca Starch & 0.0355638 \\
1.3121212 & Topioca Starch & 0.0360319 \\
1.3363636 & Topioca Starch & 0.0364949 \\
1.3606061 & Topioca Starch & 0.0369527 \\
1.3848485 & Topioca Starch & 0.0374056 \\
1.4090909 & Topioca Starch & 0.0378536 \\
1.4333333 & Topioca Starch & 0.0382967 \\
1.4575758 & Topioca Starch & 0.0387350 \\
1.4818182 & Topioca Starch & 0.0391687 \\
1.5060606 & Topioca Starch & 0.0395977 \\
1.5303030 & Topioca Starch & 0.0400222 \\
1.5545455 & Topioca Starch & 0.0404422 \\
1.5787879 & Topioca Starch & 0.0408578 \\
1.6030303 & Topioca Starch & 0.0412691 \\
1.6272727 & Topioca Starch & 0.0416762 \\
1.6515152 & Topioca Starch & 0.0420790 \\
1.6757576 & Topioca Starch & 0.0424777 \\
1.7000000 & Topioca Starch & 0.0428724 \\
1.7242424 & Topioca Starch & 0.0432630 \\
1.7484848 & Topioca Starch & 0.0436497 \\
1.7727273 & Topioca Starch & 0.0440325 \\
1.7969697 & Topioca Starch & 0.0444115 \\
1.8212121 & Topioca Starch & 0.0447867 \\
1.8454545 & Topioca Starch & 0.0451582 \\
1.8696970 & Topioca Starch & 0.0455260 \\
1.8939394 & Topioca Starch & 0.0458902 \\
1.9181818 & Topioca Starch & 0.0462509 \\
1.9424242 & Topioca Starch & 0.0466080 \\
1.9666667 & Topioca Starch & 0.0469618 \\
1.9909091 & Topioca Starch & 0.0473121 \\
2.0151515 & Topioca Starch & 0.0476590 \\
2.0393939 & Topioca Starch & 0.0480027 \\
2.0636364 & Topioca Starch & 0.0483431 \\
2.0878788 & Topioca Starch & 0.0486803 \\
2.1121212 & Topioca Starch & 0.0490143 \\
2.1363636 & Topioca Starch & 0.0493452 \\
2.1606061 & Topioca Starch & 0.0496730 \\
2.1848485 & Topioca Starch & 0.0499978 \\
2.2090909 & Topioca Starch & 0.0503196 \\
2.2333333 & Topioca Starch & 0.0506384 \\
2.2575758 & Topioca Starch & 0.0509544 \\
2.2818182 & Topioca Starch & 0.0512674 \\
2.3060606 & Topioca Starch & 0.0515776 \\
2.3303030 & Topioca Starch & 0.0518851 \\
2.3545455 & Topioca Starch & 0.0521897 \\
2.3787879 & Topioca Starch & 0.0524917 \\
2.4030303 & Topioca Starch & 0.0527910 \\
2.4272727 & Topioca Starch & 0.0530876 \\
2.4515152 & Topioca Starch & 0.0533816 \\
2.4757576 & Topioca Starch & 0.0536730 \\
2.5000000 & Topioca Starch & 0.0539619 \\
\end{longtable}
}

\subsection{Plot the Substrate Saturation
Curves}\label{plot-the-substrate-saturation-curves}

\begin{Shaded}
\begin{Highlighting}[]
\FunctionTok{ggplot}\NormalTok{(v0\_data, }\FunctionTok{aes}\NormalTok{(}\AttributeTok{x =}\NormalTok{ sub\_conc, }\AttributeTok{y =}\NormalTok{ v0, }\AttributeTok{color =}\NormalTok{ Starch)) }\SpecialCharTok{+}
  \FunctionTok{geom\_point}\NormalTok{(}\AttributeTok{alpha =} \FloatTok{0.5}\NormalTok{, }\AttributeTok{size =} \DecValTok{2}\NormalTok{) }\SpecialCharTok{+}
  \CommentTok{\# Fit the non{-}linear Michaelis{-}Menten curve onto the plot}
  \FunctionTok{stat\_smooth}\NormalTok{(}
    \AttributeTok{method =} \StringTok{"nls"}\NormalTok{, }\CommentTok{\# non{-}linear least squares}
    \AttributeTok{formula =}\NormalTok{ y }\SpecialCharTok{\textasciitilde{}}\NormalTok{ (Vmax }\SpecialCharTok{*}\NormalTok{ x) }\SpecialCharTok{/}\NormalTok{ (Km }\SpecialCharTok{+}\NormalTok{ x),}
    \AttributeTok{method.args =} \FunctionTok{list}\NormalTok{(}\AttributeTok{start =} \FunctionTok{list}\NormalTok{(}\AttributeTok{Vmax =} \FloatTok{0.05}\NormalTok{, }\AttributeTok{Km =} \DecValTok{1}\NormalTok{)),}
    \AttributeTok{se =} \ConstantTok{FALSE}\NormalTok{,}
    \AttributeTok{linewidth =} \DecValTok{1}
\NormalTok{  ) }\SpecialCharTok{+}
  \FunctionTok{labs}\NormalTok{(}
    \AttributeTok{x =} \StringTok{"Substrate Concentration (\%)"}\NormalTok{,}
    \AttributeTok{y =} \FunctionTok{expression}\NormalTok{(}\StringTok{"Initial Velocity ("} \SpecialCharTok{*}\NormalTok{ Delta }\SpecialCharTok{*} \StringTok{" Abs / min)"}\NormalTok{),}
    \AttributeTok{title =} \StringTok{"Michaelis{-}Menten Kinetics of Alpha{-}Amylase"}\NormalTok{,}
    \AttributeTok{subtitle =} \StringTok{"Substrate saturation curves for various starches"}\NormalTok{,}
    \AttributeTok{color =} \ConstantTok{NULL}
\NormalTok{  ) }\SpecialCharTok{+}
  \FunctionTok{theme\_minimal}\NormalTok{() }\SpecialCharTok{+}
  \FunctionTok{theme}\NormalTok{(}
    \AttributeTok{text =} \FunctionTok{element\_text}\NormalTok{(}
      \AttributeTok{family =}\NormalTok{ mfont,}
      \AttributeTok{color =} \StringTok{"black"}\NormalTok{,}
      \AttributeTok{size =} \DecValTok{12}
\NormalTok{      ),}
    \AttributeTok{axis.line =} \FunctionTok{element\_line}\NormalTok{(}\AttributeTok{color =} \StringTok{"black"}\NormalTok{),}
    \AttributeTok{axis.text =} \FunctionTok{element\_text}\NormalTok{(}\AttributeTok{color =} \StringTok{"black"}\NormalTok{),}
    \AttributeTok{legend.position =} \StringTok{"bottom"}\NormalTok{,}
    \AttributeTok{legend.margin =} \FunctionTok{margin}\NormalTok{(}\SpecialCharTok{{-}}\DecValTok{10}\NormalTok{, }\DecValTok{0}\NormalTok{, }\DecValTok{0}\NormalTok{, }\DecValTok{0}\NormalTok{),}
    \AttributeTok{plot.margin =} \FunctionTok{margin}\NormalTok{(}\DecValTok{5}\NormalTok{, }\DecValTok{5}\NormalTok{, }\DecValTok{5}\NormalTok{, }\DecValTok{5}\NormalTok{)}
\NormalTok{)}
\end{Highlighting}
\end{Shaded}

\begin{center}\includegraphics[width=1\linewidth]{../Figs/Plot the Substrate Saturation Curves-1} \end{center}

\begin{Shaded}
\begin{Highlighting}[]
\DocumentationTok{\#\#\#\#\#\#\#\#\#\#\#\#\#\#\# }



\CommentTok{\# Km estimates from summary(fit\_grouped) to overlay vertical lines}
\NormalTok{km\_df }\OtherTok{\textless{}{-}} \FunctionTok{data.frame}\NormalTok{(}
  \AttributeTok{Starch =} \FunctionTok{c}\NormalTok{(}\StringTok{"Soluble Starch"}\NormalTok{, }\StringTok{"Corn Starch"}\NormalTok{, }
             \StringTok{"Rice Starch"}\NormalTok{, }\StringTok{"Topioca Starch"}\NormalTok{),}
  \AttributeTok{Km =} \FunctionTok{c}\NormalTok{(}\FloatTok{3.816}\NormalTok{, }\FloatTok{2.939}\NormalTok{, }\FloatTok{1.409}\NormalTok{, }\FloatTok{3.051}\NormalTok{)}
\NormalTok{)}

\FunctionTok{kable}\NormalTok{(km\_df)}
\end{Highlighting}
\end{Shaded}

{\def\LTcaptype{none} % do not increment counter
\begin{longtable}[]{@{}lr@{}}
\toprule\noalign{}
Starch & Km \\
\midrule\noalign{}
\endhead
\bottomrule\noalign{}
\endlastfoot
Soluble Starch & 3.816 \\
Corn Starch & 2.939 \\
Rice Starch & 1.409 \\
Topioca Starch & 3.051 \\
\end{longtable}
}

\begin{Shaded}
\begin{Highlighting}[]
\FunctionTok{ggplot}\NormalTok{() }\SpecialCharTok{+}
  \CommentTok{\# Raw data points}
  \FunctionTok{geom\_point}\NormalTok{(}\AttributeTok{data =}\NormalTok{ v0\_data, }\FunctionTok{aes}\NormalTok{(}\AttributeTok{x =}\NormalTok{ sub\_conc, }\AttributeTok{y =}\NormalTok{ v0, }
                                 \AttributeTok{color =}\NormalTok{ Starch), }\AttributeTok{size =} \FloatTok{2.5}\NormalTok{) }\SpecialCharTok{+}
  \CommentTok{\# Fitted Michaelis{-}Menten lines}
  \FunctionTok{geom\_line}\NormalTok{(}\AttributeTok{data =}\NormalTok{ pred\_grid, }\FunctionTok{aes}\NormalTok{(}\AttributeTok{x =}\NormalTok{ sub\_conc, }\AttributeTok{y =}\NormalTok{ pred\_v0, }
                                  \AttributeTok{color =}\NormalTok{ Starch), }
            \AttributeTok{linewidth =} \DecValTok{1}\NormalTok{) }\SpecialCharTok{+}
  \CommentTok{\# Vertical dashed lines indicating Km}
  \FunctionTok{geom\_vline}\NormalTok{(}\AttributeTok{data =}\NormalTok{ km\_df, }\FunctionTok{aes}\NormalTok{(}\AttributeTok{xintercept =}\NormalTok{ Km, }\AttributeTok{color =}\NormalTok{ Starch), }
             \AttributeTok{linetype =} \StringTok{"dashed"}\NormalTok{, }\AttributeTok{alpha =} \FloatTok{0.7}\NormalTok{) }\SpecialCharTok{+}
  \FunctionTok{facet\_wrap}\NormalTok{(}\SpecialCharTok{\textasciitilde{}}\NormalTok{ Starch) }\SpecialCharTok{+}
  \FunctionTok{labs}\NormalTok{(}
    \AttributeTok{x =} \StringTok{"Substrate Concentration (\%)"}\NormalTok{,}
    \AttributeTok{y =} \FunctionTok{expression}\NormalTok{(}\StringTok{"Initial Velocity "} \SpecialCharTok{*}\NormalTok{ v[}\DecValTok{0}\NormalTok{] }\SpecialCharTok{*} \StringTok{" ("} \SpecialCharTok{*}\NormalTok{ Delta }\SpecialCharTok{*} \StringTok{"Abs/min)"}\NormalTok{),}
    \AttributeTok{title =} \StringTok{"Michaelis{-}Menten Fit by Starch Type"}\NormalTok{,}
    \AttributeTok{subtitle =} \StringTok{"Dashed lines indicate estimated Km (half{-}saturation) values"}
\NormalTok{  ) }\SpecialCharTok{+}
  \FunctionTok{theme\_bw}\NormalTok{() }\SpecialCharTok{+}
  \FunctionTok{theme}\NormalTok{(}
    \AttributeTok{text =} \FunctionTok{element\_text}\NormalTok{(}
      \AttributeTok{family =}\NormalTok{ mfont,}
      \AttributeTok{color =} \StringTok{"black"}\NormalTok{,}
      \AttributeTok{size =} \DecValTok{12}
\NormalTok{    ),}
    \AttributeTok{axis.line =} \FunctionTok{element\_line}\NormalTok{(}\AttributeTok{color =} \StringTok{"black"}\NormalTok{),}
    \AttributeTok{axis.text =} \FunctionTok{element\_text}\NormalTok{(}\AttributeTok{color =} \StringTok{"black"}\NormalTok{),}
    \AttributeTok{legend.margin =} \FunctionTok{margin}\NormalTok{(}\SpecialCharTok{{-}}\DecValTok{10}\NormalTok{, }\DecValTok{0}\NormalTok{, }\DecValTok{0}\NormalTok{, }\DecValTok{0}\NormalTok{),}
    \AttributeTok{plot.margin =} \FunctionTok{margin}\NormalTok{(}\DecValTok{5}\NormalTok{, }\DecValTok{5}\NormalTok{, }\DecValTok{5}\NormalTok{, }\DecValTok{5}\NormalTok{),}
    \AttributeTok{legend.position =} \StringTok{"none"}
\NormalTok{    )}
\end{Highlighting}
\end{Shaded}

\begin{center}\includegraphics[width=1\linewidth]{../Figs/Plot the Substrate Saturation Curves-2} \end{center}

\begin{Shaded}
\begin{Highlighting}[]
\CommentTok{\# Tidy extracted parameter estimates from fit\_grouped}
\NormalTok{param\_df }\OtherTok{\textless{}{-}} \FunctionTok{data.frame}\NormalTok{(}
  \AttributeTok{Starch =} \FunctionTok{rep}\NormalTok{(}\FunctionTok{c}\NormalTok{(}\StringTok{"Soluble Starch"}\NormalTok{, }\StringTok{"Corn Starch"}\NormalTok{, }
                 \StringTok{"Rice Starch"}\NormalTok{, }\StringTok{"Topioca Starch"}\NormalTok{), }\DecValTok{2}\NormalTok{),}
  \AttributeTok{Parameter =} \FunctionTok{c}\NormalTok{(}\FunctionTok{rep}\NormalTok{(}\StringTok{"Vmax (d)"}\NormalTok{, }\DecValTok{4}\NormalTok{), }\FunctionTok{rep}\NormalTok{(}\StringTok{"Km (e)"}\NormalTok{, }\DecValTok{4}\NormalTok{)),}
  \AttributeTok{Estimate =} \FunctionTok{c}\NormalTok{(}\FloatTok{0.1389}\NormalTok{, }\FloatTok{0.1237}\NormalTok{, }\FloatTok{0.0772}\NormalTok{, }\FloatTok{0.1198}\NormalTok{, }\FloatTok{3.816}\NormalTok{, }\FloatTok{2.939}\NormalTok{, }\FloatTok{1.409}\NormalTok{, }\FloatTok{3.051}\NormalTok{),}
  \AttributeTok{SE =} \FunctionTok{c}\NormalTok{(}\FloatTok{0.0531}\NormalTok{, }\FloatTok{0.0366}\NormalTok{, }\FloatTok{0.0134}\NormalTok{, }\FloatTok{0.0387}\NormalTok{, }\FloatTok{2.045}\NormalTok{, }\FloatTok{1.315}\NormalTok{, }\FloatTok{0.483}\NormalTok{, }\FloatTok{1.471}\NormalTok{)}
\NormalTok{)}

\FunctionTok{kable}\NormalTok{(param\_df)}
\end{Highlighting}
\end{Shaded}

{\def\LTcaptype{none} % do not increment counter
\begin{longtable}[]{@{}llrr@{}}
\toprule\noalign{}
Starch & Parameter & Estimate & SE \\
\midrule\noalign{}
\endhead
\bottomrule\noalign{}
\endlastfoot
Soluble Starch & Vmax (d) & 0.1389 & 0.0531 \\
Corn Starch & Vmax (d) & 0.1237 & 0.0366 \\
Rice Starch & Vmax (d) & 0.0772 & 0.0134 \\
Topioca Starch & Vmax (d) & 0.1198 & 0.0387 \\
Soluble Starch & Km (e) & 3.8160 & 2.0450 \\
Corn Starch & Km (e) & 2.9390 & 1.3150 \\
Rice Starch & Km (e) & 1.4090 & 0.4830 \\
Topioca Starch & Km (e) & 3.0510 & 1.4710 \\
\end{longtable}
}

\begin{Shaded}
\begin{Highlighting}[]
\FunctionTok{ggplot}\NormalTok{(param\_df, }\FunctionTok{aes}\NormalTok{(}\AttributeTok{x =}\NormalTok{ Starch, }\AttributeTok{y =}\NormalTok{ Estimate, }\AttributeTok{fill =}\NormalTok{ Starch)) }\SpecialCharTok{+}
  \FunctionTok{geom\_col}\NormalTok{(}\AttributeTok{width =} \FloatTok{0.6}\NormalTok{, }\AttributeTok{alpha =} \FloatTok{0.85}\NormalTok{, }\AttributeTok{color =} \StringTok{"black"}\NormalTok{) }\SpecialCharTok{+}
  \FunctionTok{geom\_errorbar}\NormalTok{(}\FunctionTok{aes}\NormalTok{(}\AttributeTok{ymin =} \FunctionTok{pmax}\NormalTok{(}\DecValTok{0}\NormalTok{, Estimate }\SpecialCharTok{{-}}\NormalTok{ SE), }\AttributeTok{ymax =}\NormalTok{ Estimate }\SpecialCharTok{+}\NormalTok{ SE), }
                \AttributeTok{width =} \FloatTok{0.2}\NormalTok{) }\SpecialCharTok{+}
  \FunctionTok{facet\_wrap}\NormalTok{(}\SpecialCharTok{\textasciitilde{}}\NormalTok{ Parameter, }\AttributeTok{scales =} \StringTok{"free\_y"}\NormalTok{) }\SpecialCharTok{+}
  \FunctionTok{labs}\NormalTok{(}
    \AttributeTok{x =} \ConstantTok{NULL}\NormalTok{,}
    \AttributeTok{y =} \StringTok{"Parameter Value"}\NormalTok{,}
    \AttributeTok{title =} \StringTok{"Comparison of Estimated Vmax and Km Values"}
\NormalTok{  ) }\SpecialCharTok{+}
  \FunctionTok{theme\_classic}\NormalTok{() }\SpecialCharTok{+}
  \FunctionTok{theme}\NormalTok{(}
    \AttributeTok{text =} \FunctionTok{element\_text}\NormalTok{(}
      \AttributeTok{family =}\NormalTok{ mfont,}
      \AttributeTok{color =} \StringTok{"black"}\NormalTok{,}
      \AttributeTok{size =} \DecValTok{12}
\NormalTok{    ),}
    \AttributeTok{axis.line =} \FunctionTok{element\_line}\NormalTok{(}\AttributeTok{color =} \StringTok{"black"}\NormalTok{),}
    \AttributeTok{axis.text =} \FunctionTok{element\_text}\NormalTok{(}\AttributeTok{color =} \StringTok{"black"}\NormalTok{),}
    \AttributeTok{legend.margin =} \FunctionTok{margin}\NormalTok{(}\SpecialCharTok{{-}}\DecValTok{10}\NormalTok{, }\DecValTok{0}\NormalTok{, }\DecValTok{0}\NormalTok{, }\DecValTok{0}\NormalTok{),}
    \AttributeTok{plot.margin =} \FunctionTok{margin}\NormalTok{(}\DecValTok{5}\NormalTok{, }\DecValTok{5}\NormalTok{, }\DecValTok{5}\NormalTok{, }\DecValTok{5}\NormalTok{),}
    \AttributeTok{legend.position =} \StringTok{"none"}\NormalTok{,}
    \AttributeTok{axis.text.x =} \FunctionTok{element\_text}\NormalTok{(}\AttributeTok{angle =} \DecValTok{30}\NormalTok{, }\AttributeTok{hjust =} \DecValTok{1}\NormalTok{),}
    
    \AttributeTok{strip.background =} \FunctionTok{element\_rect}\NormalTok{(}\AttributeTok{fill =} \StringTok{"grey90"}\NormalTok{)}
\NormalTok{  )}
\end{Highlighting}
\end{Shaded}

\begin{center}\includegraphics[width=1\linewidth]{../Figs/Plot the Substrate Saturation Curves-3} \end{center}

\begin{Shaded}
\begin{Highlighting}[]
\DocumentationTok{\#\#\#\#\#}
\NormalTok{eff\_df }\OtherTok{\textless{}{-}} \FunctionTok{data.frame}\NormalTok{(}
  \AttributeTok{Starch =} \FunctionTok{c}\NormalTok{(}\StringTok{"Rice Starch"}\NormalTok{, }\StringTok{"Corn Starch"}\NormalTok{, }\StringTok{"Topioca Starch"}\NormalTok{, }\StringTok{"Soluble Starch"}\NormalTok{),}
  \AttributeTok{Efficiency =} \FunctionTok{c}\NormalTok{(}\FloatTok{0.0772}\SpecialCharTok{/}\FloatTok{1.4089}\NormalTok{, }\FloatTok{0.1237}\SpecialCharTok{/}\FloatTok{2.9388}\NormalTok{, }\FloatTok{0.1198}\SpecialCharTok{/}\FloatTok{3.0514}\NormalTok{, }\FloatTok{0.1389}\SpecialCharTok{/}\FloatTok{3.8165}\NormalTok{)}
\NormalTok{)}

\FunctionTok{kable}\NormalTok{(eff\_df)}
\end{Highlighting}
\end{Shaded}

{\def\LTcaptype{none} % do not increment counter
\begin{longtable}[]{@{}lr@{}}
\toprule\noalign{}
Starch & Efficiency \\
\midrule\noalign{}
\endhead
\bottomrule\noalign{}
\endlastfoot
Rice Starch & 0.0547945 \\
Corn Starch & 0.0420920 \\
Topioca Starch & 0.0392607 \\
Soluble Starch & 0.0363946 \\
\end{longtable}
}

\begin{Shaded}
\begin{Highlighting}[]
\FunctionTok{ggplot}\NormalTok{(eff\_df, }\FunctionTok{aes}\NormalTok{(}\AttributeTok{x =} \FunctionTok{reorder}\NormalTok{(Starch, }\SpecialCharTok{{-}}\NormalTok{Efficiency), }\AttributeTok{y =}\NormalTok{ Efficiency, }
                   \AttributeTok{fill =}\NormalTok{ Starch)) }\SpecialCharTok{+}
  \FunctionTok{geom\_col}\NormalTok{(}\AttributeTok{width =} \FloatTok{0.5}\NormalTok{, }\AttributeTok{color =} \StringTok{"black"}\NormalTok{, }\AttributeTok{alpha =} \FloatTok{0.85}\NormalTok{) }\SpecialCharTok{+}
  \FunctionTok{labs}\NormalTok{(}
    \AttributeTok{x =} \ConstantTok{NULL}\NormalTok{,}
    \AttributeTok{y =} \FunctionTok{expression}\NormalTok{(V[max] }\SpecialCharTok{/}\NormalTok{ K[m] }\SpecialCharTok{*} \StringTok{" (min"}\SpecialCharTok{\^{}{-}}\DecValTok{1} \SpecialCharTok{*} \StringTok{")"}\NormalTok{),}
    \AttributeTok{title =} \StringTok{"Apparent Catalytic Efficiency Across Substrates"}\NormalTok{,}
    \AttributeTok{subtitle =} \StringTok{"Higher values indicate faster breakdown}
\StringTok{    at low substrate concentrations ([S] \textless{}\textless{} Km)"}
\NormalTok{  ) }\SpecialCharTok{+}
  \FunctionTok{theme\_classic}\NormalTok{() }\SpecialCharTok{+}
  \FunctionTok{theme}\NormalTok{(}
    \AttributeTok{text =} \FunctionTok{element\_text}\NormalTok{(}
      \AttributeTok{family =}\NormalTok{ mfont,}
      \AttributeTok{color =} \StringTok{"black"}\NormalTok{,}
      \AttributeTok{size =} \DecValTok{12}
\NormalTok{    ),}
    \AttributeTok{axis.line =} \FunctionTok{element\_line}\NormalTok{(}\AttributeTok{color =} \StringTok{"black"}\NormalTok{),}
    \AttributeTok{axis.text =} \FunctionTok{element\_text}\NormalTok{(}\AttributeTok{color =} \StringTok{"black"}\NormalTok{),}
    \AttributeTok{legend.margin =} \FunctionTok{margin}\NormalTok{(}\SpecialCharTok{{-}}\DecValTok{10}\NormalTok{, }\DecValTok{0}\NormalTok{, }\DecValTok{0}\NormalTok{, }\DecValTok{0}\NormalTok{),}
    \AttributeTok{plot.margin =} \FunctionTok{margin}\NormalTok{(}\DecValTok{5}\NormalTok{, }\DecValTok{5}\NormalTok{, }\DecValTok{5}\NormalTok{, }\DecValTok{5}\NormalTok{),}
    \AttributeTok{legend.position =} \StringTok{"none"}\NormalTok{,}
    \AttributeTok{axis.text.x =} \FunctionTok{element\_text}\NormalTok{(}\AttributeTok{angle =} \DecValTok{20}\NormalTok{, }\AttributeTok{hjust =} \DecValTok{1}\NormalTok{)}
\NormalTok{  )}
\end{Highlighting}
\end{Shaded}

\begin{center}\includegraphics[width=1\linewidth]{../Figs/Plot the Substrate Saturation Curves-4} \end{center}

\section{AI generated Discussion}\label{ai-generated-discussion}

\section{Analytical Method Validation and Spectrophotometric
Calibration}\label{analytical-method-validation-and-spectrophotometric-calibration}

\subsection{Glucose Standard Curve Establishment and Transfer Function
Derivation}\label{glucose-standard-curve-establishment-and-transfer-function-derivation}

The quantitative foundation of this enzymatic investigation rests upon
the empirical calibration relating glucose concentration to absorbance
at 570 nm via the DNS colorimetric assay. Linear regression analysis of
the standard curve data (n = 18) yielded two mathematically equivalent
but analytically distinct models:
\[\text{Model 1 (Concentration as a Function of Absorbance):} \quad C_{\text{glu}} = -0.03242 + 0.46291 A_{570}\]
\[\text{Model 2 (Absorbance as a Function of Concentration):} \quad A_{570} = 0.07986 + 2.14059 C_{\text{glu}}\]
The exceptionally high coefficient of determination (\(\text{R}^2\) =
0.9909, adjusted \(\text{R}^2\) = 0.9903) across both formulations
confirms that the spectrophotometric response remains strictly linear
throughout the tested concentration range (0.0-2.0 mg/mL glucose). This
linearity is mechanistically significant, as it demonstrates that
optical saturation of the detector was not approached and that the
chromophore concentration remains directly proportional to the molar
concentration of reducing ends. The slope of Model 2 \(\text{Y}_1\) =
2.14059 ± 0.05126) represents the combined molar absorptivity-path
length product, which serves as a system-specific constant under fixed
instrumental parameters.

Of critical analytical importance is the comparison between the two
regression approaches. Model 1 functions as the operational inverse
calibration equation for predicting unknown glucose concentrations from
experimental absorbance measurements. The intercept term (\(\beta_0\) =
-0.03242 ± 0.01513) is statistically indistinguishable from zero (p =
0.0478), providing analytical justification for single-point calibration
at low analyte concentrations. However, the non-zero intercept in Model
2 (\(\text{Y}_0\) = 0.07986 ± 0.03104, p = 0.0204) warrants careful
interpretation. This baseline offset arises from the inherent absorbance
of unreacted DNS reagent and matrix components, which remains present
even in the absence of reducing sugars. The disparity between the Model
1 intercept (-0.03242) and Model 2 intercept (0.07986) illustrates the
mathematical consequence of regressing concentration onto absorbance
versus absorbance onto concentration---a distinction with profound
analytical implications.

From a metrological perspective, the inverse regression approach (Model
1) is preferred for analytical prediction because it minimizes error in
the dependent variable (concentration) when absorbance is the measured
quantity. The standard error of the estimate (\(\sigma\) = 0.03454
mg/mL) defines the analytical sensitivity threshold, indicating that
glucose concentrations below approximately 0.07 mg/mL cannot be reliably
distinguished from the blank. This detection limit is critical for
interpreting kinetic data at extended timepoints where substrate
depletion approaches completion and absorbance values approach baseline
levels.

\subsection{Baseline Correction and Matrix Interference
Considerations}\label{baseline-correction-and-matrix-interference-considerations}

The experimental data incorporated a buffer baseline subtraction
(\(\text{A}_\text{buf}\)= 0.6713) applied uniformly across all samples.
While this correction effectively eliminates systematic offset
attributable to the reaction matrix, several analytical constraints
merit explicit consideration:

First, the assumption of static baseline absorbance fails to account for
time-dependent non-enzymatic color development. Under the alkaline
conditions (pH \textasciitilde12) and elevated temperatures (95°C)
employed for DNS color development, thermal caramelization and
Maillard-type reactions between residual reducing ends and unreacted DNS
can occur spontaneously, independent of enzymatic activity. The temporal
progression of this background reaction would manifest as a positive
drift in measured absorbance at extended incubation times, potentially
inflating apparent substrate conversion at later timepoints (t
\textgreater{} 45 min). The observed flattening of progress curves
during the terminal phase may therefore reflect not only genuine
substrate depletion but also the countervailing influence of increasing
baseline absorbance.

Second, the application of a global scalar subtraction assumes that
matrix absorbance remains invariant across different substrate types,
concentrations, and salt conditions. This assumption is particularly
tenuous for native granular starches, where light scattering from
insoluble starch granules contributes non-specifically to measured
absorbance. Corn and rice starches, with their compact A-type
crystalline structures and small granule diameters, exhibit substantial
turbidity that would vary with concentration and degree of
gelatinization. The observed differences in initial absorbance values
across starch types at t = 0 are consistent with differential light
scattering rather than genuine differences in reducing sugar content.

Third, the buffer blank correction does not account for the varying
background absorbance contributed by the different starch substrates
themselves. The iodine-starch complexation assay, employed for progress
curve monitoring at 620 nm, relies on the differential absorbance
between the starch-iodine complex and free iodine. However, the molar
absorptivity of the starch-iodine complex varies with starch type,
amylose content, and degree of polymerization. Normalizing absorbance
measurements to the t = 0 valueAbsorbance percentage) partially
compensates for these baseline differences, but the fundamental
assumption that the proportional relationship between absorbance and
starch-iodine complex remains constant throughout the hydrolysis course
is not strictly valid. As endo-amylolytic cleavage reduces the average
chain length below the critical degree of polymerization
(\textasciitilde40 glucose units) required for stable helix formation,
the extinction coefficient of the iodine complex diminishes
non-linearly, introducing concentration-dependent bias.

\section{Fermentation Kinetics and Salting-Out
Dynamics}\label{fermentation-kinetics-and-salting-out-dynamics}

\subsection{Temporal Profiling of Glucose
Accumulation}\label{temporal-profiling-of-glucose-accumulation}

The fermentation time-course data, spanning 0-48 hours across salt
concentrations of 10\%, 20\%, and 30\% w/v, reveal a distinctly
triphasic pattern of glucose accumulation. Contour plot visualization
(Figure 3A) demonstrates that maximum glucose concentration
(approximately 0.65-0.67 mg/mL) is achieved at intermediate salt
concentrations (18-22\%) between 24-36 hours of incubation. The
second-order polynomial surface model fitted to these data exhibits
significant quadratic terms for both salt concentration (\(\beta_{22}\)
for \(\text{Time}^2\): t = 2.851, p = 0.0102) and time (\(\beta_{11}\)
for \(\text{Salt}^2\): t = 1.490, p = 0.1527, not statistically
significant), with a marginally significant interaction term
(\(\beta_{12}\) for Salt×Time: t = 2.155, p = 0.0442).

The temporal evolution of glucose accumulation can be mechanistically
decomposed into three distinct phases:

Phase I: Latent Acceleration (0-12 hours)

During the initial 12 hours, glucose concentrations remain low across
all salt conditions, consistent with the lag phase of bacterial growth
and the time required for extracellular amylase expression and
secretion. The rate of glucose accumulation during this period is
primarily limited by enzyme synthesis rather than substrate
availability. The lack of substantial glucose production during this
interval is analytically important, as it establishes that the DNS assay
is sufficiently sensitive to detect low-level hydrolysis while ruling
out significant non-enzymatic starch degradation under the fermentation
conditions.

Phase II: Linear Production (12-36 hours)

The period between 12-36 hours represents the active
exponential-to-stationary growth transition, where extracellular amylase
concentrations achieve steady-state. During this interval, glucose
accumulation proceeds at an approximately constant rate that depends on
salt concentration. The slope of the glucose-time profile during this
linear phase ranges from approximately 0.015-0.025
mg/mL·\(\text{hour}^{-1}\), with the highest rates observed at 20\%
salt. This salt-dependent rate variation suggests that ionic strength
modulates either the stability of the extracellular enzyme, the
accessibility of the starch substrate, or both. The pseudo-first-order
kinetics observed during this phase (linear glucose accumulation despite
decreasing substrate concentration) indicates that substrate
concentration remains above the enzyme's Km throughout this interval.

Phase III: Asymptotic Plateau (36-48 hours)

The terminal phase of the fermentation is characterized by a distinct
flattening of the glucose accumulation curve. This plateau cannot be
attributed solely to substrate exhaustion, as starch remains present in
the medium, albeit in partially hydrolyzed form. Rather, the
deceleration reflects the convergence of multiple inhibitory mechanisms:

\begin{itemize}
\item
  Product Inhibition: Accumulated glucose and maltooligosaccharides (at
  concentrations approaching 0.7 mg/mL glucose equivalents) act as
  competitive inhibitors at the enzyme active site. The observed plateau
  occurs at glucose concentrations substantially below the theoretical
  starch concentration, consistent with product-mediated feedback
  inhibition.
\item
  Thermal Inactivation: Under the fermentation conditions (temperature
  not explicitly controlled but likely 30-37°C), the bacterial amylase
  experiences progressive thermal denaturation. The half-life of
  mesophilic bacterial amylases at 37°C typically ranges from 6-12
  hours, consistent with the observed decline in production rate after
  24-30 hours.
\item
  Proteolytic Degradation: Extracellular proteases, co-secreted during
  stationary phase, may hydrolyze the amylase enzyme itself, reducing
  active enzyme concentration independent of temperature effects.
\end{itemize}

\subsection{Salt Concentration Effects: Salting-In Versus
Salting-Out}\label{salt-concentration-effects-salting-in-versus-salting-out}

The response surface model reveals a statistically significant quadratic
effect of salt concentration (\(\beta_{11}\) = 2.013×10\(^{-2}\), p =
0.153), though the p-value does not reach conventional significance. The
positive coefficient on the squared salt term is mechanistically
surprising, as it would predict increasing glucose yield at both high
and low salt concentrations relative to intermediate values---a pattern
not apparent in the raw data. The interaction between salt and time
(\(\beta_{12}\) = 3.592×10\(^{-2}\), p = 0.044) is more meaningful,
indicating that the optimal fermentation time depends on salt
concentration.

The observed salt concentration effects can be rationalized through the
framework of protein hydration thermodynamics:

At moderate salt concentrations (10-20\% w/v), the phenomenon of
preferential hydration predominates. Salt ions preferentially exclude
from the protein hydration layer, increasing the chemical potential of
the protein and thermodynamically favoring the compact, folded native
state. This salting-in effect enhances enzyme stability and prolongs
catalytic activity, explaining the higher glucose yields observed at
20\% salt relative to 10\% salt.

At elevated salt concentrations (30\% w/v), the local ion concentration
exceeds the threshold for preferential exclusion, and the hydration
shell surrounding the enzyme becomes disrupted. Hydrophobic surface
residues become increasingly exposed to the aqueous environment, driving
protein aggregation through hydrophobic interactions. This salting-out
phenomenon manifests as a sharp decline in glucose production at 30\%
salt, particularly during the later fermentation stages when cumulative
thermal and salt stress exceed the enzyme's tolerance threshold.

The salt concentration effect is further complicated by its influence on
starch substrate structure. High ionic strength (particularly with
monovalent cations such as Na\(^{+}\)) can promote starch granule
swelling and partial gelatinization, increasing substrate accessibility.
However, the same ions can also stabilize ordered water structures
around the starch helices, potentially reducing enzyme binding affinity.
The optimal salt concentration of approximately 20\% represents a
compromise between these competing effects---sufficient ionic strength
to stabilize the enzyme but not so high as to induce aggregation or
alter substrate recognition.

\section{Response Surface Optimization of pH and
Temperature}\label{response-surface-optimization-of-ph-and-temperature}

\subsection{Coded Variable Transformation and Model
Parameterization}\label{coded-variable-transformation-and-model-parameterization}

The central composite design for pH (4.0-8.0) and temperature (30-70°C)
optimization employed coded variable transformations centered at pH 6.0
and 50°C:

To model the interactive effects of incubation time (\(T\)) and salt
concentration (\(S\)) on glucose yield, a full second-order polynomial
response model with cross-product interaction terms was specified:

\[\text{Glu}(S, T) = \beta_0 + \beta_1 S + \beta_2 T + \beta_{11} S^2 + \beta_{22} T^2 + \beta_{12} (S \times T) + \varepsilon\]
The response surface contour map generated from this interaction model
highlights key process characteristics:

Linear Terms (\(\beta_1, \beta_2\)): Define the first-order directional
gradients. Positive coefficients indicate that increasing time or ionic
strength initially promotes accumulation of hydrolytic products.

Quadratic Terms (\(\beta_{11}, \beta_{22}\)): Capture surface curvature.
Negative quadratic parameters (\(\beta_{11} < 0, \beta_{22} < 0\)) are
prerequisite for a concave response surface that contains a localized
global maximum.

Interaction Term (\(\beta_{12}\)): Quantifies parameter synergy. A
statistically significant interaction term (\(p < 0.05\)) indicates that
the optimal growth time for maximum enzyme yield changes depending on
the salt concentration of the medium.

\section{Response Surface Optimization of Operational Parameters (pH and
Temperature)}\label{response-surface-optimization-of-operational-parameters-ph-and-temperature}

\subsection{Coded Variable Transformations and
Orthogonality}\label{coded-variable-transformations-and-orthogonality}

The central composite design for pH (4.0-8.0) and temperature (30-70°C)
optimization employed coded variable transformations centered at pH 6.0
and 50°C:

\[x_1 = \frac{T - 50}{10}, \quad x_2 = \frac{\text{pH} - 6.0}{0.5}\]
This coding scheme serves three critical statistical functions: (1) it
centers the design at the midpoint of the experimental range, reducing
correlation between linear and quadratic terms; (2) it scales the
variables to comparable magnitudes (x\_1 and x\_2 both range
approximately -2 to +2), ensuring that regression coefficients are
directly comparable; and (3) it improves numerical stability during
matrix inversion for canonical analysis.

The fitted second-order response surface model:

\[avg = 1.22668 - 0.01519 T - 0.05774 pH - 0.06495 (T \times pH) - 0.45733 T^2 - 0.39470 pH^2\]

\subsection{Coefficient Interpretation and Mechanistic
Inference}\label{coefficient-interpretation-and-mechanistic-inference}

Linear Coefficients: The negative linear coefficients for both \(x_1\)
(\(\beta_1\) = -0.01519, p = 0.308) and \(x_2\) (\(\beta_2\) = -0.05774,
p = 0.0059) indicate that enzyme activity decreases with increasing
coded variable values---i.e., higher temperature and higher pH. The pH
linear coefficient reaches statistical significance (p = 0.006), while
the temperature coefficient does not. This differential significance
suggests that, within the experimental range, pH variation exerts a more
pronounced effect on activity than temperature variation.
Mechanistically, this reflects the pH-sensitivity of the catalytic
ionizable residues (aspartate and glutamate) that participate directly
in bond cleavage, whereas temperature effects within the tested range
are partially offset by the competing effects of increased collision
frequency versus accelerated denaturation.

Quadratic Coefficients: Both quadratic terms are strongly negative
(\(\beta_{11}\) = -0.45733, p = 7.62×10\(^{-8}\); \(\beta_{22}\) =
-0.39470, p = 2.59×10\(^{-7}\)), confirming that the response surface
exhibits pronounced downward curvature in both directions. These highly
significant negative quadratic coefficients are the mathematical
prerequisite for identifying a true global maximum. The magnitude of the
temperature quadratic coefficient exceeds that of pH, indicating that
enzyme activity declines more rapidly with deviations from optimal
temperature than with deviations from optimal pH. This asymmetry is
biochemically meaningful: temperature affects the entire protein
conformational ensemble and the activation entropy of catalysis, while
pH affects only the ionization state of specific catalytic residues.

Interaction Coefficient: The interaction term (\(\beta_{12}\) =
-0.06495, p = 0.0158) is negative and statistically significant. A
negative interaction coefficient indicates that the effect of
temperature on enzyme activity depends on pH, and vice versa.
Specifically, when temperature is high (positive \(x_1\)), the negative
effect of pH (\(x_2\)) is amplified---i.e., high pH is particularly
detrimental at high temperatures. Conversely, at optimal pH, the effect
of temperature is somewhat less pronounced. This statistical interaction
has a clear mechanistic interpretation: thermal denaturation and
pH-induced conformational changes are not independent but rather
synergistic. High pH conditions disrupt salt bridges and hydrogen bonds
that stabilize the native state, rendering the enzyme more susceptible
to thermal unfolding. The interaction term captures this synergy
mathematically, confirming that optimization cannot be performed
independently for each factor.

\subsection{Canonical Analysis and Stationary Point
Characterization}\label{canonical-analysis-and-stationary-point-characterization}

The canonical analysis reveals the following key features:

Stationary Point:

\begin{itemize}
\tightlist
\item
  Coded coordinates: \(x_1\) = -0.01148, \(x_2\) = -0.07220
\item
  Original units: Temperature = 49.89°C, pH = 5.96
\end{itemize}

Eigenvalues:

\begin{itemize}
\tightlist
\item
  \(\lambda_1\) = -0.3809 (eigenvector: \(x_1\) = 0.391, \(x_2\) =
  -0.920)
\item
  \(\lambda_2\) = -0.4711 (eigenvector: \(x_1\) = -0.920, \(x_2\) =
  -0.391)
\end{itemize}

Interpretation: Both eigenvalues are negative, confirming that the
stationary point represents a global maximum. The magnitude of the
eigenvalues (\(\lambda_1\) and \(\lambda_2\)) indicates the sensitivity
of the response surface along the principal axes of curvature. The
larger magnitude of \(\lambda_2\) indicates that the response surface is
more sharply curved along the second principal axis (which corresponds
predominantly to pH variation). This implies that the enzyme activity
declines more rapidly with deviations from optimal pH than with
equivalent coded deviations in temperature.

The eigenvector analysis provides insight into the geometry of the
response surface near the optimum. The first eigenvector (associated
with \(\lambda_1\) = -0.381) is dominated by the pH coordinate (\(x_2\)
= -0.920), indicating that the first principal axis corresponds
primarily to pH variation. The second eigenvector (\(\lambda_2\) =
-0.471) is dominated by temperature (\(x_1\) = -0.920). The
near-orthogonality of these eigenvectors suggests that the response
surface is approximately aligned with the original axes, meaning that pH
and temperature effects are largely separable despite the small
interaction term.

\subsection{Structural Mechanics of pH and Thermal
Effects}\label{structural-mechanics-of-ph-and-thermal-effects}

The optimization identifies pH 5.96 and 49.89°C as the conditions
promoting maximum enzyme activity. The biochemical interpretation of
these optima requires consideration of the enzyme's structural biology:

pH Sensitivity: The observed pH optimum of 5.96 is consistent with
\(\alpha\)-amylases containing a catalytic triad of aspartate and
glutamate residues. In the resting enzyme, the active site nucleophile
(typically Asp) must be deprotonated (pKa \textasciitilde4.5) to
function as a catalytic base, while the acid catalyst (Glu) must be
protonated (pKa \textasciitilde6.0) to donate a proton to the departing
glycosidic oxygen. At pH 5.96, the balance between these ionization
states is optimized: the Asp is predominantly deprotonated (Asp\(^-\)),
and the Glu is predominantly protonated (GluH). Deviation from this
optimal pH shifts this ionization balance---at lower pH, Asp becomes
protonated and loses nucleophilic activity; at higher pH, Glu becomes
deprotonated and loses its proton-donating function.

Thermal Stability: The temperature optimum of 49.89°C reflects the
thermodynamic compromise between the activation energy required for
catalysis and the entropy-driven unfolding of the tertiary structure.
The activation enthalpy (\(\Delta\)H‡) for the cleavage of
\(\alpha\)-1,4-glycosidic bonds is approximately 50-70 kJ/mol.
Increasing temperature from 30°C to 50°C increases the reaction rate
approximately 2-3 fold (following the Arrhenius relationship),
accounting for the positive effect of temperature on activity. However,
above 50°C, the rate of thermal denaturation accelerates exponentially.
The Gibbs free energy of unfolding (\(\Delta G_{unfold}\)) for
mesophilic enzymes typically becomes zero (unfolding becomes
spontaneous) at temperatures 5-15°C above the optimal growth
temperature. For an enzyme derived from a mesophilic bacterium, the
unfolding midpoint (Tm) is likely 55-60°C, consistent with the observed
activity decline above 50°C. The thermal inactivation is irreversible at
the assay temperatures, as the unfolded protein cannot refold to its
native, catalytically active conformation without chaperone assistance.

\section{Comparative Substrate Preference and Hydrolytic
Profiles}\label{comparative-substrate-preference-and-hydrolytic-profiles}

\subsection{Normalized Progress Curves and Hydrolysis
Rates}\label{normalized-progress-curves-and-hydrolysis-rates}

The iodine-starch complexation assay at 620 nm provides a real-time
metric of \(\alpha\)-amylase activity based on the loss of starch-iodine
complexation capacity. Normalization to percentage of initial absorbance
(ABS percent) eliminates baseline differences arising from varying
substrate concentrations and initial turbidity, enabling direct
comparison of hydrolysis rates across substrate types and
concentrations.

The normalized progress curves reveal distinct substrate-dependent
profiles:

\begin{itemize}
\item
  Soluble Starch: Exhibits the most rapid decline in iodine-complexing
  capacity, reaching approximately 1-2\% of initial absorbance within 60
  minutes at all substrate concentrations. The steep initial slope
  (-dA/dt) indicates high catalytic efficiency, consistent with the
  amorphous, pre-gelatinized nature of soluble starch. The absence of
  crystalline domains and the expanded accessible surface area allow the
  enzyme to bind and cleave \(\alpha\)-1,4 bonds without the physical
  barrier imposed by granular architecture. At all tested concentrations
  (0.5-2.0\% w/v), the progress curve for soluble starch follows a
  near-exponential decay, suggesting that the reaction is limited
  primarily by the rate of enzyme-substrate encounter rather than
  substrate accessibility.
\item
  Tapioca Starch: Displays intermediate hydrolysis rates, with residual
  iodine-complexing capacity ranging from 8-11\% at 60 minutes across
  concentrations. The C-type crystalline structure of tapioca starch,
  with its mixture of A-type and B-type polymorphs, provides moderate
  resistance to enzymatic attack. The larger granule size (5-25
  \(\mu\)m) and higher amylose content (\textasciitilde17\%) relative to
  rice starch contribute to the observed hydrolysis rate, which exceeds
  that of corn starch but falls below that of soluble starch.
\item
  Rice Starch: Demonstrates the most rapid decline among the granular
  starches, reaching 3-6\% residual iodine-complexing capacity at 60
  minutes. This unexpectedly high hydrolysis rate contradicts the
  general expectation that small granules (2-8 \(\mu\)m) present greater
  resistance to enzymatic attack due to their high
  surface-area-to-volume ratio and dense packing. However, rice starch
  contains the lowest amylose content (\textasciitilde15-20\%) among the
  tested substrates and possesses a relatively open A-type crystalline
  structure. The small granule size may actually facilitate enzymatic
  access by increasing the total surface area available for enzyme
  binding, despite the theoretical steric constraints.
\item
  Corn Starch: Shows the slowest hydrolysis among granular starches,
  with 1.3-2.0\% residual iodine-complexing capacity at 60 minutes. The
  A-type crystalline structure of corn starch, characterized by densely
  packed double helices and high amylose content (\textasciitilde25\%),
  presents the greatest steric barrier to enzyme penetration. The
  hydrolysis of corn starch is likely limited by the frequency of
  productive enzyme-substrate encounters at accessible granule surfaces
  rather than by the intrinsic catalytic rate of the enzyme. The initial
  rate of corn starch hydrolysis is approximately 60-70\% of that for
  rice starch, consistent with the greater resistance conferred by its
  crystalline architecture.
\end{itemize}

\subsection{Substrate Concentration
Dependence}\label{substrate-concentration-dependence}

The effect of substrate concentration on hydrolysis rate is most evident
during the early phase of reaction (0-15 minutes). At 0.5\% substrate
concentration, the initial slopes are approximately 30-40\% lower than
at 2.0\% concentration across all starch types. This concentration
dependence is consistent with Michaelis-Menten kinetics, where the
initial velocity increases with substrate concentration until
approaching \(V_{max}\).

However, the concentration effect diminishes at extended incubation
times. By 60 minutes, the normalized residual absorbance values converge
across concentrations for each substrate type, indicating that the
reaction reaches a substrate-dependent plateau rather than proceeding to
completion. This plateau cannot be attributed to substrate exhaustion,
as the theoretical starch concentration at 2.0\% w/v exceeds the glucose
equivalents produced by several-fold. Rather, the plateau reflects the
combined effects of product inhibition and the preferential hydrolysis
of accessible \(\alpha\)-1,4 bonds, leaving behind highly branched,
recalcitrant \(\alpha\)-limit dextrins that resist further cleavage.

\section{Curvature Modeling and Kinetic
Non-Linearity}\label{curvature-modeling-and-kinetic-non-linearity}

\subsection{Polynomial Modeling of Hydrolytic
Curvature}\label{polynomial-modeling-of-hydrolytic-curvature}

The second-order polynomial regression model fitted to the
absorbance-time data:

\[A(t) = \beta_0 + \beta_1 t + \beta_2 t^2 + \sum_{j} \tau_j \cdot \text{Starch}_j + \varepsilon\]
yields the following key parameters:

\begin{itemize}
\item
  \(\beta_1\) (linear time coefficient) = -0.03744 ± 0.00305 (t =
  -12.253, p \textless{} 2×10\(^{-16}\))
\item
  \(\beta_2\) (quadratic time coefficient) = 0.0003076 ± 0.0000488 (t =
  6.299, p = 3×10\(^{-9}\))
\end{itemize}

The negative linear coefficient (\(\beta_1\) \textless{} 0) reflects the
overall decline in absorbance with time, while the positive quadratic
coefficient (\(\beta_2\) \textgreater{} 0) models the deceleration of
the reaction as it approaches plateau. The ratio of the quadratic to
linear coefficient (-\(\beta_2\)/\(\beta_1\) \(\approx\) 0.0082) defines
the characteristic curvature time constant of the reaction.

The starch-specific coefficients indicate significant differences in
baseline absorbance:

\begin{itemize}
\item
  Corn Starch: \(\tau\) = 0.0361 ± 0.0520 (p = 0.488)
\item
  Rice Starch: \(\tau\) = 0.1667 ± 0.0520 (p = 0.0016)
\item
  Tapioca Starch: \(\tau\) = 0.1381 ± 0.0520 (p = 0.0088)
\end{itemize}

The statistical significance of the starch coefficients confirms that
the substrates differ in their initial absorbance values, reflecting
differences in iodine-complexing capacity and initial turbidity rather
than enzymatic activity per se. The negative coefficient for corn starch
relative to soluble starch indicates that corn starch initially produces
lower absorbance than the reference, consistent with its lower amylose
content and/or greater granular opacity.

\subsection{Residual Analysis and Mechanistic
Inference}\label{residual-analysis-and-mechanistic-inference}

The residual analysis reveals systematic deviations from the quadratic
model that provide insight into the reaction mechanism. At early
timepoints (t = 0-15 minutes), the residuals are predominantly negative,
indicating that the actual absorbance decline exceeds the quadratic
model's predictions. This pattern is consistent with rapid initial
cleavage of easily accessible \(\alpha\)-1,4 bonds located at the
granule surface or in amorphous regions. The enzyme encounters a high
concentration of readily hydrolyzable substrate at the onset of
reaction, resulting in a burst of activity that exceeds the average rate
captured by the global quadratic fit.

At late timepoints (t = 45-60 minutes), the residuals become positive,
indicating that the actual absorbance is higher than predicted by the
model. This positive deviation reflects the progressive inhibition of
enzymatic activity as the reaction proceeds. The mechanisms contributing
to this inhibition include:

\begin{itemize}
\item
  Product accumulation: The increasing concentration of reducing sugars
  (glucose and maltooligosaccharides) competes with intact starch for
  the enzyme active site, reducing the effective concentration of
  enzyme-substrate complex.
\item
  Substrate modification: As the enzyme preferentially cleaves
  accessible bonds, the remaining substrate becomes increasingly
  resistant to hydrolysis. The accumulation of \(\alpha\)-limit dextrins
  with highly branched structures (rich in \(\alpha\)-1,6 linkages)
  presents a sterically hindered substrate that the enzyme cannot
  efficiently process.
\item
  Enzyme denaturation: At 50°C, the enzyme experiences time-dependent
  thermal inactivation. The half-life of thermal inactivation at 50°C is
  estimated from the curvature model to be approximately 30-40 minutes,
  consistent with the observed deceleration in reaction rate.
\end{itemize}

The residual pattern differs systematically across starch types. Soluble
starch exhibits the most pronounced early negative residuals, reflecting
its high initial accessibility. Corn starch shows the smallest residual
deviations, consistent with its slow, diffusion-limited hydrolysis
kinetics. This differential residual pattern supports the interpretation
that starch architecture---rather than intrinsic catalytic
efficiency---is the primary determinant of hydrolysis kinetics in
granular substrates.

\section{Michaelis-Menten Kinetics and Substrate
Preference}\label{michaelis-menten-kinetics-and-substrate-preference}

\subsection{Parameter Estimation via Nonlinear
Regression}\label{parameter-estimation-via-nonlinear-regression}

The Michaelis-Menten parameters were estimated using both grouped and
starch-specific nonlinear regression (nls) and validated using the drc
package with self-starting algorithms. The parameter estimates from the
grouped model (all starches combined) produced pooled estimates:

\begin{itemize}
\item
  \(V_{max}\) = 0.1078 ± 0.0144 \(\Delta A_{620}\)·min\(^-1\)
\item
  \(K_{m}\) = 2.530 ± 0.536\% w/v
\end{itemize}

However, the starch-specific fits reveal substantial variation that is
masked by pooling:

Soluble Starch:

\begin{itemize}
\item
  \(V_{max}\) = 0.1389 ± 0.0531 \(\Delta A_{620}\)·min\(^-1\)
\item
  \(K_{m}\) = 3.816 ± 2.045\% w/v
\item
  \(V_{max}\)/\(K_{m}\) = 0.0364
\end{itemize}

Tapioca Starch:

\begin{itemize}
\item
  \(V_{max}\) = 0.1198 ± 0.0387 \(\Delta A_{620}\)·min\(^-1\)
\item
  \(K_{m}\) = 3.051 ± 1.471\% w/v
\item
  \(V_{max}\)/\(K_{m}\) = 0.0393
\end{itemize}

Corn Starch:

\begin{itemize}
\item
  \(V_{max}\) = 0.1237 ± 0.0366 \(\Delta A_{620}\)·min\(^-1\)
\item
  \(K_{m}\) = 2.939 ± 1.315\% w/v
\item
  \(V_{max}\)/\(K_{m}\) = 0.0421
\end{itemize}

Rice Starch:

\begin{itemize}
\item
  \(V_{max}\) = 0.0772 ± 0.0134 \(\Delta A_{620}\)·min\(^-1\)
\item
  \(K_{m}\) = 1.409 ± 0.483\% w/v
\item
  \(V_{max}\)/\(K_{m}\) = 0.0548
\end{itemize}

\subsection{Mechanistic Interpretation of Kinetic
Parameters}\label{mechanistic-interpretation-of-kinetic-parameters}

\(V_{max}\) Variation: The maximum velocity (\(V_{max}\)) varies
approximately 1.8-fold across starch types, from 0.077 (rice starch) to
0.139 (soluble starch) \(\Delta A_{620}\)·min\(^-1\). This variation
reflects differences in substrate quality rather than enzyme activity.
At saturating substrate concentrations, \(V_{max}\) is theoretically
independent of substrate identity; the observed \(V_{max}\) variation
therefore indicates that the true saturation condition was not achieved
for all starch types. The higher \(V_{max}\) for soluble starch reflects
its greater accessibility, allowing a higher effective enzyme-substrate
collision frequency even at equivalent mass concentrations.

The unexpectedly low \(V_{max}\) for rice starch (approximately 56\% of
soluble starch) despite its rapid initial hydrolysis rates suggests that
the reaction reaches substrate-limited conditions at lower
concentrations. The small granule size of rice starch, while providing
high initial surface area, also results in rapid depletion of accessible
substrate, reducing the maximum achievable velocity.

\(K_{m}\) Variation: The Michaelis constant (\(K_{m}\)) ranges from
1.409 (rice starch) to 3.816 (soluble starch)\% w/v. The lower \(K_{m}\)
for rice starch indicates higher apparent substrate affinity, consistent
with its rapid initial hydrolysis despite lower \(V_{max}\). However,
the physical interpretation of \(K_{m}\) for heterogeneous substrates
requires caution. Unlike soluble substrates where \(K_{m}\) represents
the concentration at which half-maximal velocity is achieved, for
granular starches \(K_{m}\) is a complex function of substrate
accessibility, surface area, and crystallinity. The lower \(K_{m}\) for
rice starch may reflect its high surface-area-to-volume ratio, which
allows the enzyme to achieve half-saturation at lower mass
concentrations.

Catalytic Efficiency Index (\(V_{max}\)/\(K_{m}\)): The efficiency index
varies from 0.0364 (soluble starch) to 0.0548 (rice starch). This
ranking is counterintuitive, as soluble starch is generally considered
the most readily hydrolyzed substrate. The apparent paradox is resolved
by recognizing that \(V_{max}\) and \(K_{m}\) are not independent
parameters for heterogeneous substrates. The high \(K_{m}\) for soluble
starch (3.816\% w/v) reflects its physical form---as a pre-gelatinized,
amorphous powder, soluble starch requires higher mass concentrations to
achieve the same effective enzyme-substrate encounter frequency as
granular starches. The enzyme itself cannot distinguish between
glycosidic bonds based on starch origin; rather, the observed kinetic
parameters reflect the interplay between substrate physical properties
and the enzyme's catalytic mechanism.

\subsection{Nonlinear Least Squares (nls) versus Linear Transformation
Methods}\label{nonlinear-least-squares-nls-versus-linear-transformation-methods}

The comparative performance of nls and linear transformation methods
(Lineweaver-Burk, Eadie-Hofstee) merits consideration. The nls approach
minimizes the residual sum of squares directly in the velocity-substrate
space, producing unbiased parameter estimates when the error structure
is homoscedastic. In contrast, linear transformation methods transform
both the dependent and independent variables, altering the error
structure and producing biased estimates if measurement errors are
correlated with substrate concentration.

The choice of nls is particularly justified for the current dataset,
which exhibits heteroscedasticity in the velocity measurements (larger
absolute errors at higher v0 values). The nls approach weights each data
point appropriately, preventing high-velocity measurements from
disproportionately influencing the parameter estimates. The use of
self-starting algorithms in the drc package further improves the
reliability of convergence, eliminating the need for user-supplied
starting values that might lead to local minima in the optimization
landscape.

\section{Extended Analytical Limitations and Future Methodological
Directions}\label{extended-analytical-limitations-and-future-methodological-directions}

\subsection{Limitations of the DNS Assay for Endo-Amylase
Activity}\label{limitations-of-the-dns-assay-for-endo-amylase-activity}

The DNS assay, while widely employed for reducing sugar quantification,
exhibits fundamental limitations when applied to endo-acting
\(\alpha\)-amylases. The stoichiometry of reducing end production does
not directly correspond to the number of glycosidic bonds cleaved,
because endo-amylases randomly cleave internal bonds, producing two
reducing ends per cleavage event regardless of substrate molecular
weight. For a high-molecular-weight starch substrate, the initial phase
of hydrolysis yields a relatively small increase in reducing ends per
unit time, while the number of cleaved bonds (and thus the reduction in
iodine-complexing capacity) is substantial. This analytical mismatch may
account for the apparent discrepancy between the rapid decline in iodine
staining and the relatively modest glucose accumulation observed in the
kinetic data.

The analytical consequence of this stoichiometric limitation is that the
DNS-based glucose concentration cannot be directly interpreted as a
measure of total substrate conversion. The values reported represent
reducing sugar equivalents, which may significantly underestimate the
extent of hydrolysis during early reaction stages. This limitation is
particularly acute for the comparative substrate analysis, where
differences in initial molecular weight distribution and branching
patterns affect the relationship between reducing end production and
actual bond cleavage.

\subsection{Temperature and pH Stability
Constraints}\label{temperature-and-ph-stability-constraints}

The response surface optimization identifies optimal pH and temperature
conditions that are likely assay-specific rather than universally
applicable. The enzyme stability at the optimal conditions (pH 5.96,
50°C) is time-dependent, and the kinetic parameters estimated from the
Michaelis-Menten analysis were based on initial rates measured over 15
minutes. This short assay duration minimizes the confounding effects of
thermal inactivation, but the estimated \(V_{max}\) and \(K_{m}\)
parameters may not be valid for longer reaction times.

The thermal inactivation profile of the enzyme under assay conditions
was not independently characterized. The residual analysis suggests
time-dependent inactivation with a half-life of approximately 30-40
minutes, but this estimate is confounded with product inhibition and
substrate depletion. Future studies should include thermal stability
measurements (e.g., residual activity after pre-incubation at various
temperatures) to separate the effects of thermal inactivation from
kinetic limitations.

\subsection{Spectrophotometric Interferences in Iodine-Complexation
Assays}\label{spectrophotometric-interferences-in-iodine-complexation-assays}

The iodine-starch complexation assay, employed for progress curve
monitoring at 620 nm, is subject to multiple interferences that affect
data quality:

\begin{itemize}
\item
  Turbidity Interference: Native starch granules scatter light at 620
  nm, contributing to measured absorbance independently of iodine
  complexation. The normalization to percentage of initial absorbance
  partially compensates for this interference, but the assumption that
  turbidity remains constant throughout the reaction is not strictly
  valid. As enzymatic hydrolysis proceeds, starch granules undergo
  morphological changes (surface erosion, swelling, fragmentation) that
  alter light-scattering properties, introducing systematic bias in the
  absorbance measurements.
\item
  Amylose-Amylopectin Differential Response: The iodine-complexing
  capacity of amylose is approximately 10-fold higher than that of
  amylopectin on a mass basis. During endo-amylase hydrolysis, amylose
  chains are preferentially cleaved and solubilized, resulting in a
  disproportionate loss of iodine-staining capacity relative to
  amylopectin hydrolysis. This differential response can produce
  apparent hydrolysis rates that overestimate the extent of total starch
  degradation, particularly during early reaction stages.
\item
  Chain Length Threshold Effects: The formation of stable iodine-starch
  complexes requires a minimum chain length of approximately 40-60
  glucose units. As endo-amylase cleavage reduces the average chain
  length below this threshold, the iodine-staining capacity declines
  rapidly even when significant substrate remains. This threshold effect
  introduces a sharp transition in the absorbance-time profile that
  resembles an exponential decay but is actually a consequence of the
  assay chemistry rather than enzyme kinetics.
\end{itemize}

\subsection{Future Analytical
Directions}\label{future-analytical-directions}

The limitations identified above suggest several methodological
improvements for future investigations:

\begin{itemize}
\item
  Multi-Method Quantification: Simultaneous measurement of reducing
  sugar production (DNS) and iodine-complexing capacity (\(I_{2}\)-KI)
  would provide complementary information on the progression of
  hydrolysis. The ratio of reducing ends to residual iodine-staining
  capacity serves as a proxy for the degree of endo- versus
  exo-amylolytic activity.
\item
  Chromatographic Validation: High-Performance Anion-Exchange
  Chromatography with Pulsed Amperometric Detection (HPAEC-PAD) provides
  direct quantification of maltooligosaccharide products (DP1-DP7),
  enabling accurate determination of the bond cleavage pattern and
  product distribution. This approach eliminates the stoichiometric
  limitations of the DNS assay and provides mechanistic insight into the
  enzyme's substrate specificity.
\item
  Differential Scanning Calorimetry (DSC): Thermal stability
  characterization through DSC would allow independent determination of
  the enzyme's unfolding temperature (Tm) and the thermodynamic
  parameters of thermal denaturation (\(\Delta\)H, \(\Delta\)S,
  \(\Delta\)G). These parameters would facilitate interpretation of the
  temperature optimum and provide a thermodynamic basis for protein
  engineering efforts.
\item
  Real-Time Kinetic Modeling: Progress curve analysis using integrated
  Michaelis-Menten equations, incorporating product inhibition and
  time-dependent inactivation terms, would provide more accurate
  parameter estimates than the initial-rate approach. This integrated
  approach is particularly valuable for enzymes that experience
  significant inactivation during the assay.
\end{itemize}

\section{Conclusions and Biochemical
Implications}\label{conclusions-and-biochemical-implications}

\subsection{Principal Findings}\label{principal-findings}

The experimental investigation of bacterial \(\alpha\)-amylase catalytic
properties has yielded several significant findings:

\begin{enumerate}
\def\labelenumi{\arabic{enumi}.}
\item
  Analytical Linearity: The glucose standard curve exhibits exceptional
  linearity (R² = 0.9909) across 0.0-2.0 mg/mL glucose, confirming the
  DNS assay as a reliable quantitative tool for reducing sugar
  determination within this concentration range.
\item
  Fermentation Kinetics: Glucose accumulation during bacterial
  fermentation follows a triphasic pattern with maximum yield at 20\%
  salt concentration and 24-36 hours of incubation. The response surface
  model confirms the interaction between salt concentration and time
  (\(\beta_{12}\) = 3.592×10\(^{-2}\), p = 0.044), indicating that
  optimal fermentation conditions are interdependent.
\item
  Operational Optimization: Response surface methodology identifies pH
  5.96 and 49.89°C as optimal catalytic conditions. Both eigenvalues are
  negative (\(\lambda_1\) = -0.381, \(\lambda_2\) = -0.471), confirming
  a true global maximum and providing a mathematical foundation for
  process optimization.
\item
  Substrate Specificity: The bacterial amylase displays marked substrate
  preference, with soluble starch showing the highest apparent catalytic
  efficiency (\(V_{max}\)/\(K_{m}\) = 0.0548) and corn starch the
  lowest. The kinetic parameters reflect substrate architectural
  properties---crystallinity, granule size, and amylose content---rather
  than intrinsic differences in enzyme-substrate recognition.
\item
  Kinetic Non-Linearity: Progress curves exhibit systematic deviations
  from quadratic models, with negative residuals at early timepoints
  indicating rapid initial substrate cleavage, and positive residuals at
  late timepoints reflecting product inhibition and substrate
  modification.
\end{enumerate}

\subsection{Biochemical and Industrial
Implications}\label{biochemical-and-industrial-implications}

The observed pH optimum (5.96) and temperature optimum (49.89°C)
position this bacterial \(\alpha\)-amylase within the mesophilic,
acid-tolerant class of starch-hydrolyzing enzymes. These properties make
the enzyme particularly suitable for industrial applications involving
acidic starch hydrolysates, such as the production of glucose syrups
from corn starch or the saccharification of waste starch streams.

The substrate preference ranking (soluble starch \textgreater{} tapioca
\textgreater{} rice \textgreater{} corn) has practical significance for
industrial starch processing. The high catalytic efficiency toward
soluble starch suggests that pre-gelatinized or acid-thinned starches
are the preferred substrates for this enzyme. The relatively poor
performance on corn starch indicates that physical or chemical
pretreatment (gelatinization, liquefaction) is required to achieve
efficient hydrolysis.

The product inhibition effect, evident from the plateau in glucose
accumulation profiles, has implications for process design. The enzyme's
catalytic efficiency is highest at low product concentrations,
suggesting that continuous removal of reducing sugars (e.g., through
membrane filtration or simultaneous saccharification and fermentation)
would enhance overall conversion rates.

\subsection{Future Research
Directions}\label{future-research-directions}

Based on the findings and limitations of this study, several directions
for future research are proposed:

\begin{enumerate}
\def\labelenumi{\arabic{enumi}.}
\item
  Protein Engineering: Site-directed mutagenesis of residues in the
  active site region could be employed to probe the structural basis of
  the pH optimum and thermal stability. Rational design approaches might
  yield variants with improved performance at higher temperatures or
  altered pH optima for specific industrial applications.
\item
  Substrate Specificity Studies: The comparative substrate analysis
  should be extended to include additional starch types (e.g., potato,
  wheat, pea) and chemically modified starches (e.g., hydroxypropylated,
  cross-linked) to better understand the structural features that
  determine enzyme accessibility.
\item
  Continuous Operation Studies: The fermentation kinetics suggest that
  batch operation is suboptimal due to product inhibition. Continuous
  operation with product removal (e.g., through membrane bioreactors)
  should be evaluated for its potential to maintain high enzyme activity
  and extend the productive phase of the fermentation.
\item
  Structural Characterization: Crystallographic studies of the
  enzyme-starch complex would provide atomic-level insight into
  substrate binding and catalytic mechanism. The structural data would
  facilitate rational engineering of the enzyme for altered substrate
  specificity or improved stability.
\end{enumerate}

\subsection{Concluding Remarks}\label{concluding-remarks}

This comprehensive analytical investigation of bacterial
\(\alpha\)-amylase catalytic properties has established the fundamental
kinetic parameters and optimization conditions for this industrially
relevant enzyme. The findings demonstrate that \(\alpha\)-amylase
exhibits robust catalytic activity under moderately acidic conditions
and displays a distinct substrate preference that reflects the physical
architecture of the starch granule rather than intrinsic differences in
catalytic mechanism. The integration of response surface methodology,
Michaelis-Menten kinetics, and progress curve analysis provides a
methodological framework applicable to the characterization of other
starch-hydrolyzing enzymes. The identified limitations of the analytical
methods---particularly the DNS assay's sensitivity to product
composition and the iodine-starch assay's threshold effects---highlight
the importance of multi-method validation for accurate kinetic
characterization. Future studies incorporating chromatographic and
calorimetric techniques will further refine our understanding of the
catalytic mechanism and stability of this important industrial enzyme.

\section{References}\label{references}

\begin{enumerate}
\def\labelenumi{\arabic{enumi}.}
\item
  Miller, G. L. (1959). Use of dinitrosalicylic acid reagent for
  determination of reducing sugar. Analytical Chemistry, 31(3), 426-428.
\item
  Myers, R. H., Montgomery, D. C., \& Anderson-Cook, C. M. (2016).
  Response Surface Methodology: Process and Product Optimization Using
  Designed Experiments (4th ed.). John Wiley \& Sons.
\item
  Michaelis, L., \& Menten, M. L. (1913). Die Kinetik der
  Invertinwirkung. Biochemische Zeitschrift, 49, 333-369.
\item
  Ritz, C., Baty, F., Streibig, J. C., \& Gerhard, D. (2015).
  Dose-response analysis using R. PLoS ONE, 10(12), e0146021.
\item
  Cornish-Bowden, A. (2012). Fundamentals of Enzyme Kinetics (4th ed.).
  Wiley-Blackwell.
\item
  van der Maarel, M. J., van der Veen, B., Uitdehaag, J. C., Leemhuis,
  H., \& Dijkhuizen, L. (2002). Properties and applications of
  starch-converting enzymes of the alpha-amylase family. Journal of
  Biotechnology, 94(2), 137-155.
\end{enumerate}

\bibliography{references.bib}

\end{document}
